%% Journal of Intelligent Manufacturing -- Springer Nature template (sn-jnl), author-year references.
%DIF LATEXDIFF DIFFERENCE FILE
%DIF DEL old_sametitle.tex   Fri Oct 23 17:49:53 2026
%DIF ADD main.tex            Fri Oct 23 17:49:07 2026
%% Line numbers for review (option lineno); remove the option for the final version.
\documentclass[pdflatex,sn-basic,lineno]{sn-jnl}

\usepackage{graphicx}
\usepackage{multirow}
\usepackage{amsmath,amssymb,amsfonts}
\usepackage{amsthm}
\usepackage{booktabs}
\usepackage{array}
\usepackage{tikz}
\usetikzlibrary{arrows.meta,positioning,calc}
\usepackage{siunitx}
\usepackage{manyfoot}   % the class's begin-document footnote hook expects this package to be loaded
\usepackage{placeins}   % \FloatBarrier keeps floats from drifting past sections
\graphicspath{{figures/}}

\raggedbottom

%% Float placement: allow figures and tables to share pages with text and with each other.
\renewcommand{\topfraction}{0.9}
\renewcommand{\bottomfraction}{0.7}
\renewcommand{\textfraction}{0.07}
\renewcommand{\floatpagefraction}{0.8}
\setcounter{topnumber}{3}
\setcounter{totalnumber}{4}
\setlength{\textfloatsep}{10pt plus 3pt minus 2pt}
\setlength{\floatsep}{10pt plus 3pt minus 2pt}
\setlength{\intextsep}{10pt plus 3pt minus 2pt}
\setlength{\belowcaptionskip}{3pt}
%% Float pages fill from the top instead of centring their floats.
\makeatletter
\setlength{\@fptop}{0pt}
\setlength{\@fpsep}{12pt plus 2fil}
%DIF 35c35
%DIF < % the appendix tables fill two float pages, which are built after the appendix ends: hence global
%DIF -------
% the appendix tables fill float pages, which are built after the appendix ends: hence global %DIF > 
%DIF -------
\newcommand{\appendixfloatpages}{\global\@fpsep=6pt plus 2fil\relax}
\setlength{\@fpbot}{0pt plus 1fil}
\makeatother
%% Table body size: the class's own 8 bp table font (its \footnotesize is 7 pt and its \scriptsize 9 pt).
\newcommand{\tabsize}{\fontsize{8}{9.5}\selectfont}
%DIF 41c41
%DIF < \newcommand{\tabsizesmall}{\footnotesize}% the appendix tables (7/8 pt), set on two float pages
%DIF -------
\newcommand{\tabsizesmall}{\footnotesize}% the appendix tables (7/8 pt) %DIF > 
%DIF -------
%% Reference list: the class sets it at normal size with entries 1em apart; the 8 bp table font with
%% hanging indents keeps the entries distinct on fewer pages.
\setlength{\bibsep}{0pt}
\makeatletter
\def\bibfont{\reset@font\fontfamily{\rmdefault}\fontsize{8}{9.5}\selectfont}
\makeatother
%% DOIs in the reference list as full links (https://doi.org/...), as the journal asks.
\newcommand{\doi}[1]{\href{https://doi.org/#1}{\nolinkurl{https://doi.org/#1}}}
%% Ragged-right paragraph column for text tables.
\newcolumntype{L}[1]{>{\raggedright\arraybackslash\hspace{0pt}}p{#1}}
%DIF 52a52-53
%% Online Resource 1 (supplementary material) is cited as ESM S1--S20. %DIF > 
\newcommand{\esm}[1]{Online Resource~1, #1} %DIF > 
%DIF -------

%% Set \anontrue to produce an anonymised copy (no author block, no identifying declarations).
\newcommand{\pending}[1]{\textit{[#1 pending]}}
\newif\ifanon
\anonfalse
%DIF PREAMBLE EXTENSION ADDED BY LATEXDIFF
%DIF UNDERLINE PREAMBLE %DIF PREAMBLE
\RequirePackage[normalem]{ulem} %DIF PREAMBLE
\RequirePackage{color}\definecolor{RED}{rgb}{1,0,0}\definecolor{BLUE}{rgb}{0,0,1} %DIF PREAMBLE
\providecommand{\DIFadd}[1]{{\protect\color{blue}\uwave{#1}}} %DIF PREAMBLE
\providecommand{\DIFdel}[1]{{\protect\color{red}\sout{#1}}}                      %DIF PREAMBLE
%DIF SAFE PREAMBLE %DIF PREAMBLE
\providecommand{\DIFaddbegin}{} %DIF PREAMBLE
\providecommand{\DIFaddend}{} %DIF PREAMBLE
\providecommand{\DIFdelbegin}{} %DIF PREAMBLE
\providecommand{\DIFdelend}{} %DIF PREAMBLE
\providecommand{\DIFmodbegin}{} %DIF PREAMBLE
\providecommand{\DIFmodend}{} %DIF PREAMBLE
%DIF FLOATSAFE PREAMBLE %DIF PREAMBLE
\providecommand{\DIFaddFL}[1]{\DIFadd{#1}} %DIF PREAMBLE
\providecommand{\DIFdelFL}[1]{\DIFdel{#1}} %DIF PREAMBLE
\providecommand{\DIFaddbeginFL}{} %DIF PREAMBLE
\providecommand{\DIFaddendFL}{} %DIF PREAMBLE
\providecommand{\DIFdelbeginFL}{} %DIF PREAMBLE
\providecommand{\DIFdelendFL}{} %DIF PREAMBLE
\newcommand{\DIFscaledelfig}{0.5}
%DIF HIGHLIGHTGRAPHICS PREAMBLE %DIF PREAMBLE
\RequirePackage{settobox} %DIF PREAMBLE
\RequirePackage{letltxmacro} %DIF PREAMBLE
\newsavebox{\DIFdelgraphicsbox} %DIF PREAMBLE
\newlength{\DIFdelgraphicswidth} %DIF PREAMBLE
\newlength{\DIFdelgraphicsheight} %DIF PREAMBLE
% store original definition of \includegraphics %DIF PREAMBLE
\LetLtxMacro{\DIFOincludegraphics}{\includegraphics} %DIF PREAMBLE
\newcommand{\DIFaddincludegraphics}[2][]{{\color{blue}\fbox{\DIFOincludegraphics[#1]{#2}}}} %DIF PREAMBLE
\newcommand{\DIFdelincludegraphics}[2][]{% %DIF PREAMBLE
\sbox{\DIFdelgraphicsbox}{\DIFOincludegraphics[#1]{#2}}% %DIF PREAMBLE
\settoboxwidth{\DIFdelgraphicswidth}{\DIFdelgraphicsbox} %DIF PREAMBLE
\settoboxtotalheight{\DIFdelgraphicsheight}{\DIFdelgraphicsbox} %DIF PREAMBLE
\scalebox{\DIFscaledelfig}{% %DIF PREAMBLE
\parbox[b]{\DIFdelgraphicswidth}{\usebox{\DIFdelgraphicsbox}\\[-\baselineskip] \rule{\DIFdelgraphicswidth}{0em}}\llap{\resizebox{\DIFdelgraphicswidth}{\DIFdelgraphicsheight}{% %DIF PREAMBLE
\setlength{\unitlength}{\DIFdelgraphicswidth}% %DIF PREAMBLE
\begin{picture}(1,1)% %DIF PREAMBLE
\thicklines\linethickness{2pt} %DIF PREAMBLE
{\color[rgb]{1,0,0}\put(0,0){\framebox(1,1){}}}% %DIF PREAMBLE
{\color[rgb]{1,0,0}\put(0,0){\line( 1,1){1}}}% %DIF PREAMBLE
{\color[rgb]{1,0,0}\put(0,1){\line(1,-1){1}}}% %DIF PREAMBLE
\end{picture}% %DIF PREAMBLE
}\hspace*{3pt}}} %DIF PREAMBLE
} %DIF PREAMBLE
\LetLtxMacro{\DIFOaddbegin}{\DIFaddbegin} %DIF PREAMBLE
\LetLtxMacro{\DIFOaddend}{\DIFaddend} %DIF PREAMBLE
\LetLtxMacro{\DIFOdelbegin}{\DIFdelbegin} %DIF PREAMBLE
\LetLtxMacro{\DIFOdelend}{\DIFdelend} %DIF PREAMBLE
\DeclareRobustCommand{\DIFaddbegin}{\DIFOaddbegin \let\includegraphics\DIFaddincludegraphics} %DIF PREAMBLE
\DeclareRobustCommand{\DIFaddend}{\DIFOaddend \let\includegraphics\DIFOincludegraphics} %DIF PREAMBLE
\DeclareRobustCommand{\DIFdelbegin}{\DIFOdelbegin \let\includegraphics\DIFdelincludegraphics} %DIF PREAMBLE
\DeclareRobustCommand{\DIFdelend}{\DIFOaddend \let\includegraphics\DIFOincludegraphics} %DIF PREAMBLE
\LetLtxMacro{\DIFOaddbeginFL}{\DIFaddbeginFL} %DIF PREAMBLE
\LetLtxMacro{\DIFOaddendFL}{\DIFaddendFL} %DIF PREAMBLE
\LetLtxMacro{\DIFOdelbeginFL}{\DIFdelbeginFL} %DIF PREAMBLE
\LetLtxMacro{\DIFOdelendFL}{\DIFdelendFL} %DIF PREAMBLE
\DeclareRobustCommand{\DIFaddbeginFL}{\DIFOaddbeginFL \let\includegraphics\DIFaddincludegraphics} %DIF PREAMBLE
\DeclareRobustCommand{\DIFaddendFL}{\DIFOaddendFL \let\includegraphics\DIFOincludegraphics} %DIF PREAMBLE
\DeclareRobustCommand{\DIFdelbeginFL}{\DIFOdelbeginFL \let\includegraphics\DIFdelincludegraphics} %DIF PREAMBLE
\DeclareRobustCommand{\DIFdelendFL}{\DIFOaddendFL \let\includegraphics\DIFOincludegraphics} %DIF PREAMBLE
%DIF COLORLISTINGS PREAMBLE %DIF PREAMBLE
\RequirePackage{listings} %DIF PREAMBLE
\RequirePackage{color} %DIF PREAMBLE
\lstdefinelanguage{DIFcode}{ %DIF PREAMBLE
%DIF DIFCODE_UNDERLINE %DIF PREAMBLE
  moredelim=[il][\color{red}\sout]{\%DIF\ <\ }, %DIF PREAMBLE
  moredelim=[il][\color{blue}\uwave]{\%DIF\ >\ } %DIF PREAMBLE
} %DIF PREAMBLE
\lstdefinestyle{DIFverbatimstyle}{ %DIF PREAMBLE
	language=DIFcode, %DIF PREAMBLE
	basicstyle=\ttfamily, %DIF PREAMBLE
	columns=fullflexible, %DIF PREAMBLE
	keepspaces=true %DIF PREAMBLE
} %DIF PREAMBLE
\lstnewenvironment{DIFverbatim}{\lstset{style=DIFverbatimstyle}}{} %DIF PREAMBLE
\lstnewenvironment{DIFverbatim*}{\lstset{style=DIFverbatimstyle,showspaces=true}}{} %DIF PREAMBLE
%DIF END PREAMBLE EXTENSION ADDED BY LATEXDIFF

\begin{document}

\title[When can a prediction replace a trial?]{When can a prediction replace a physical trial? A production-referenced evaluation of design-stage decision rules in deep drawing}

\ifanon
\author*[1]{\fnm{Anonymised} \sur{for review}}
\affil*[1]{\orgname{Affiliation withheld for review}}
\else
\author*[1]{\fnm{H\"useyin Tayyer} \sur{Canseven}}\email{huseyin.canseven@lut.fi}

\affil*[1]{\orgdiv{Department of Electrical Engineering, LUT School of Energy Systems}, \orgname{LUT University}, \orgaddress{\street{Yliopistonkatu 34}, \city{Lappeenranta}, \postcode{53850}, \country{Finland}}}
\fi

%% BEGIN ABSTRACT
\abstract{\DIFdelbegin \DIFdel{Design-stage manufacturability decisions rest increasingly on simulation }\DIFdelend \DIFaddbegin \DIFadd{Simulation }\DIFaddend and machine learning \DIFaddbegin \DIFadd{inform design-stage manufacturability decisions}\DIFaddend , \DIFdelbegin \DIFdel{whose fitness to decide is rarely evaluated against production}\DIFdelend \DIFaddbegin \DIFadd{yet they are judged by prediction error, not by whether their verdicts could replace a physical trial}\DIFaddend . \DIFdelbegin \DIFdel{The proposed framework's }\DIFdelend \DIFaddbegin \DIFadd{This paper proposes a production-referenced evaluation of design-stage decision rules. Its }\DIFaddend unit, the \DIFdelbegin \DIFdel{production }\DIFdelend \DIFaddbegin \DIFadd{within-run }\DIFaddend floor, is the difference that drift and scatter of production and measurement produce between two batches of one design within a run. For each \DIFaddbegin \DIFadd{evidence }\DIFaddend relation of a new design to produced \DIFdelbegin \DIFdel{evidence}\DIFdelend \DIFaddbegin \DIFadd{ones (a sibling, a new process setting or a new family)}\DIFaddend , a rule \DIFdelbegin \DIFdel{returning }\DIFdelend \DIFaddbegin \DIFadd{that returns }\DIFaddend meets, fails or trial \DIFdelbegin \DIFdel{has a }\DIFdelend \DIFaddbegin \DIFadd{is characterised by its }\DIFaddend decisive distance\DIFdelbegin \DIFdel{in floors}\DIFdelend , beyond which at least 95\,\% of its verdicts are correct, and \DIFdelbegin \DIFdel{a }\DIFdelend \DIFaddbegin \DIFadd{its }\DIFaddend safe distance, beyond which at most 5\,\% are wrong; a \DIFdelbegin \DIFdel{procedure adds }\DIFdelend guard \DIFdelbegin \DIFdel{bands for }\DIFdelend \DIFaddbegin \DIFadd{band turns them into }\DIFaddend single verdicts. \DIFdelbegin \DIFdel{Fourteen rules were evaluated at 95\,\% conformance on open data of 18 deep-drawn alternatives of 500 parts, mostly }\DIFdelend \DIFaddbegin \DIFadd{An exploratory application to one laboratory deep-drawing }\DIFaddend process\DIFdelbegin \DIFdel{settings}\DIFdelend , with \DIFdelbegin \DIFdel{matched simulations (a draw-in floor: 0.05\,mm). With a }\DIFdelend \DIFaddbegin \DIFadd{alternatives }\DIFaddend produced \DIFdelbegin \DIFdel{sibling}\DIFdelend \DIFaddbegin \DIFadd{once and an uncalibrated simulation dataset}\DIFaddend , \DIFdelbegin \DIFdel{the }\DIFdelend \DIFaddbegin \DIFadd{scored simulation-based, learned and production-record rules. The }\DIFaddend nearest produced setting\DIFdelbegin \DIFdel{was right for requirements over }\DIFdelend \DIFaddbegin \DIFadd{, committed to in advance, decided at }\DIFaddend 3.4 floors \DIFdelbegin \DIFdel{from the truth (near-replicates 0.85, resolvable siblings 4.9). For }\DIFdelend \DIFaddbegin \DIFadd{with a sibling and 8.5 for }\DIFaddend a new \DIFaddbegin \DIFadd{process }\DIFaddend setting\DIFaddbegin \DIFadd{. No simulation-based or learned rule was more decisive; a Gaussian process specified for small data reduced to a force trend. Intervals calibrated on }\DIFaddend the \DIFdelbegin \DIFdel{best rules needed 4.7 floors interpolating}\DIFdelend \DIFaddbegin \DIFadd{same relation in the other family}\DIFaddend , \DIFdelbegin \DIFdel{4.9 extrapolating to 500\,kN and 8.9 to 100\,kN}\DIFdelend \DIFaddbegin \DIFadd{not learning}\DIFaddend , \DIFdelbegin \DIFdel{with a stroke-speed change; }\DIFdelend \DIFaddbegin \DIFadd{made the production record safe }\DIFaddend for \DIFaddbegin \DIFadd{new settings. For }\DIFaddend a new family \DIFdelbegin \DIFdel{the best was safe at 16--17 }\DIFdelend \DIFaddbegin \DIFadd{only simulation-based rules decided, none closer than 30 }\DIFaddend floors\DIFdelbegin \DIFdel{in either family's floor}\DIFdelend . \DIFdelbegin \DIFdel{Requirements at }\DIFdelend \DIFaddbegin \DIFadd{Within }\DIFaddend a \DIFdelbegin \DIFdel{performance index of 1.33 lay inside every decisive distance except }\DIFdelend \DIFaddbegin \DIFadd{family, }\DIFaddend a \DIFdelbegin \DIFdel{near-replicate's. With six }\DIFdelend \DIFaddbegin \DIFadd{production record indexed by its relation }\DIFaddend to \DIFdelbegin \DIFdel{nine produced alternatives at two or three forces, }\DIFdelend the \DIFdelbegin \DIFdel{simulation as used, offset or rescaled, helped only across families; learned 90\,\% intervals covered 85--96\,\% with }\DIFdelend \DIFaddbegin \DIFadd{coming design is the baseline }\DIFaddend a \DIFdelbegin \DIFdel{sibling, under 50\,\% extrapolating or across families. This gives AI-based }\DIFdelend \DIFaddbegin \DIFadd{learned }\DIFaddend manufacturability \DIFdelbegin \DIFdel{support a production-referenced evaluation standard}\DIFdelend \DIFaddbegin \DIFadd{check has to beat}\DIFaddend .}
%% END ABSTRACT

\DIFdelbegin %DIFDELCMD < \keywords{machine learning, design for manufacturing, predictive quality, uncertainty quantification, operating characteristic, deep drawing}
%DIFDELCMD < %%%
\DIFdelend \DIFaddbegin \keywords{decision support, design for manufacturing, machine learning, uncertainty quantification, selective prediction, deep drawing}
\DIFaddend 

\maketitle

%% ----------------------------------------------------------------------
\section{Introduction}\label{sec:intro}

Whether a design alternative can be made to its requirements is decided long before it is made. In sheet-metal forming the decision rests on a finite-element simulation \citep{banabic2010sheet}: the team compares the predicted characteristics with the requirements and releases the alternative, changes it or sends it to a physical trial. \DIFdelbegin \DIFdel{Simulation has moved such decisions forward \mbox{%DIFAUXCMD
\citep{thomke1998simulation}}\hskip0pt%DIFAUXCMD
, and artificial intelligence makes the }\DIFdelend \DIFaddbegin \DIFadd{Artificial intelligence makes this }\DIFaddend loop faster: surrogates replace simulations, models learn manufacturability from earlier and synthetic designs \DIFdelbegin \DIFdel{, checks are automated \mbox{%DIFAUXCMD
\citep{yan2026machine, ghadai2018learning, zhang2018featurenet, ballegeer2026bendfm}}\hskip0pt%DIFAUXCMD
}\DIFdelend \DIFaddbegin \DIFadd{\mbox{%DIFAUXCMD
\citep{yan2026machine, ghadai2018learning, ballegeer2026bendfm}}\hskip0pt%DIFAUXCMD
}\DIFaddend , and predictive-quality models and virtual metrology estimate \DIFdelbegin \DIFdel{product }\DIFdelend quality from process data \citep{tercan2022machine, dreyfus2022virtual}. Such models are judged by their prediction error on held-out data, which leaves open the question the team has to answer: \emph{can this prediction be relied on to decide\DIFaddbegin \DIFadd{, or is a trial still needed}\DIFaddend ?} Mechanical design cannot be verified by simulation the way integrated circuits are \citep{whitney1996why}, and \DIFdelbegin \DIFdel{production variation, with which development organisations still contend \mbox{%DIFAUXCMD
\citep{thornton2000more}}\hskip0pt%DIFAUXCMD
, adds a further reason: whether a prediction-based decision is credible is a property of production , and production evidence is scarce at the design stage. }%DIFDELCMD < 

%DIFDELCMD < %%%
\DIFdel{This paper treats the design-stage manufacturability decision itself as the object of study}\DIFdelend \DIFaddbegin \DIFadd{the answer depends on how much production has been seen and how the new design relates to it.
}

\DIFadd{Uncertainty quantification supplies intervals \mbox{%DIFAUXCMD
\citep{nemani2023uncertainty}}\hskip0pt%DIFAUXCMD
, and virtual metrology attaches a reliance index to each prediction \mbox{%DIFAUXCMD
\citep{cheng2008evaluating}}\hskip0pt%DIFAUXCMD
, but neither says how close to a requirement a verdict can be relied on, in a unit that production defines, for each relation of a new design to the produced evidence. This paper asks three questions}\DIFaddend . How can the fitness of a design-stage prediction to decide be expressed on a \DIFdelbegin \DIFdel{common }\DIFdelend production-referenced scale, for any predictor and characteristic (RQ1)? How does it depend on the amount of production evidence and on how the new design relates to \DIFdelbegin \DIFdel{what has been produced }\DIFdelend \DIFaddbegin \DIFadd{it }\DIFaddend (RQ2)? And\DIFdelbegin \DIFdel{what do simulation and machine learning contribute once production evidence exists }\DIFdelend \DIFaddbegin \DIFadd{, with six to nine produced alternatives, what do a generic simulation dataset, Gaussian processes, gradient boosting and calibrated intervals contribute against a production-record baseline }\DIFaddend (RQ3)?
\DIFdelbegin \DIFdel{Uncertainty quantification for machine learning supplies intervals \mbox{%DIFAUXCMD
\citep{nemani2023uncertainty}}\hskip0pt%DIFAUXCMD
, but AI-based decision support, decision-based design \mbox{%DIFAUXCMD
\citep{hazelrigg1998framework}}\hskip0pt%DIFAUXCMD
, design margins \mbox{%DIFAUXCMD
\citep{eckert2019design} }\hskip0pt%DIFAUXCMD
and test planning \mbox{%DIFAUXCMD
\citep{thomke2001sequential, salado2018mathematical, tahera2019testing} }\hskip0pt%DIFAUXCMD
need to know how close to a requirement a prediction can be relied on.
}\DIFdelend 

The \DIFdelbegin \DIFdel{framework proposed here has the }\DIFdelend \DIFaddbegin \DIFadd{proposed framework has the }\DIFaddend \emph{\DIFdelbegin \DIFdel{production }\DIFdelend \DIFaddbegin \DIFadd{within-run }\DIFaddend floor} as its unit: the difference that drift and scatter of production and measurement produce between two batches of one design within a run\DIFdelbegin \DIFdel{, a floor in the sense of a noise floor}\DIFdelend . A rule that turns the design-stage evidence into a verdict (\emph{meets}, \emph{fails} or a physical \emph{trial}) is \DIFdelbegin \DIFdel{characterised}\DIFdelend \DIFaddbegin \DIFadd{scored}\DIFaddend , for each \emph{evidence relation} between a new design and the produced alternatives, by two points on its operating \DIFdelbegin \DIFdel{characteristics in floors}\DIFdelend \DIFaddbegin \DIFadd{characteristic}\DIFaddend , its \emph{decisive} and \emph{safe distances}\DIFdelbegin \DIFdel{, and }\DIFdelend \DIFaddbegin \DIFadd{; }\DIFaddend a procedure turns them into a guard band for single verdicts \DIFaddbegin \DIFadd{and prices them by the costs of errors and trials}\DIFaddend . The decisions in scope are those of \DIFaddbegin \DIFadd{process planning, the approval of process changes in series production and the qualification of AI-based tools that would take part in them, for }\DIFaddend a team that has produced variants of a design family\DIFdelbegin \DIFdel{, in process planning and tryout and when production knowledge is carried to a new family}\DIFdelend \DIFaddbegin \DIFadd{; the tryout of a new die, in which the tool itself changes between loops, is not}\DIFaddend .

The \DIFdelbegin \DIFdel{evaluation }\DIFdelend \DIFaddbegin \DIFadd{application }\DIFaddend uses open data: 9,000 deep-drawn and cut cups of dual-phase steel, produced as 18 design--process alternatives of 500 consecutive parts each \DIFaddbegin \DIFadd{on one laboratory tool set }\DIFaddend \citep{baum2026rddac}, with finite-element simulations of the same process \DIFdelbegin \DIFdel{\mbox{%DIFAUXCMD
\citep{baum2025ddacs}}\hskip0pt%DIFAUXCMD
. Within each of two cup geometries the alternatives differ in blank-holder force and lubrication on one tool, so two of the three relations concern process settings, and the geometries give two transfers between families. Each alternative is judged as a new design by fourteen rules at a 95\,\% conformance level, against what its 500 parts did}\DIFdelend \DIFaddbegin \DIFadd{from a generic dataset \mbox{%DIFAUXCMD
\citep{baum2025ddacs}}\hskip0pt%DIFAUXCMD
. Each alternative was produced once, and the simulations were not calibrated to the process}\DIFaddend . The study \DIFaddbegin \DIFadd{is exploratory: it }\DIFaddend shows whether the framework discriminates between rules and evidence situations \DIFdelbegin \DIFdel{, not whether design teams decide better with it}\DIFdelend \DIFaddbegin \DIFadd{on real data}\DIFaddend , and its findings \DIFdelbegin \DIFdel{on one dataset }\DIFdelend are hypotheses for confirmation. The contributions are the framework (Section~\ref{sec:method})\DIFdelbegin \DIFdel{, }\DIFdelend \DIFaddbegin \DIFadd{; }\DIFaddend an evaluation design that separates evidence relations\DIFdelbegin \DIFdel{and }\DIFdelend \DIFaddbegin \DIFadd{, }\DIFaddend isolates what the simulation adds \DIFaddbegin \DIFadd{and tests the learners against a production-record baseline }\DIFaddend (Section~\ref{sec:demo})\DIFdelbegin \DIFdel{, empirical answers for deep drawing }\DIFdelend \DIFaddbegin \DIFadd{; findings for this process }\DIFaddend (Section~\ref{sec:results})\DIFdelbegin \DIFdel{, and guidance and a reporting standard for evaluating simulation- and AI-based manufacturability support, including classifiers and language-model checks }\DIFdelend \DIFaddbegin \DIFadd{; and a proposed reporting checklist with an open scoring function }\DIFaddend (Section~\DIFdelbegin \DIFdel{\ref{sec:discussion}}\DIFdelend \DIFaddbegin \DIFadd{\ref{sec:reuse}}\DIFaddend ).

%% ----------------------------------------------------------------------
\section{Background and related work}\label{sec:background}

\DIFdelbegin %DIFDELCMD < \subsection{Machine learning for manufacturability and quality}
%DIFDELCMD < %%%
\DIFdelend \DIFaddbegin \subsection{Machine learning for manufacturability, predictive quality and virtual metrology}
\DIFaddend 

Design for manufacturing \DIFdelbegin \DIFdel{(DFM) }\DIFdelend codifies what a process can make into rules \DIFdelbegin \DIFdel{, guidelines }\DIFdelend and cost models \citep{boothroyd2010product}\DIFdelbegin \DIFdel{, and quantitative manufacturability models have long driven design decisions \mbox{%DIFAUXCMD
\citep{latrobe2003design}}\hskip0pt%DIFAUXCMD
. Machine }\DIFdelend \DIFaddbegin \DIFadd{; machine }\DIFaddend learning learns manufacturability from CAD models \citep{ghadai2018learning, zhang2018featurenet}, \DIFdelbegin \DIFdel{also from }\DIFdelend synthetic sheet-metal datasets \citep{ballegeer2026bendfm}, \DIFdelbegin \DIFdel{from }\DIFdelend process data and simulated design spaces \citep{yan2026machine}\DIFdelbegin \DIFdel{, and generative models propose designs directly \mbox{%DIFAUXCMD
\citep{regenwetter2022deep}}\hskip0pt%DIFAUXCMD
}\DIFdelend . In production, predictive-quality models and virtual metrology estimate quality from process data \citep{wuest2016machine, tercan2022machine, dreyfus2022virtual}\DIFdelbegin \DIFdel{, their reproducibility and reuse have become research questions \mbox{%DIFAUXCMD
\citep{xie2025reproducible, xie2026reusability}}\hskip0pt%DIFAUXCMD
, and for deep drawing, }\DIFdelend \DIFaddbegin \DIFadd{; their uncertainty is quantified by Bayesian, ensemble and conformal methods \mbox{%DIFAUXCMD
\citep{nemani2023uncertainty, heddoub2026conformal}}\hskip0pt%DIFAUXCMD
. Automatic virtual metrology goes further and attaches to every prediction a reliance index and a global similarity index that flags process data unlike the training data \mbox{%DIFAUXCMD
\citep{cheng2008evaluating, cheng2012developing}}\hskip0pt%DIFAUXCMD
, the closest analogue in this literature to the guard band and the applicability guards of this paper. Models are carried across products and processes by transfer learning and domain adaptation \mbox{%DIFAUXCMD
\citep{maschler2021deep}}\hskip0pt%DIFAUXCMD
, also from simulation to real parts \mbox{%DIFAUXCMD
\citep{tercan2018transfer}}\hskip0pt%DIFAUXCMD
. In sheet forming, machine learning predicts defects and springback from simulated and measured data \mbox{%DIFAUXCMD
\citep{cao2024artificial, heinzelmann2025datadriven}}\hskip0pt%DIFAUXCMD
, and }\DIFaddend discrepancies between synthetic and real data are being addressed for \DIFdelbegin \DIFdel{AI applications \mbox{%DIFAUXCMD
\citep{baum2026addressing}}\hskip0pt%DIFAUXCMD
. Learned }\DIFdelend \DIFaddbegin \DIFadd{deep drawing \mbox{%DIFAUXCMD
\citep{baum2026addressing}}\hskip0pt%DIFAUXCMD
. These }\DIFaddend models are evaluated by held-out accuracy \DIFdelbegin \DIFdel{on labelled designs or parts}\DIFdelend \DIFaddbegin \DIFadd{or coverage}\DIFaddend , and language \DIFdelbegin \DIFdel{and vision-language }\DIFdelend models by case studies \citep{makatura2024how} and \DIFdelbegin \DIFdel{benchmark questions }\DIFdelend \DIFaddbegin \DIFadd{benchmarks }\DIFaddend \citep{picard2025concept}, not by the decisions they \DIFdelbegin \DIFdel{support }\DIFdelend \DIFaddbegin \DIFadd{would take }\DIFaddend against the production of the alternatives they are asked about.

\DIFdelbegin %DIFDELCMD < \subsection{Design decisions, design types and production evidence}
%DIFDELCMD < 

%DIFDELCMD < %%%
\DIFdel{Decision-based design casts design as a choice between alternatives with uncertain outcomes \mbox{%DIFAUXCMD
\citep{hazelrigg1998framework}}\hskip0pt%DIFAUXCMD
. How much a team knows about a new alternative depends on the kind of design, variant, adaptive or original \mbox{%DIFAUXCMD
\citep{pahl2007engineering}}\hskip0pt%DIFAUXCMD
, organised across products by platform design \mbox{%DIFAUXCMD
\citep{simpson2006product}}\hskip0pt%DIFAUXCMD
; the evidence relations of this paper are their counterparts for production evidence, the design type concerning the novelty of the solution and the evidence relation the novelty of its realisation relative to what has been produced. Production knowledge enters design through process-capability data \mbox{%DIFAUXCMD
\citep{thornton2004variation, tata1999process}}\hskip0pt%DIFAUXCMD
. Unresolved uncertainty is absorbed by margins \mbox{%DIFAUXCMD
\citep{eckert2019design}}\hskip0pt%DIFAUXCMD
, whose part that covers uncertainty is distinct from the part that absorbs change \mbox{%DIFAUXCMD
\citep{eckert2017safety}}\hskip0pt%DIFAUXCMD
. Tests and prototypes, the alternative to a margin, are modelled as a trade between the costs of tests and of errors \mbox{%DIFAUXCMD
\citep{thomke1998simulation, thomke2001sequential, salado2018mathematical} }\hskip0pt%DIFAUXCMD
and studied as design activities \mbox{%DIFAUXCMD
\citep{tahera2019testing}}\hskip0pt%DIFAUXCMD
. When upstream information is fit to act on is the question of overlapped development \mbox{%DIFAUXCMD
\citep{krishnan1997model, terwiesch2002exchanging}}\hskip0pt%DIFAUXCMD
, and set-based concurrent engineering codifies production knowledge as feasibility limits \mbox{%DIFAUXCMD
\citep{sobek1999toyota}}\hskip0pt%DIFAUXCMD
. All of these assume a measure of how far a design-stage prediction can be relied on; the framework supplies one in a unit set by production.
}%DIFDELCMD < 

%DIFDELCMD < \subsection{Simulation credibility, calibration and validation}
%DIFDELCMD < %%%
\DIFdelend \DIFaddbegin \subsection{Simulation credibility and the validation of forming simulation}
\DIFaddend 

Verification and validation ask whether a model is accurate enough for its intended use\DIFdelbegin \DIFdel{; the distinction between its validation and application domains \mbox{%DIFAUXCMD
\citep{oberkampf2010verification}}\hskip0pt%DIFAUXCMD
is the question that the evidence relation asks of a design, and risk-informed }\DIFdelend \DIFaddbegin \DIFadd{, distinguishing its validation from its application domain \mbox{%DIFAUXCMD
\citep{oberkampf2010verification}}\hskip0pt%DIFAUXCMD
; }\DIFaddend credibility assessment scales the required evidence with the consequence of a wrong decision \citep{asme2018assessing}. Bayesian calibration \DIFdelbegin \DIFdel{adds a discrepancy model and calibration parameters to a simulator \mbox{%DIFAUXCMD
\citep{kennedy2001bayesian}}\hskip0pt%DIFAUXCMD
, }\DIFdelend \DIFaddbegin \DIFadd{estimates calibration parameters together with a model discrepancy \mbox{%DIFAUXCMD
\citep{kennedy2001bayesian, higdon2004combining}}\hskip0pt%DIFAUXCMD
, without which calibrated parameters are biased and over-confident \mbox{%DIFAUXCMD
\citep{brynjarsdottir2014learning, arendt2012quantification}}\hskip0pt%DIFAUXCMD
; }\DIFaddend multi-fidelity models \DIFdelbegin \DIFdel{scale a cheap model to an expensive one \mbox{%DIFAUXCMD
\citep{kennedy2000predicting}}\hskip0pt%DIFAUXCMD
, and a model may be relied on only within its valid input domain \mbox{%DIFAUXCMD
\citep{malak2010using}}\hskip0pt%DIFAUXCMD
. Conformal prediction gives distribution-free intervals under exchangeability \mbox{%DIFAUXCMD
\citep{lei2018distribution, barber2021predictive}}\hskip0pt%DIFAUXCMD
, also under covariate shift \mbox{%DIFAUXCMD
\citep{tibshirani2019conformal} }\hskip0pt%DIFAUXCMD
and for decision risks \mbox{%DIFAUXCMD
\citep{angelopoulos2024conformal}}\hskip0pt%DIFAUXCMD
, and have entered quality assurance \mbox{%DIFAUXCMD
\citep{heddoub2026conformal} }\hskip0pt%DIFAUXCMD
and engineering design \mbox{%DIFAUXCMD
\citep{zong2026trustworthy}}\hskip0pt%DIFAUXCMD
. In sheet-metal forming, friction \mbox{%DIFAUXCMD
\citep{hol2012advanced} }\hskip0pt%DIFAUXCMD
and the unloading behaviour of }\DIFdelend \DIFaddbegin \DIFadd{combine a cheap and an expensive model through a fitted scale and a discrepancy \mbox{%DIFAUXCMD
\citep{kennedy2000predicting, legratiet2014recursive}}\hskip0pt%DIFAUXCMD
. In sheet forming, the friction law \mbox{%DIFAUXCMD
\citep{hol2012advanced}}\hskip0pt%DIFAUXCMD
, the unloading and hardening behaviour that govern springback \mbox{%DIFAUXCMD
\citep{yoshida2002model, wagoner2013advanced} }\hskip0pt%DIFAUXCMD
and the elastic deflection of the tools \mbox{%DIFAUXCMD
\citep{pilthammar2018introduction} }\hskip0pt%DIFAUXCMD
are }\DIFaddend the \DIFdelbegin \DIFdel{material \mbox{%DIFAUXCMD
\citep{yoshida2002model} }\hskip0pt%DIFAUXCMD
are }\DIFdelend main sources of simulation error\DIFdelbegin \DIFdel{, simulation datasets of deep drawing }\DIFdelend \DIFaddbegin \DIFadd{; robustness analysis propagates the scatter of material and process inputs \mbox{%DIFAUXCMD
\citep{wiebenga2014effect, chen2022robustness}}\hskip0pt%DIFAUXCMD
, and draw-in is measured and controlled in line in series production \mbox{%DIFAUXCMD
\citep{havinga2020exploiting, fischer2018implementation}}\hskip0pt%DIFAUXCMD
. Datasets of deep-drawing simulations }\DIFaddend support surrogate and transfer learning \citep{baum2025ddacs, heinzelmann2025comprehensive}, and the deviation between simulation and parts has been decomposed statistically \citep{baum2026statistical}\DIFdelbegin \DIFdel{. How }\DIFdelend \DIFaddbegin \DIFadd{; how }\DIFaddend close to a requirement \DIFdelbegin \DIFdel{such intervals and accuracy statements }\DIFdelend \DIFaddbegin \DIFadd{a calibrated or learned prediction }\DIFaddend can be relied on to decide \DIFdelbegin \DIFdel{, and how much production evidence that takes, }\DIFdelend remains open.

\DIFdelbegin %DIFDELCMD < \subsection{Decisions with abstention and their operating characteristics}
%DIFDELCMD < %%%
\DIFdelend \DIFaddbegin \subsection{Decisions with abstention and their evaluation}
\DIFaddend 

In conformity assessment a measured value proves conformance, proves non-conformance or leaves the decision open, with guard bands set by the measurement uncertainty \DIFdelbegin \DIFdel{and the risks of false acceptance and rejection as the measures of a decision rule \mbox{%DIFAUXCMD
\citep{iso2017decision}}\hskip0pt%DIFAUXCMD
. In selection procedures and acceptance sampling }\DIFdelend \DIFaddbegin \DIFadd{\mbox{%DIFAUXCMD
\citep{iso2017decision}}\hskip0pt%DIFAUXCMD
; the global consumer's and producer's risks follow from a prior distribution of the measurand \mbox{%DIFAUXCMD
\citep{jcgm2012evaluation}}\hskip0pt%DIFAUXCMD
. In acceptance sampling and selection procedures }\DIFaddend the operating characteristic gives the probability of a decision against the true quality, with an indifference zone inside which no decision is required \citep{bechhofer1954single}. A classifier with a reject option trades errors against abstentions \citep{chow1970optimum, geifman2017selective}, \DIFaddbegin \DIFadd{as does regression with a reject option \mbox{%DIFAUXCMD
\citep{zaoui2020regression}}\hskip0pt%DIFAUXCMD
; }\DIFaddend learning to defer passes difficult cases to an expert \citep{madras2018predict}\DIFaddbegin \DIFadd{. Cost-sensitive evaluation \mbox{%DIFAUXCMD
\citep{elkan2001foundations} }\hskip0pt%DIFAUXCMD
and decision curves \mbox{%DIFAUXCMD
\citep{vickers2006decision} }\hskip0pt%DIFAUXCMD
price the errors}\DIFaddend , and appropriate reliance on automation requires that its reliability be visible to its users \citep{lee2004trust}, \DIFdelbegin \DIFdel{also in design teams working with AI \mbox{%DIFAUXCMD
\citep{zhang2021cautionary}}\hskip0pt%DIFAUXCMD
}\DIFdelend \DIFaddbegin \DIFadd{though displayed confidence can calibrate trust without improving joint accuracy \mbox{%DIFAUXCMD
\citep{zhang2020effect, zhang2021cautionary}}\hskip0pt%DIFAUXCMD
}\DIFaddend . The decisive and safe distances are operating characteristics of this kind, and the \DIFdelbegin \DIFdel{procedure's }\DIFdelend guard band is the design-stage counterpart of a conformity-assessment one\DIFaddbegin \DIFadd{, whose dependence on a prior of requirement distances mirrors that of the global risks}\DIFaddend .

\DIFdelbegin %DIFDELCMD < \subsection{Production variation and measurement}
%DIFDELCMD < %%%
\DIFdelend \DIFaddbegin \subsection{Production variation, measurement and design decisions}
\DIFaddend 

Statistical quality control compares short- and long-term variation with tolerances through capability and performance indices\DIFdelbegin \DIFdel{\mbox{%DIFAUXCMD
\citep{montgomery2019introduction}}\hskip0pt%DIFAUXCMD
, with release criteria such as a performance index of at least }\DIFdelend \DIFaddbegin \DIFadd{, with recommended minima for existing processes \mbox{%DIFAUXCMD
\citep{montgomery2019introduction}}\hskip0pt%DIFAUXCMD
; in an initial process study an index above 1.67 meets the acceptance criteria and one between }\DIFaddend 1.33 \DIFdelbegin \DIFdel{\mbox{%DIFAUXCMD
\citep{aiag2006production} }\hskip0pt%DIFAUXCMD
}\DIFdelend \DIFaddbegin \DIFadd{and 1.67 may be acceptable after a review by the customer \mbox{%DIFAUXCMD
\citep{aiag2006production}}\hskip0pt%DIFAUXCMD
, }\DIFaddend and time-dependent models \DIFdelbegin \DIFdel{for }\DIFdelend \DIFaddbegin \DIFadd{describe }\DIFaddend drifting processes \citep{iso2026process}. Precision experiments define the repeatability limit, the difference that two results under repeatability conditions exceed with a probability of 5\,\% \citep{iso2025repeatability}, and intermediate precision under conditions that differ in time \DIFdelbegin \DIFdel{, operator or equipment }\DIFdelend \citep{iso2023intermediate}; measurement systems analysis separates repeatability from \DIFaddbegin \DIFadd{bias, linearity and }\DIFaddend stability \citep{aiag2010measurement}. \DIFdelbegin \DIFdel{The production floor is a precision limit of this kind for batch centres within a run (Section~\ref{sec:floor}); what is new is its use not to judge whether two results agree but as the unit in which a design decision rule is scored and its guard band set.
}\DIFdelend \DIFaddbegin \DIFadd{In engineering design, decisions between alternatives with uncertain outcomes \mbox{%DIFAUXCMD
\citep{hazelrigg1998framework} }\hskip0pt%DIFAUXCMD
are taken with margins \mbox{%DIFAUXCMD
\citep{eckert2019design} }\hskip0pt%DIFAUXCMD
or with tests whose cost is traded against that of errors \mbox{%DIFAUXCMD
\citep{thomke2001sequential, salado2018mathematical}}\hskip0pt%DIFAUXCMD
, and what a team knows about a new alternative depends on whether it is a variant, an adaptation or an original design \mbox{%DIFAUXCMD
\citep{pahl2007engineering}}\hskip0pt%DIFAUXCMD
, of which the evidence relations are the counterparts for production evidence.
}

\DIFaddend Table~\ref{tab:positioning} places the framework among \DIFdelbegin \DIFdel{them: it does not replace calibration or conformal prediction but scores them as decision rules against production in a common unit}\DIFdelend \DIFaddbegin \DIFadd{these methods. It adopts the indifference zone of operating characteristics, the guard band of conformity assessment, the selective risk of abstaining predictors and grouped cross-validation \mbox{%DIFAUXCMD
\citep{roberts2017cross}}\hskip0pt%DIFAUXCMD
. What is new is the production-defined unit in which every rule and characteristic is scored, the qualification by evidence relation, and their combination into a procedure for design-stage decisions that any predictor can enter}\DIFaddend .

\begin{table}[t]
\tabsize
\setlength{\tabcolsep}{3.5pt}
\caption{Positioning among methods that support or evaluate design-stage manufacturability decisions}\label{tab:positioning}
\DIFdelbeginFL %DIFDELCMD < \begin{tabular}{@{}L{3.0cm}L{2.5cm}L{3.0cm}L{1.9cm}L{1.6cm}@{}}
%DIFDELCMD < \toprule
%DIFDELCMD < Approach & Evidence & Output & Scored against production & Abstains \\
%DIFDELCMD < \midrule
%DIFDELCMD < DFM rules and cost models \citep{boothroyd2010product} & codified process knowledge & guideline compliance, cost & no & no \\
%DIFDELCMD < Design margins \citep{eckert2019design} & experience, analysis & buffer to the requirement & rarely & no \\
%DIFDELCMD < V\&V, credibility assessment \citep{oberkampf2010verification, asme2018assessing} & physics model, validation tests & accuracy in physical units; required evidence & on test cases & can require tests \\
%DIFDELCMD < Bayesian calibration \citep{kennedy2001bayesian} & simulations, experiments & predictive distribution & not as decisions & implicitly \\
%DIFDELCMD < Conformal prediction \citep{lei2018distribution} & exchangeable calibration data & interval with coverage & coverage; per decision possible & yes \\
%DIFDELCMD < Conformity assessment, OC curves \citep{iso2017decision, bechhofer1954single} & measurement and its uncertainty & accept, reject, undecided; risks & on the measured item & yes \\
%DIFDELCMD < Selective classification \citep{chow1970optimum, geifman2017selective} & labelled data & class or reject; risk--coverage & on held-out data & yes \\
%DIFDELCMD < ML manufacturability models \citep{ghadai2018learning, yan2026machine} & historical or simulated designs & class or value & on held-out data & rarely \\
%DIFDELCMD < This paper & simulations, produced alternatives & verdict; decisive and safe distance in floors by evidence relation & yes, per verdict & yes \\
%DIFDELCMD < \botrule
%DIFDELCMD < \end{tabular}
%DIFDELCMD < %%%
\DIFdelendFL \DIFaddbeginFL \begin{tabular}{@{}L{3.2cm}L{2.5cm}L{2.9cm}L{1.9cm}L{1.5cm}@{}}
\toprule
Approach & Evidence & Output & Scored against production & Abstains \\
\midrule
DFM rules and cost models \citep{boothroyd2010product} & codified process knowledge & guideline compliance, cost & no & no \\
V\&V, credibility assessment \citep{oberkampf2010verification, asme2018assessing} & physics model, validation tests & accuracy; required evidence & on test cases & can require tests \\
Bayesian calibration \citep{kennedy2001bayesian} & simulations, experiments & predictive distribution & not as decisions & implicitly \\
Conformal prediction \citep{lei2018distribution} & exchangeable calibration data & interval with coverage & coverage & yes \\
Virtual metrology with reliance index \citep{cheng2008evaluating} & process data of produced items & value, reliance and similarity index & on held-out production & flags \\
Conformity assessment, OC curves \citep{iso2017decision, bechhofer1954single} & measurement and its uncertainty & accept, reject, undecided; risks & on the measured item & yes \\
Selective classification and regression \citep{chow1970optimum, zaoui2020regression} & labelled data & value or reject; risk--coverage & on held-out data & yes \\
ML manufacturability models \citep{ghadai2018learning, yan2026machine} & historical or simulated designs & class or value & on held-out data & rarely \\
This paper & simulations, produced alternatives & verdict; distances in floors by evidence relation & yes, per verdict & yes \\
\botrule
\end{tabular}
\DIFaddendFL \end{table}

%% ----------------------------------------------------------------------
\section{The framework}\label{sec:method}

\DIFdelbegin %DIFDELCMD < \subsection{The decision problem and its scope}%%%
\DIFdelend \DIFaddbegin \subsection{The decision problem, its scope and its loss}\DIFaddend \label{sec:problem}

A design alternative $a$ is a combination of product geometry and process settings produced as a series of parts, over which a quality characteristic $c$ varies. A requirement is an upper limit $R$ on the characteristic, or on its absolute value for characteristics whose target is zero, that a share $p$ of the parts must meet: alternative $a$ is \emph{adequate} for $c$ when the $p$-quantile of its parts satisfies $q_p(a,c)\le R$. \DIFdelbegin \DIFdel{Lower or two-sided limits follow by a change of sign or by the absolute deviation from the target. }\DIFdelend At the design stage $q_p$ is unknown. A decision rule $r$ maps the available evidence to an interval $[L_r, U_r]$ for $q_p$ and to a verdict
\begin{equation}
V_r = \begin{cases} \text{meets} & U_r \le R,\\ \text{fails} & L_r > R,\\ \text{trial} & \text{otherwise.} \end{cases}
\label{eq:verdict}
\end{equation}
The three verdicts mirror the decision rules of conformity assessment \citep{iso2017decision}\DIFdelbegin \DIFdel{, here for the uncertainty of a design-stage prediction}\DIFdelend . A rule that returns a point ($L_r=U_r$) always decides. The evidence consists of simulations of $a$ and of the production record of other alternatives, the \emph{calibration alternatives}. The \DIFdelbegin \DIFdel{demonstration decides on }\DIFdelend \DIFaddbegin \DIFadd{verdict costs $c_\mathrm{FA}$ when a failing alternative is accepted, $c_\mathrm{FR}$ when an adequate one is rejected and $c_\mathrm{T}$ when it is sent to trial, so a rule's expected cost per verdict is
}\begin{equation}
\DIFadd{C = c_\mathrm{FA}P_\mathrm{FA} + c_\mathrm{FR}P_\mathrm{FR} + c_\mathrm{T}P_\mathrm{T},
\label{eq:cost}
}\end{equation}
\DIFadd{which the two distances of Section~\ref{sec:distances} summarise at a fixed share of errors. Three limits of scope apply: the demonstration uses upper limits on the 95th percentile (}\DIFaddend $p=0.95$, \DIFdelbegin \DIFdel{which }\DIFdelend for a normal characteristic \DIFdelbegin \DIFdel{corresponds to }\DIFdelend a performance index of about 0.55\DIFdelbegin \DIFdel{. Release criteria of series production, a performance index of 1.33 or more \mbox{%DIFAUXCMD
\citep{aiag2006production}}\hskip0pt%DIFAUXCMD
, concern nonconforming fractions in parts per million that a series of 500 parts cannot estimate without a parametric tail; Section~\ref{sec:res-requirements} locates such limits relative to the distances, without deciding them}\DIFdelend \DIFaddbegin \DIFadd{), lower limits following by a change of sign; verdicts are per characteristic, a design-level verdict over all characteristics being reported as an extension (Section~\ref{sec:res-robust}); and attribute outcomes such as splits and wrinkles, which need a requirement on a rate, are not covered}\DIFaddend .

\DIFdelbegin %DIFDELCMD < \subsection{Production floor, effect-to-scatter ratio and minimum resolvable change}%%%
\DIFdelend \DIFaddbegin \subsection{The within-run floor}\DIFaddend \label{sec:floor}

Production evidence comes as series of consecutive parts, each cut into batches of $B$ consecutive parts \DIFdelbegin \DIFdel{, and each batch is }\DIFdelend summarised by a robust centre $m$ (here the 10\,\% trimmed mean). The \DIFdelbegin \DIFdel{production }\DIFdelend \DIFaddbegin \DIFadd{within-run }\DIFaddend floor of characteristic $c$\DIFaddbegin \DIFadd{, the floor for short, }\DIFaddend is
\begin{equation}
F_c = Q_\alpha\bigl\{\,|m_i - m_j| : i< j \text{ batches of the same alternative}\,\bigr\},
\label{eq:floor}
\end{equation}
the $\alpha$-quantile of the difference between two batches of one design. Under a stationary normal model with a standard deviation $\sigma_\mathrm{b}$ between batch means and $\sigma_\mathrm{w}$ between parts, $F_c = 1.96\sqrt{2}\sqrt{\sigma_\mathrm{b}^2+\sigma_\mathrm{w}^2/B}$ at $\alpha=0.95$, the form of the repeatability limit \DIFdelbegin \DIFdel{of precision experiments \mbox{%DIFAUXCMD
\citep{iso2025repeatability}}\hskip0pt%DIFAUXCMD
. Because }\DIFdelend \DIFaddbegin \DIFadd{\mbox{%DIFAUXCMD
\citep{iso2025repeatability, iso1994accuracy}}\hskip0pt%DIFAUXCMD
; because }\DIFaddend it pools batch pairs up to 450 parts apart and grows with their lag\DIFdelbegin \DIFdel{(Section~\ref{sec:res-floor})}\DIFdelend , it is closer to a time-different intermediate-precision limit \citep{iso2023intermediate} of batch centres within a run\DIFdelbegin \DIFdel{; in the terms of ISO 22514-2 the runs follow a time-dependent model with varying location \mbox{%DIFAUXCMD
\citep{iso2026process}}\hskip0pt%DIFAUXCMD
}\DIFdelend . The reference protocol \DIFdelbegin \DIFdel{of this paper }\DIFdelend is $B=50$, ten batches of a 500-part run, all pairs of batches and $\alpha=0.95$\DIFdelbegin \DIFdel{, with per-geometry floors where geometries differ}\DIFdelend \DIFaddbegin \DIFadd{; every characteristic is scored in the floor of its own geometry, and the pooled floor serves only to convert distances to physical units}\DIFaddend . The floor rests on three assumptions: batches of one design differ only by drift and scatter within a run (A1); the variation within the runs from which it is estimated is like that within future runs (A2); and the measuring system is repeatable and its \DIFdelbegin \DIFdel{measurement }\DIFdelend resolution small against the floor (A3\DIFdelbegin \DIFdel{, checked in Section~\ref{sec:res-floor}}\DIFdelend ).

Four properties matter for its use. It has the unit of the characteristic, so distances in floors \DIFdelbegin \DIFdel{are on a common scale for }\DIFdelend \DIFaddbegin \DIFadd{put }\DIFaddend all characteristics and predictors \DIFaddbegin \DIFadd{on a common scale}\DIFaddend , though equal distances do not mean equal nonconforming fractions\DIFdelbegin \DIFdel{(Table~\ref{tab:floors})}\DIFdelend . It is a property of a protocol\DIFdelbegin \DIFdel{: under drift it grows with the time between the batches compared and shrinks with $B$}\DIFdelend , and $\sigma_\mathrm{b}$ and $\sigma_\mathrm{w}$ are reported so that it can be recomputed. It contains whatever drifts between batches, process and measuring system alike\DIFdelbegin \DIFdel{, so }\DIFdelend \DIFaddbegin \DIFadd{: }\DIFaddend without a reference artefact scanned through the series it is a production-and-measurement floor\DIFaddbegin \DIFadd{, and not the shop floor}\DIFaddend . And it excludes the variation between runs, which makes it a lower bound of the variation of series production. A new family, whose floor is unknown before it is produced, is scored in the floor of the produced family\DIFdelbegin \DIFdel{, and because the floor resolves batch centres, the reproducibility of the decided quantile is reported alongside. }%DIFDELCMD < 

%DIFDELCMD < %%%
\DIFdelend \DIFaddbegin \DIFadd{. }\DIFaddend Two alternatives $a$ and $b$ are resolvable for $c$ when their effect-to-scatter ratio \DIFdelbegin \begin{displaymath}
\DIFdel{\mathrm{ESR}_c(a,b) = |\bar m_a - \bar m_b| \,/\, F_c
%DIFDELCMD < \label{eq:esr}%%%
}\end{displaymath}%DIFAUXCMD
\DIFdelend \DIFaddbegin \DIFadd{$\mathrm{ESR}_c(a,b) = |\bar m_a - \bar m_b|/F_c$ }\DIFaddend is at least one, $\bar m$ being the centre of all parts of an alternative; \DIFdelbegin \DIFdel{this is a practical threshold, not a test. For a continuous design or process }\DIFdelend \DIFaddbegin \DIFadd{for a continuous }\DIFaddend parameter $x$ with \DIFdelbegin \DIFdel{a local sensitivity $s_c = \partial \bar m/\partial x$, }\DIFdelend \DIFaddbegin \DIFadd{local sensitivity $s_c=\partial\bar m/\partial x$ }\DIFaddend the minimum resolvable change is $F_c/|s_c|$\DIFdelbegin \DIFdel{: a change of $x$ by less moves the centre of the characteristic by less than two batches of one design differ. It is local and concerns centres; a change below it can still move the tail}\DIFdelend .

\subsection{Decision rules}\label{sec:rules}

\DIFdelbegin \DIFdel{Fourteen rules span what a design team can do with its evidence }\DIFdelend \DIFaddbegin \DIFadd{Eight rules carry the main results }\DIFaddend (Table~\ref{tab:rules}); \DIFdelbegin \DIFdel{the first seven and the force trend M1sn carry the main results. Beside the nominal simulation $s_0$ }\DIFdelend \DIFaddbegin \DIFadd{ten further rules, with and without the simulation, are defined in Appendix~\ref{app:appendix}. The nominal simulation }\DIFaddend (M0) \DIFdelbegin \DIFdel{, the bias-corrected simulation (M1) adds the median offset between the calibration alternatives' $q_p$ and their $s_0$; the }\DIFdelend \DIFaddbegin \DIFadd{is the simulation at nominal sheet thickness and friction. The }\DIFaddend scaled simulation (M1s) fits $q_p = a + b\,s_0$ to \DIFdelbegin \DIFdel{them }\DIFdelend \DIFaddbegin \DIFadd{the calibration alternatives }\DIFaddend by least squares, \DIFdelbegin \DIFdel{a linear }\DIFdelend \DIFaddbegin \DIFadd{the scale term of a }\DIFaddend multi-fidelity correction \DIFdelbegin \DIFdel{\mbox{%DIFAUXCMD
\citep{kennedy2000predicting}}\hskip0pt%DIFAUXCMD
}\DIFdelend \DIFaddbegin \DIFadd{without its discrepancy}\DIFaddend . The envelope rule (M2) adds the median offset \DIFdelbegin \DIFdel{and the largest absolute residual }\DIFdelend of the calibration alternatives \DIFdelbegin \DIFdel{to the upper end }\DIFdelend \DIFaddbegin \DIFadd{and their largest absolute residual to the top }\DIFaddend of the simulation envelope over \DIFdelbegin \DIFdel{plausible incoming conditions, which vary sheet thickness and friction but not material properties. }\DIFdelend \DIFaddbegin \DIFadd{the dataset's thicknesses and friction coefficients. The GP-corrected envelope (}\DIFaddend M5\DIFaddbegin \DIFadd{) }\DIFaddend models the offset between $q_p$ and the envelope as a Gaussian process \citep{rasmussen2006gaussian} \DIFdelbegin \DIFdel{with a 90\,\% predictive interval, whose noise variance is }\DIFdelend \DIFaddbegin \DIFadd{fitted by maximum marginal likelihood, its noise variance }\DIFaddend bounded below by the variance between the calibration \DIFdelbegin \DIFdel{alternatives' }\DIFdelend series of one geometry and force \DIFdelbegin \DIFdel{, so that a new series of a produced setting is not predicted more precisely than produced series reproduce each other; M5n is the same process of $q_p$ without the simulation}\DIFdelend \DIFaddbegin \DIFadd{(the noise-variance bound)}\DIFaddend . The nearest produced setting (NN) is \DIFaddbegin \DIFadd{the mean $q_p$ of the calibration alternatives at the nearest force, }\DIFaddend a capability-database look-up\DIFdelbegin \DIFdel{that uses neither simulation nor learning. The seven further rules include a friction-calibrated simulation \mbox{%DIFAUXCMD
\citep{kennedy2001bayesian} }\hskip0pt%DIFAUXCMD
(Mc), part-level gradient boosting \mbox{%DIFAUXCMD
\citep{friedman2001greedy} }\hskip0pt%DIFAUXCMD
with a jackknife+ interval \mbox{%DIFAUXCMD
\citep{barber2021predictive}}\hskip0pt%DIFAUXCMD
, with and without the simulation (M3}\DIFdelend , \DIFdelbegin \DIFdel{M3n) , and }\DIFdelend \DIFaddbegin \DIFadd{and the force trend (M1sn) }\DIFaddend the least-squares line of $q_p$ in force\DIFdelbegin \DIFdel{(M1sn). The nominal simulation and the envelope do not depend on the lubrication pattern, so }\DIFdelend \DIFaddbegin \DIFadd{, the counterpart of M1s without the simulation, which }\DIFaddend within a geometry \DIFdelbegin \DIFdel{they are functions }\DIFdelend \DIFaddbegin \DIFadd{is a function }\DIFaddend of the force\DIFdelbegin \DIFdel{alone: M1sn is the counterpart of M1s without the simulation, and a lubrication variant gives M0, M1, M1s, M2 and M5 no simulated information that distinguishes it from its siblings.
}\DIFdelend \DIFaddbegin \DIFadd{.
}\DIFaddend 

\DIFdelbegin \DIFdel{With $n$ calibration alternatives, the largest-residual margin covers a new exchangeable residual with probability $n/(n+1)$ for a fixed offset and at least $(n-1)/(n+1)$ for an estimated one; it is the split-conformal quantile \mbox{%DIFAUXCMD
\citep{lei2018distribution} }\hskip0pt%DIFAUXCMD
at the highest attainable level, not a }\DIFdelend \DIFaddbegin \DIFadd{The trend GP (M7n) is a learner specified for six to nine produced alternatives: a Gaussian process of $q_p$ with a linear mean in force and lubrication whose coefficients are integrated out (universal kriging), weakly informative priors on length scales and noise instead of hard bounds, and a noise term made of the known sampling variance of each $q_p$ and an estimated between-run nugget that absorbs no modelled effect. Its counterpart with the simulation (M7) is the autoregressive multi-fidelity model \mbox{%DIFAUXCMD
\citep{kennedy2000predicting}}\hskip0pt%DIFAUXCMD
, with the simulation in the mean and a GP discrepancy. Finally, intervals are separated from predictors: P+RM puts around the point prediction of a rule P a relation-matched interval, the }\DIFaddend 90\,\DIFdelbegin \DIFdel{\% interval. A distribution-free interval cannot certify a miss rate below $1/(n+1)$, so a miss rate of 5\,}\DIFdelend \% \DIFdelbegin \DIFdel{needs at least 19 produced alternatives in }\DIFdelend \DIFaddbegin \DIFadd{split-conformal quantile \mbox{%DIFAUXCMD
\citep{lei2018distribution} }\hskip0pt%DIFAUXCMD
of P's absolute errors in floors over the cases of }\DIFaddend the \DIFdelbegin \DIFdel{relevant relation. With six to nine calibration alternatives the jackknife}\DIFdelend \DIFaddbegin \DIFadd{same evidence relation in the other family, a group-conditional construction \mbox{%DIFAUXCMD
\citep{vovk2005algorithmic}}\hskip0pt%DIFAUXCMD
. The residuals are borrowed because a family's own calibration set holds no case of its new-setting relation; they guarantee coverage only if the two families' errors are exchangeable, which the data cannot establish \mbox{%DIFAUXCMD
\citep{barber2023conformal}}\hskip0pt%DIFAUXCMD
. NN}\DIFaddend +\DIFdelbegin \DIFdel{bounds at the level of 0.9 fall at or beyond the extreme leave-one-out values, which are used instead; under exchangeability these guarantee a coverage of $1-2/(n+1)$, 0.71 to 0.80 \mbox{%DIFAUXCMD
\citep{barber2021predictive}}\hskip0pt%DIFAUXCMD
, approximately because the residuals are centred on their median. The Gaussian-process interval is model-based. The level trades decisiveness for safety, as the cost of errors would set the threshold of a reject option \mbox{%DIFAUXCMD
\citep{chow1970optimum}}\hskip0pt%DIFAUXCMD
; 0.9 is used here, 0.8 and 0.95 in the robustness analysis}\DIFdelend \DIFaddbegin \DIFadd{RM is the production-record baseline with such an interval}\DIFaddend .

\begin{table}[t]
\tabsize
\setlength{\tabcolsep}{3.5pt}
\DIFdelbeginFL %DIFDELCMD < \caption{The decision rules; the first seven and M1sn carry the main results}%%%
\DIFdelendFL \DIFaddbeginFL \caption{The rules of the main results; ten further rules are defined in Table~\ref{tab:further}}\DIFaddendFL \label{tab:rules}
\DIFdelbeginFL %DIFDELCMD < \begin{tabular}{@{}L{3.5cm}L{4.9cm}L{1.9cm}L{1.4cm}@{}}
%DIFDELCMD < \toprule
%DIFDELCMD < Rule & Prediction of $q_p$ & Interval & Evidence \\
%DIFDELCMD < \midrule
%DIFDELCMD < M0 nominal simulation & simulation at nominal thickness and friction & point & none \\
%DIFDELCMD < M1 bias-corrected simulation & M0 plus the median offset & point & $q_p$ \\
%DIFDELCMD < M1s scaled simulation & $a+b\,$M0, least squares & point & $q_p$ \\
%DIFDELCMD < M2 envelope + residual margin & envelope top plus the median offset & largest residual & $q_p$ \\
%DIFDELCMD < M5 GP-corrected envelope & envelope plus a GP offset over the descriptors & GP, 90\,\% & $q_p$ \\
%DIFDELCMD < M5n GP of $q_p$ & GP of $q_p$ over the descriptors & GP, 90\,\% & $q_p$ \\
%DIFDELCMD < NN nearest produced setting & mean $q_p$ at the nearest produced setting & point & $q_p$ \\
%DIFDELCMD < \midrule
%DIFDELCMD < M0w process-window worst case & simulated range over the incoming conditions & range & none \\
%DIFDELCMD < Mc friction-calibrated simulation & friction calibrated on the draw-in, plus offset & point & centres, $q_p$ \\
%DIFDELCMD < M3 simulation + boosting & matched simulation plus boosted discrepancy & jackknife+ & parts \\
%DIFDELCMD < M3n boosting, no simulation & boosted model of the characteristic & jackknife+ & parts \\
%DIFDELCMD < M1n median of the produced & median $q_p$ of the calibration alternatives & point & $q_p$ \\
%DIFDELCMD < M1sn force trend & $a+b\,$force, least squares & point & $q_p$ \\
%DIFDELCMD < M2n median + residual margin & M1n with the largest-residual margin & largest residual & $q_p$ \\
%DIFDELCMD < \botrule
%DIFDELCMD < \end{tabular}
%DIFDELCMD < %%%
\DIFdelendFL \DIFaddbeginFL \begin{tabular}{@{}L{3.6cm}L{4.6cm}L{2.3cm}L{1.4cm}@{}}
\toprule
Rule & Prediction of $q_p$ & Interval (level) & Evidence \\
\midrule
M0 nominal simulation & simulation at nominal thickness and friction & point & none \\
M1s scaled simulation$^\dagger$ & $a+b\,$M0, least squares & point & $q_p$ \\
M2 envelope + residual margin & envelope top plus the median offset & largest residual ($\frac{n-1}{n+1}$) & $q_p$ \\
M5 GP-corrected envelope & envelope plus a GP offset, ML-II & GP (nominal 0.9) & $q_p$ \\
NN nearest produced setting$^\dagger$ & mean $q_p$ at the nearest produced force & point & $q_p$ \\
M1sn force trend$^\dagger$ & $a+b\,$force, least squares & point & $q_p$ \\
M7n trend GP$^\ddagger$ & linear mean plus GP, MAP under priors & GP (nominal 0.9) & $q_p$ \\
NN+RM$^\ddagger$ & NN & relation-matched (0.9) & $q_p$ \\
\botrule
\end{tabular}
\DIFaddendFL \footnotetext{Fitted on the calibration alternatives\DIFdelbeginFL \DIFdelFL{. Descriptors}\DIFdelendFL \DIFaddbeginFL \DIFaddFL{; descriptors}\DIFaddendFL : geometry indicator, blank-holder force and lubrication rank\DIFdelbeginFL \DIFdelFL{(the order of the oil amount)}\DIFdelendFL . \DIFdelbeginFL \DIFdelFL{GP}\DIFdelendFL \DIFaddbeginFL \DIFaddFL{Pre-specified unless marked}\DIFaddendFL : \DIFdelbeginFL \DIFdelFL{Gaussian process. M3 and M3n at part level; M3n is named M4 in }\DIFdelendFL \DIFaddbeginFL \DIFaddFL{$^\dagger$ added after }\DIFaddendFL the \DIFdelbeginFL \DIFdelFL{released code}\DIFdelendFL \DIFaddbeginFL \DIFaddFL{first results; $^\ddagger$ added for this version}\DIFaddendFL . \DIFdelbeginFL \DIFdelFL{M6}\DIFdelendFL \DIFaddbeginFL \DIFaddFL{Levels: guaranteed (largest residual}\DIFaddendFL , \DIFaddbeginFL \DIFaddFL{$n$ calibration alternatives: 8 with }\DIFaddendFL a \DIFdelbeginFL \DIFdelFL{conformalised GP}\DIFdelendFL \DIFaddbeginFL \DIFaddFL{sibling}\DIFaddendFL , \DIFdelbeginFL \DIFdelFL{is }\DIFdelendFL \DIFaddbeginFL \DIFaddFL{6 for }\DIFaddendFL a \DIFdelbeginFL \DIFdelFL{robustness variant (Table~\ref{tab:distances}}\DIFdelendFL \DIFaddbeginFL \DIFaddFL{new setting, 9 for a new family}\DIFaddendFL ) \DIFaddbeginFL \DIFaddFL{or nominal}\DIFaddendFL .}
\end{table}

\subsection{Operating characteristics: decisive and safe distances}\label{sec:distances}

A rule is scored by judging each produced alternative as a new design at requirements placed at a known distance from the truth,
\begin{equation}
R = q_p(a,c) + d\,F_c ,
\label{eq:requirement}
\end{equation}
so that $a$ is adequate exactly when $d\ge 0$. \DIFdelbegin \DIFdel{For a case $(a,c)$ with interval $[L,U]$, let $u=(U-q_p)/F_c$ and $v=(q_p-L)/F_c$ be the distances by which the interval excludes the truth from above and below. At $d=+\delta$ the verdict is correct when $u\le\delta$ and a false reject when $-v>\delta$; at $d=-\delta$ it is correct when $v<\delta$ and a false accept when $-u\ge\delta$. }\DIFdelend Let $\kappa(\delta)$ and $\omega(\delta)$ be the shares of correct and of wrong verdicts among all feasible verdicts at $|d|=\delta$: every case gives one verdict on an adequate alternative and, where \DIFdelbegin \DIFdel{$q_p-\delta F_c\ge0$}\DIFdelend \DIFaddbegin \DIFadd{the requirement is not negative}\DIFaddend , one on a failing alternative\DIFdelbegin \DIFdel{, and every verdict has equal weight. As $\delta$ grows, failing-side requirements become infeasible for characteristics close to zero, so the failing side carries half the weight up to a few floors and a third at about 100 floors (98 of 126 failing-side requirements remain feasible at 16 floors and 63 at 108)}\DIFdelend . The \emph{decisive distance} $\delta_\mathrm{d}$ \DIFdelbegin \DIFdel{of a rule }\DIFdelend is the smallest $\delta$ beyond which $\kappa\ge 0.95$, and \DIFdelbegin \DIFdel{its }\DIFdelend \DIFaddbegin \DIFadd{the }\DIFaddend \emph{safe distance} $\delta_\mathrm{s}$ the smallest $\delta$ beyond which $\omega\le 0.05$ \DIFdelbegin \DIFdel{.
}%DIFDELCMD < 

%DIFDELCMD < %%%
\DIFdel{The distances are exact up to the 0.05-floor grid on which they are evaluated: $\kappa$ and $\omega$ change only where $\delta$ crosses one of the case's $u$, $v$, $-u$ or $-v$, so the distances are order statistics of these exclusion distances; }\DIFdelend \DIFaddbegin \DIFadd{(Fig.~\ref{fig:schematic}a). In plain terms, }\DIFaddend for a rule that always decides, $\delta_\mathrm{d}=\delta_\mathrm{s}$ is close to the \DIFdelbegin \DIFdel{90\,\% quantile }\DIFdelend \DIFaddbegin \DIFadd{90th percentile }\DIFaddend of its absolute error in floors\DIFdelbegin \DIFdel{. Always $\delta_\mathrm{s}\le\delta_\mathrm{d}$}\DIFdelend \DIFaddbegin \DIFadd{; an interval rule trades decisiveness for safety}\DIFaddend , and between \DIFdelbegin \DIFdel{the two the rule }\DIFdelend \DIFaddbegin \DIFadd{$\delta_\mathrm{s}$ and $\delta_\mathrm{d}$ it }\DIFaddend sends designs to trial rather than misleading\DIFdelbegin \DIFdel{; one-sided distances for failing alternatives (false accepts) and adequate ones (false rejects) matter when the errors cost differently. The distances }\DIFdelend \DIFaddbegin \DIFadd{. Both }\DIFaddend are points on operating characteristics over a population of design situations\DIFdelbegin \DIFdel{: }\DIFdelend \DIFaddbegin \DIFadd{; }\DIFaddend $\delta_\mathrm{d}$ is the analogue of the indifference zone of a selection procedure \citep{bechhofer1954single} at a probability of correct selection of 0.95.
\DIFdelbegin \DIFdel{A trial is not error-free: a short run reproduces the 95th percentile }\DIFdelend \DIFaddbegin 

\DIFadd{The distances are exact up to the 0.05-floor grid on which they are evaluated, because $\kappa$ and $\omega$ change only where $\delta$ crosses the distance by which a case's interval excludes its truth. They are order statistics, so a few verdicts far out can move them: with 126 cases, some 240 verdicts, the 95\,\% threshold rests on about a dozen. Distances below about one floor are reported as ``$<$1'', since production reproduces the decided quantile }\DIFaddend only to about one floor within a run (Section~\ref{sec:res-floor})\DIFdelbegin \DIFdel{and, being a run, also carries the variation between runs, so no decision is reliable closer to the requirement than the trial replacing it}\DIFdelend \DIFaddbegin \DIFadd{. The distances weigh false accepts and false rejects equally, and as $\delta$ grows, failing-side requirements become infeasible for characteristics close to zero, so the failing side carries half the weight up to a few floors and a third at about 100 floors; the safe distances on the false-accept and the false-reject side are therefore reported as well}\DIFaddend .

\DIFdelbegin %DIFDELCMD < \subsection{Evidence relation, calibration budget and applicability guard}%%%
\DIFdelend \DIFaddbegin \begin{figure}[t]
\centering
\includegraphics[width=\textwidth]{F_schematic}
\caption{Operating characteristics of one rule (M7n, extrapolated setting): verdict shares against the requirement distance, with $\delta_\mathrm{d}$ and $\delta_\mathrm{s}$ (a), and correct share among decided verdicts against the margin, with $g$ (b)}\label{fig:schematic}
\end{figure}

\subsection{Evidence relation and calibration budget}\DIFaddend \label{sec:budget}

How a new design relates to the produced alternatives is part of the decision situation. With a \emph{sibling}, an alternative at the same process setting has been produced that differs only in the factor varied at a fixed setting (here the lubrication pattern\DIFdelbegin \DIFdel{), which production may not resolve}\DIFdelend \DIFaddbegin \DIFadd{, which the nominal simulation does not represent)}\DIFaddend ; a \emph{new process setting} lies between produced settings (interpolation) or outside them (extrapolation); a \emph{new family} shares no produced alternative. \DIFaddbegin \DIFadd{A team knows, before it decides, whether the produced siblings of a new variant differ resolvably from each other, so the sibling relation is stratified by this decision-time information, not by the new design's own production. }\DIFaddend The calibration budget is the dependence of the distances on the number $k$ of produced alternatives\DIFdelbegin \DIFdel{, over all subsets of size $k$ classified by relation}\DIFdelend . An applicability guard flags a \DIFdelbegin \DIFdel{new design that lies }\DIFdelend \DIFaddbegin \DIFadd{design }\DIFaddend outside the calibration\DIFaddbegin \DIFadd{, in the spirit of a valid input domain \mbox{%DIFAUXCMD
\citep{malak2010using}}\hskip0pt%DIFAUXCMD
}\DIFaddend : its geometry is not calibrated\DIFdelbegin \DIFdel{(family guard)}\DIFdelend , a descriptor lies outside the calibrated range\DIFdelbegin \DIFdel{(range guard)}\DIFdelend , or its simulated $q_p$ lies outside \DIFaddbegin \DIFadd{that of }\DIFaddend the calibration alternatives\DIFdelbegin \DIFdel{' simulated values (novelty guard), in the spirit of a valid input domain \mbox{%DIFAUXCMD
\citep{malak2010using}}\hskip0pt%DIFAUXCMD
}\DIFdelend .

\subsection{Procedure}\label{sec:procedure}

Figure~\ref{fig:procedure} turns the framework into six steps\DIFdelbegin \DIFdel{: }\DIFdelend \DIFaddbegin \DIFadd{, }\DIFaddend steps 1 to 5 once per design family \DIFdelbegin \DIFdel{, }\DIFdelend \DIFaddbegin \DIFadd{and }\DIFaddend repeated when production evidence is added, and step 6 for every new design. Steps 1 to 4 frame the decision (family, alternatives, characteristics, form of the requirement, \DIFdelbegin \DIFdel{conformance level}\DIFdelend \DIFaddbegin \DIFadd{costs}\DIFaddend ), measure the floor with a stated protocol, state which differences production resolves\DIFaddbegin \DIFadd{, }\DIFaddend and calibrate the candidate rules with and without the simulation.
\DIFaddbegin 

\DIFaddend Step 5 qualifies \DIFdelbegin \DIFdel{them }\DIFdelend \DIFaddbegin \DIFadd{the rules of the family }\DIFaddend by their distances for each evidence relation and fixes for each rule and relation a \emph{guard band} $g$: the smallest observable margin between the predicted interval and the requirement beyond which at least 95\,\% of the decided verdicts \DIFdelbegin \DIFdel{were }\DIFdelend \DIFaddbegin \DIFadd{are }\DIFaddend correct. Being conditional on what is observed, $g$ depends on the distribution of requirement distances that the team expects, \DIFdelbegin \DIFdel{not on rule and relation alone}\DIFdelend \DIFaddbegin \DIFadd{its prior, as the global risks of conformity assessment do \mbox{%DIFAUXCMD
\citep{jcgm2012evaluation}}\hskip0pt%DIFAUXCMD
}\DIFaddend . The reference prior weights the 47 distances of the \DIFdelbegin \DIFdel{requirement }\DIFdelend \DIFaddbegin \DIFadd{evaluation }\DIFaddend grid within $\pm20$ floors equally (0, $\pm0.5$ to $\pm10$ in steps of 0.5, $\pm12$, $\pm15$\DIFdelbegin \DIFdel{and }\DIFdelend \DIFaddbegin \DIFadd{, }\DIFaddend $\pm20$\DIFdelbegin \DIFdel{; 87\,\% of the weight within $\pm10$), drops infeasible requirements and searches $g$ in steps of 0.25 floors up to 20, beyond which the rules send most designs within the prior's range to trial; it is a reference for requirements set from }\DIFdelend \DIFaddbegin \DIFadd{), requirements set }\DIFaddend a few to some tens of floors from the true 95th percentile \DIFdelbegin \DIFdel{, denser where verdicts are hardest, which }\DIFdelend \DIFaddbegin \DIFadd{($\pm20$ floors are about $\pm1$\,mm of mid-side draw-in, $\pm1.6^\circ$ of wall angle, $\pm0.2$\,mm of cup depth and $\pm0.07$--0.16\,mm of dome); a release-like prior, the grid within $\pm5$ floors, stands for requirements near capability limits, and }\DIFaddend a team replaces \DIFaddbegin \DIFadd{both }\DIFaddend by the distribution of its own requirements\DIFdelbegin \DIFdel{. }\DIFdelend \DIFaddbegin \DIFadd{, for example the margins between past drawing limits and achieved 95th percentiles. The band is estimated per family and validated out of sample, leaving out one alternative and, for new settings, one force cell.
}

\DIFaddend Step 6 establishes the relation of the new design, applies the rule qualified for it, widens the rule's interval by $g$ floors and returns meets or fails with the margin in floors, or a trial when no rule is qualified or the requirement lies inside the widened interval. \DIFdelbegin \DIFdel{The guard }\DIFdelend \DIFaddbegin \DIFadd{In pseudocode, step 5 is: for each rule and relation, collect the held-out verdicts over the prior; for $g = 0, 0.25, \ldots, 20$, compute the share correct among verdicts decided with a margin of at least $g$; take the smallest $g$ beyond which it stays at or above 0.95. Step 6 is: $V = \text{meets}$ if $U+gF\le R$, fails if $L-gF>R$, trial otherwise. The }\DIFaddend band of the relation takes the place of the range and family guards, which here coincide with the extrapolated and new-family relations\DIFdelbegin \DIFdel{; the novelty guard is reported as a diagnostic, not applied, since most of its refusals would have been right; Section~\ref{sec:res-procedure} reports $g$ under other priors and works one decision through.
}\DIFdelend \DIFaddbegin \DIFadd{.
}\DIFaddend 

\begin{figure}[t]
\centering
\input{fig_procedure}
\DIFdelbeginFL %DIFDELCMD < \caption{The procedure: steps 1--5 once per design family, step 6 for every new design; dashed boxes are the evidence it uses}%%%
\DIFdelendFL \DIFaddbeginFL \caption{The procedure: steps 1--5 once per design family, step 6 for every new design; dashed boxes are the evidence and inputs it uses}\DIFaddendFL \label{fig:procedure}
\end{figure}

\DIFaddbegin \subsection{A worked example}\label{sec:example}

\DIFadd{A process planner has to decide whether the concave cup can be run at 300\,kN with medium lubrication, a setting not yet produced, without a trial run, the six alternatives at 100 and 500\,kN having been produced. Step 1 frames the decision: the characteristic is the mid-side draw-in, by which the blank-holder force is set, the requirement (chosen for illustration) that 95\,\% of the parts draw in by at most 15.30\,mm, and a trial costs a short press run, about a twentieth of a false accept. Step 2 gives the floor from the produced series, 0.050\,mm, and step 3 shows that force changes of 21--67\,kN are resolvable. Step 4 calibrates the rules on the six alternatives. Step 5 qualifies the nearest produced setting for interpolation: it decides at 4.7 floors and needs a guard band of 1.5 floors on this family (2.0 pooled over the families); out of sample, 95\,\% of its decided verdicts stayed correct. Step 6: the nearest setting, the mean of the six alternatives at the neighbouring forces, predicts 14.99\,mm, a margin of 6.2 floors, more than its band, so the planner releases the setting. The trend GP agrees with a margin of 4.5 floors; M2, M5 and M5n, widened by their bands, contain the limit and would send the design to trial, and the nominal simulation, qualified for no relation, predicts 15.66\,mm. The 500 parts had a 95th percentile of 14.88\,mm, 8.4 floors below the limit.
}

\DIFaddend %% ----------------------------------------------------------------------
\section{Demonstration on deep drawing}\label{sec:demo}

\subsection{Alternatives and production data}\label{sec:data}

The production data are the Real Deep Drawing and Cutting dataset (RDDAC) of the University of Stuttgart \citep{baum2026rddac, baum2026statistical}. Square blanks of DP600 \DIFaddbegin \DIFadd{dual-phase }\DIFaddend steel, 210\,mm wide and about 1\,mm thick, were deep drawn into cups (OP10) and cut (OP20) on \DIFaddbegin \DIFadd{an industrial press with }\DIFaddend one tool set\DIFdelbegin \DIFdel{with two }\DIFdelend \DIFaddbegin \DIFadd{, whose two die }\DIFaddend geometries, concave and convex side walls, \DIFdelbegin \DIFdel{which }\DIFdelend differ in base curvature, wall angle and bottom radius at once. Each geometry was run at three blank-holder forces (100, 300 and 500\,kN) and three lubrication patterns \DIFaddbegin \DIFadd{applied to the blank }\DIFaddend (coarse, medium and fine, with series medians of 1.15--1.48, 0.92--1.23 and 0.85--1.13\,g/m$^2$ of oil): 18 alternatives, each produced as one series of 500 consecutive parts \DIFdelbegin \DIFdel{, the change of geometry being the only change of product design}\DIFdelend \DIFaddbegin \DIFadd{(their centres in \esm{Fig.~S1})}\DIFaddend . Every part has \DIFdelbegin \DIFdel{press signals , sheet }\DIFdelend \DIFaddbegin \DIFadd{the signals of four load cells around the die at 300\,Hz, the punch position and temperature, sheet-thickness }\DIFaddend and oil-film traverses and \DIFdelbegin \DIFdel{3D }\DIFdelend \DIFaddbegin \DIFadd{laser }\DIFaddend scans after both operations. The series warm up by 1.0--5.0\,K over their first 150 parts from punch temperatures of 20.5--29.6\,$^\circ$C, their median sheet thickness is 0.986--0.993\,mm, and the 100\,kN series ran at 47\,mm/s against 78--79\,mm/s\DIFdelbegin \DIFdel{. The }\DIFdelend \DIFaddbegin \DIFadd{, so the extrapolation to 100\,kN also changes the sliding speed and stroke rate on which the lubricant film depends. The experiments were collected between September and November 2025, in contiguous blocks of 500 per alternative; the }\DIFaddend documentation states neither the order \DIFdelbegin \DIFdel{and dates of the series }\DIFdelend \DIFaddbegin \DIFadd{of the blocks }\DIFaddend nor whether the parts were scanned in production order\DIFaddbegin \DIFadd{, and gives no material certificate, coil record, tool drawing, press or cushion type, or method of applying the lubrication patterns}\DIFaddend .

\subsection{Simulations}\label{sec:sims}

The simulations come from the Deep Drawing and Cutting Simulations dataset (DDACS) \DIFaddbegin \DIFadd{\mbox{%DIFAUXCMD
\citep{baum2025ddacs, heinzelmann2025comprehensive}}\hskip0pt%DIFAUXCMD
, a generic dataset of some 32,000 LS-DYNA simulations }\DIFaddend of the same \DIFdelbegin \DIFdel{group \mbox{%DIFAUXCMD
\citep{baum2025ddacs, heinzelmann2025comprehensive}}\hskip0pt%DIFAUXCMD
: LS-DYNA quarter models with }\DIFdelend \DIFaddbegin \DIFadd{cup family that was neither set up for nor validated against the produced process. Its documentation gives rigid tools (blank holder, die, punch), shell elements with seven integration points through the thickness and }\DIFaddend nodes about 1\,mm apart, a scaled DP600 flow curve, constant Coulomb friction\DIFaddbegin \DIFadd{, a drawing depth of 30\,mm, trimming in OP20 without modelled cutting tools, }\DIFaddend and springback after both operations; \DIFaddbegin \DIFadd{it leaves open the element formulation, whether contact accounts for the shell thickness, }\DIFaddend the yield criterion and \DIFdelbegin \DIFdel{unloading model, which govern springback \mbox{%DIFAUXCMD
\citep{yoshida2002model}}\hskip0pt%DIFAUXCMD
, are not documented, and }\DIFdelend \DIFaddbegin \DIFadd{hardening law, the blank-holder actuation and the springback procedure, and it does not represent }\DIFaddend the stroke speed\DIFdelbegin \DIFdel{is not represented}\DIFdelend . For the RDDAC tools, 396 simulations vary the sheet thickness (0.95--1.00\,mm) and the friction coefficient (0.05--0.15) at each geometry and force\DIFaddbegin \DIFadd{, all with the nominal flow curve}\DIFaddend . A part is matched to the simulation with the nearest thickness and the friction coefficient that the dataset's linear mapping assigns to its oil film, an assumption of the dataset's authors on which \DIFdelbegin \DIFdel{Mc, }\DIFdelend M3\DIFdelbegin \DIFdel{and every matched simulation }\DIFdelend \DIFaddbegin \DIFadd{, M3n and the matched simulations }\DIFaddend depend. The nominal simulation of an alternative is at 0.99\,mm, the grid value nearest the median thickness, and friction 0.10\DIFaddbegin \DIFadd{, the middle of the grid; Mc and Mcd calibrate the friction coefficient instead}\DIFaddend .

\DIFdelbegin %DIFDELCMD < \subsection{Quality characteristics and their measurement}%%%
\DIFdelend \DIFaddbegin \subsection{Measurement and quality characteristics}\DIFaddend \label{sec:qcs}

\DIFaddbegin \DIFadd{The parts were measured with a single-sided laser line scanner, which records a height and a luminescence image of 3200 by 2000 points, 0.077\,mm apart along the laser line ($x$) and 0.158\,mm apart between lines in the direction of motion ($y$), with height steps of 7.75\,}\textmu \DIFadd{m, the cup bottom facing the scanner (Fig.~\ref{fig:parts}). The scales are not traced to the part: the lateral one comes from the scanner specification and a 30\,mm calibration block, the line spacing from the scans themselves, partly fitted to the simulation, and the height from a 1\,mm block, against a nominal repeatability of 5\,}\textmu \DIFadd{m. Absolute values are therefore uncalibrated, which limits every absolute comparison with the simulation (Section~\ref{sec:res-sim}) but not the floor, a within-alternative difference that a constant scale error changes only in proportion.
}

\DIFaddend Seven characteristics were extracted with the same definitions from the scans and \DIFaddbegin \DIFadd{from }\DIFaddend the simulated nodes, referred to the plane of the cup bottom \DIFaddbegin \DIFadd{(Table~\ref{tab:qcs})}\DIFaddend : the mid-side and corner draw-in and the flange waviness, process indicators on the flange that the cut removes\DIFdelbegin \DIFdel{, and four characteristics a drawing would carry: }\DIFdelend \DIFaddbegin \DIFadd{; }\DIFaddend the wall angle \DIFdelbegin \DIFdel{(design angle 20$^\circ$ concave, 10$^\circ$ convex), }\DIFdelend \DIFaddbegin \DIFadd{and }\DIFaddend the cup depth, \DIFaddbegin \DIFadd{measured after drawing because the cut removes the flange that defines the depth and the cut walls facing the scanner read several degrees too steep; }\DIFaddend and the arm angle and bottom dome of the cut part, the last two with requirements on absolute values. The four sides of a part measure several of them twice, which exposes \DIFdelbegin \DIFdel{the measurement}\DIFdelend \DIFaddbegin \DIFadd{direction-dependent systematic errors}\DIFaddend : the flange width \DIFdelbegin \DIFdel{along }\DIFdelend \DIFaddbegin \DIFadd{in the direction of motion ($W_y$) reads 2.3--3.5\,mm larger than along the laser line ($W_x$) and scatters 3.7 to 17.6 times more, }\DIFaddend the \DIFdelbegin \DIFdel{scan lines scatters 4 to 18 times morethan across them, so the mid-side draw-in is taken across the scan lines}\DIFdelend \DIFaddbegin \DIFadd{two flange diagonals differ by 1.6--2.2\,mm}\DIFaddend , and the \DIFdelbegin \DIFdel{wall angle is the mean of the east and west walls}\DIFdelend \DIFaddbegin \DIFadd{south wall, shadowed, reads 7--10$^\circ$ off. The directions kept avoid the worst of them, but their residual bias is unknown}\DIFaddend . The mid-side draw-in is quantised in steps of 0.019\,mm (0.38 floors)\DIFdelbegin \DIFdel{by the pixel pitch. Differences between scans and simulations, such as the scanned surface against the shell mid-surface, enter the offsets between the two but not the floors}\DIFdelend \DIFaddbegin \DIFadd{, so that a series takes only 6 to 21 distinct values (Fig.~\ref{fig:batches})}\DIFaddend .

%DIF > % BEGIN GENERATED qcs
\DIFaddbegin \begin{table}[t]
\tabsize
\setlength{\tabcolsep}{2.4pt}
\caption{The quality characteristics and their measurement by the laser scanner}\label{tab:qcs}
\begin{tabular}{@{}L{1.75cm}L{0.75cm}L{3.15cm}rrL{3.85cm}@{}}
\toprule
Characteristic & Op. & Region; definition & Step$/F$ & Values & Direction-dependent differences \\
\midrule
Draw-in, mid-side & OP10 & flange mid-sides; $(210-W_x)/2$, $W_x$ along the laser line & 0.38 & 6--21 & $W_y-W_x$ 2.29--3.47\,mm; $W_y$, in the direction of motion, scatters 3.7--17.6 times more \\
Draw-in, corner & OP10 & flange corners; mean of the two diagonals & $<$0.01 & 491--500 & the two diagonals differ by 1.64--2.17\,mm \\
Flange waviness & OP10 & flange; residual SD about a quadratic surface & $<$0.01 & 469--494 & -- \\
Wall angle & OP10 & side walls at mid-depth; mean of east and west & $<$0.01 & 497--500 & convex: north wall $-$1.64 to $-$0.92$^\circ$ from east/west (concave $-$0.05 to +0.39$^\circ$); south wall shadowed, 6.7--10.0$^\circ$ off; after the cut, north 3.3--5.6$^\circ$ off (convex) \\
Arm angle, cut & OP20 & cut arms, centre lines; mean of the four arms & $<$0.01 & 496--500 & -- \\
Cup depth & OP10 & flange above the bottom plane; median & $<$0.01 & 486--497 & the flange plane is cut away at OP20 \\
Bottom dome, cut & OP20 & bottom sag at $r=40$\,mm; paraboloid fit & $<$0.01 & 474--497 & -- \\
\botrule
\end{tabular}
\footnotetext{\DIFaddFL{Op.: operation measured (OP10 after drawing, OP20 after cutting). Step$/F$: median step between adjacent values in one series over the floor; Values: distinct values per series. Differences from the redundancy of the four sides, ranges of the series medians (\esm{Table~S20}); the directions used avoid them, their residual bias is unknown. Absolute values are not calibrated against a reference at the size of the part.}}
\end{table}
%DIF > % END GENERATED qcs

\begin{figure}[t]
\centering
\includegraphics[width=\textwidth]{F_parts}
\caption{Scanned parts after drawing and cutting, depth below the cup bottom (a, b, d, e), and sections (c, f); lines and boxes: directions and regions of the characteristics}\label{fig:parts}
\end{figure}

\DIFaddend \subsection{Evaluation design}\label{sec:evaldesign}

Every produced alternative is judged as a new design under grouped cross-validation \citep{roberts2017cross}, one scheme per evidence relation. \emph{New variant}: the alternative alone is held out and the other eight of its geometry calibrate, two siblings at its force among them. \emph{New process setting}: the three alternatives of one force of a geometry are held out and the six others calibrate; 300\,kN is interpolated, 100 and 500\,kN extrapolated. \emph{New family}: one geometry calibrates the rules for the other (two transfers)\DIFdelbegin \DIFdel{. Pooled versions of the first two add the other geometry's alternatives to the calibration set}\DIFdelend \DIFaddbegin \DIFadd{; a few-shot variant adds one to three produced alternatives of the new family}\DIFaddend . Every calibrated rule is evaluated with and without the simulation (\DIFdelbegin \DIFdel{M1/M1n, }\DIFdelend M1s/M1sn, M2/M2n, M3/M3n, M5/M5n\DIFdelbegin \DIFdel{), isolating what the simulation contributes, and the nearest produced setting is the baseline without simulation and learning. A scope }\DIFdelend \DIFaddbegin \DIFadd{, M7/M7n). A relation }\DIFaddend comprises 126 cases (18 alternatives, 7 characteristics)\DIFaddbegin \DIFadd{, but the units of replication are the 18 alternatives, and the characteristics of one alternative are correlated}\DIFaddend .

The \DIFdelbegin \DIFdel{intervals }\DIFdelend \DIFaddbegin \DIFadd{uncertainty }\DIFaddend of the distances \DIFdelbegin \DIFdel{come from 1000 resamples of }\DIFdelend \DIFaddbegin \DIFadd{is described in three ways. The intervals of Table~\ref{tab:distances} resample }\DIFaddend the held-out alternatives within each geometry \DIFdelbegin \DIFdel{, each }\DIFdelend \DIFaddbegin \DIFadd{1000 times (those held out together for a new setting individually), each resample }\DIFaddend with a block-bootstrap replicate of every true $q_p$ \DIFdelbegin \DIFdel{, }\DIFdelend \DIFaddbegin \DIFadd{(non-overlapping blocks of 50 parts), the fitted rules held fixed; }\DIFaddend common to all rules\DIFdelbegin \DIFdel{so that differences get paired intervals ; they hold the fitted rules fixed. A second bootstrap refits every alternative-level rule after drawing }\DIFdelend \DIFaddbegin \DIFadd{, they give paired intervals of differences, descriptive and not adjusted for the many comparisons. A refitting bootstrap redraws }\DIFaddend the three lubrication patterns of each geometry with replacement \DIFdelbegin \DIFdel{, at least two of them distinct, which keeps every case in its evidence relation (Table~\ref{tab:robust}). Its 200 resamples contain 48 of the 49 distinct configurations; in three of four draws a pattern is duplicated, so that a new variant keeps one distinct sibling and five distinct calibration alternatives, and the duplicated seriesenter the noise floor of }\DIFdelend \DIFaddbegin \DIFadd{and refits every alternative-level rule, which shows the effect of losing a pattern rather than sampling uncertainty. Coverage gets cluster-bootstrap intervals over alternatives. Each truth is }\DIFaddend the \DIFdelbegin \DIFdel{Gaussian processes as identical pairs. Its intervals thus show the effect of losing a pattern, not sampling uncertainty}\DIFdelend \DIFaddbegin \DIFadd{95th percentile of a single run, so a sensitivity analysis adds to every truth a between-run component of standard deviation $\tau$ floors, from near-replicate reproducibility up to the variation between series. The selection of rules and bands is cross-fitted between the geometries, chosen on one and evaluated on the other; like any cross-validated estimate, the distances describe the fitting procedure on this design rather than one fitted rule \mbox{%DIFAUXCMD
\citep{bates2024cross}}\hskip0pt%DIFAUXCMD
}\DIFaddend .

\subsection{Computation and development record}\label{sec:computation}

The constants of the floor and \DIFaddbegin \DIFadd{of }\DIFaddend the requirement grid were \DIFdelbegin \DIFdel{fixed }\DIFdelend \DIFaddbegin \DIFadd{set }\DIFaddend before the decision results were seen, \DIFdelbegin \DIFdel{as the repository history records}\DIFdelend \DIFaddbegin \DIFadd{which the archive cannot verify, as no timestamped snapshot was deposited at the time}\DIFaddend . The rules were not \DIFaddbegin \DIFadd{fixed in advance}\DIFaddend : M0--M5 \DIFaddbegin \DIFadd{and M3n }\DIFaddend were defined first\DIFdelbegin \DIFdel{; }\DIFdelend \DIFaddbegin \DIFadd{, and everything else was added after the first results, most of it in response to reviews of earlier versions simulated with a large language model in four rounds (}\DIFaddend M0w, Mc, M1n, NN, M2n, M5n, the grouped schemes, the exact estimator\DIFdelbegin \DIFdel{and the noise floor of the Gaussian processes were added after the first results; an internal review of that version, simulated with a large language model, added }\DIFdelend \DIFaddbegin \DIFadd{, the noise-variance bound, }\DIFaddend M1s, \DIFaddbegin \DIFadd{M1sn, }\DIFaddend the jackknife+ intervals, the \DIFdelbegin \DIFdel{between-run noise floor and the analyses of Sections~\ref{sec:res-drivers} and \ref{sec:res-procedure}; and a second such review added M1sn, the guard bands under other }\DIFdelend \DIFaddbegin \DIFadd{guard bands with their }\DIFaddend priors and the new family \DIFaddbegin \DIFadd{scored }\DIFaddend in the produced family's floor\DIFaddbegin \DIFadd{). A fifth simulated round, of the version prepared for this journal, led to M7n, M7, Mcd, the relation-matched intervals, the decision-time strata, the per-family and out-of-sample guard bands, the costs, the between-run sensitivity, the cross-fitting, the few-shot transfer and the design-level verdicts}\DIFaddend . The comparison of rules is therefore exploratory and carries the optimism of a selection made on the evaluation data \DIFdelbegin \DIFdel{\mbox{%DIFAUXCMD
\citep{cawley2010over}}\hskip0pt%DIFAUXCMD
}\DIFdelend \DIFaddbegin \DIFadd{\mbox{%DIFAUXCMD
\citep{cawley2010over, varma2006bias}}\hskip0pt%DIFAUXCMD
, which the cross-fitting bounds}\DIFaddend ; every variant tried is reported in the paper\DIFdelbegin \DIFdel{or in the archived results \mbox{%DIFAUXCMD
\citep{canseven2026archive}}\hskip0pt%DIFAUXCMD
, which }\DIFdelend \DIFaddbegin \DIFadd{, Online Resource~1 or the archive \mbox{%DIFAUXCMD
\citep{canseven2026archive}}\hskip0pt%DIFAUXCMD
, mapped by }\DIFaddend Table~\ref{tab:mapping}\DIFdelbegin \DIFdel{maps to the text. A large language model (Claude, Anthropic) assisted under the author's direction with code development, literature retrieval, data analysis, drafting and these reviews; the author designed the study and verified every analysis and reference}\DIFdelend .

%% ----------------------------------------------------------------------
\section{Results}\label{sec:results}

\DIFdelbegin %DIFDELCMD < \subsection{What the production floor contains}%%%
\DIFdelend \DIFaddbegin \subsection{What the floor contains}\DIFaddend \label{sec:res-floor}

\DIFdelbegin \DIFdel{In 95\,\% of batch pairs, two }\DIFdelend \DIFaddbegin \DIFadd{Two }\DIFaddend batches of 50 consecutive parts of one alternative differ by up to 0.051\,mm in mid-side draw-in and 0.078$^\circ$ in wall angle \DIFaddbegin \DIFadd{in 95\,\% of batch pairs }\DIFaddend (Table~\ref{tab:floors})\DIFdelbegin \DIFdel{; }\DIFdelend \DIFaddbegin \DIFadd{, so }\DIFaddend one floor is 0.7 to 1.7 long-term standard deviations of the parts. The floors are set by drift within the run \DIFaddbegin \DIFadd{(Fig.~\ref{fig:batches})}\DIFaddend : without the time order they shrink by a factor of 2.1 to 4.5, and batches one position apart differ by 0.51--0.80 floors\DIFdelbegin \DIFdel{and }\DIFdelend \DIFaddbegin \DIFadd{, }\DIFaddend nine apart by 1.1--1.5\DIFdelbegin \DIFdel{(averages over the two geometries). The normal model of Section~\ref{sec:floor} reproduces the floors to within 12\,\%; for subgroups of five parts, as sampled inspection would measure, it gives floors 1.1 }\DIFdelend \DIFaddbegin \DIFadd{. They vary between runs: the floor of a single run is 0.21 }\DIFaddend to \DIFdelbegin \DIFdel{2.1 times as large (median 1.3)}\DIFdelend \DIFaddbegin \DIFadd{2.16 times the pooled floor of its geometry, and the alternative-cluster interval is wider than the batch interval. Without the warm-up of the first 150 parts the floors fall to 0.75--1.00 of the reference, which keeps the warm-up because series production starts up after every break}\DIFaddend . Independent part-to-part noise, scanner repeatability included, \DIFdelbegin \DIFdel{would produce }\DIFdelend \DIFaddbegin \DIFadd{accounts for }\DIFaddend 15--33\,\% of each floor, and the \DIFdelbegin \DIFdel{measurement }\DIFdelend resolution is small against \DIFdelbegin \DIFdel{the floor }\DIFdelend \DIFaddbegin \DIFadd{it }\DIFaddend except for the mid-side draw-in\DIFdelbegin \DIFdel{; }\DIFdelend \DIFaddbegin \DIFadd{, so }\DIFaddend A3 holds for repeatability and \DIFdelbegin \DIFdel{measurement resolution, but the stability of the }\DIFdelend \DIFaddbegin \DIFadd{resolution, while bias, linearity and stability remain untested.
}

\DIFadd{Physical candidates for the drift are the warm-up of the tools and its effect on the lubricant, the build-up of lubricant on die and blank holder over consecutive strokes, and variation of the material and of the blank size along the stack, which enters the draw-in one to one because it is computed against the nominal 210\,mm blank. Whether the drift is the process or the }\DIFaddend scanner cannot be \DIFdelbegin \DIFdel{checked }\DIFdelend \DIFaddbegin \DIFadd{settled }\DIFaddend without a reference artefact\DIFdelbegin \DIFdel{. Nor do the process signals identify the source of the drift}\DIFdelend \DIFaddbegin \DIFadd{, but two checks bear on it. No recorded process signal explains it}\DIFaddend : batch centres of press force, punch temperature, sheet thickness and oil film predict it out of fold with $R^2$ of $-0.23$ to $0.20$. \DIFaddbegin \DIFadd{And the batch centres of the flange width and of the cup width at 15\,mm depth, which punch and die largely fix, are weakly correlated within the runs ($r=-0.27$ and $+0.27$), which argues against a common lateral scale drift of the scanner; the cup width varies by 0.06--0.07\,mm between batches, so it bounds such a drift only below 0.05\,\%, and drift in height and of the edge detection remain untestable. }\DIFaddend Variation between runs is not in the floor: the three lubrication series of one geometry and force differ by \DIFdelbegin \DIFdel{up to }\DIFdelend 1.2 (waviness) to 9.8 floors (dome) \DIFdelbegin \DIFdel{, and a floor taken across them is 1.5 to 9.1 times the reference. Within }\DIFdelend \DIFaddbegin \DIFadd{at the 95th percentile of their centre differences, lubrication effects included, whereas within }\DIFaddend a run the \DIFdelbegin \DIFdel{decided statistic reproduces to about one floor (the }\DIFdelend 95th percentiles of two batches of 100 parts differ by 0.92--1.49 floors in 95\,\% of pairs\DIFdelbegin \DIFdel{), so distances below about one floor are beyond what production can confirm, and a trial run, which also carries the variation between runs, confirms less still}\DIFdelend .

%% BEGIN GENERATED floors
\begin{table}[t]
\tabsize
\DIFdelbeginFL %DIFDELCMD < \setlength{\tabcolsep}{2.0pt}
%DIFDELCMD < \caption{Production floor per characteristic, pooled and per geometry, with its components}%%%
\DIFdelendFL \DIFaddbeginFL \setlength{\tabcolsep}{3pt}
\caption{Within-run floor per characteristic, pooled and per geometry, with its spread over runs}\DIFaddendFL \label{tab:floors}
\DIFdelbeginFL %DIFDELCMD < \begin{tabular}{@{}lrrrrrrrr@{}}
%DIFDELCMD < \toprule
%DIFDELCMD < Characteristic & $F$ & 95\,\% interval & Concave & Convex & $F/\sigma_\mathrm{LT}$ & $\sigma_\mathrm{b}/F$ & $\sigma_\mathrm{w}/F$ & Drift \\
%DIFDELCMD < \midrule
%DIFDELCMD < Draw-in, mid-side (mm) & 0.051 & 0.045--0.061 & 0.050 & 0.052 & 1.0/1.4 & 0.32/0.34 & 0.91/0.66 & 3.3 \\
%DIFDELCMD < Draw-in, corner (mm) & 0.046 & 0.035--0.052 & 0.053 & 0.030 & 1.4/1.2 & 0.31/0.35 & 0.65/0.77 & 4.5 \\
%DIFDELCMD < Flange waviness (mm) & 0.0028 & 0.0024--0.0031 & 0.0037 & 0.0021 & 1.6/0.7 & 0.33/0.30 & 0.55/1.31 & 2.5 \\
%DIFDELCMD < Wall angle ($^\circ$) & 0.078 & 0.066--0.10 & 0.081 & 0.078 & 1.3/0.7 & 0.32/0.26 & 0.68/1.39 & 2.1 \\
%DIFDELCMD < Arm angle, cut ($^\circ$) & 0.11 & 0.094--0.12 & 0.12 & 0.100 & 1.7/0.8 & 0.36/0.31 & 0.46/1.22 & 2.6 \\
%DIFDELCMD < Cup depth (mm) & 0.011 & 0.0096--0.012 & 0.0087 & 0.012 & 1.4/1.6 & 0.31/0.34 & 0.65/0.52 & 3.9 \\
%DIFDELCMD < Bottom dome, cut (mm) & 0.0070 & 0.0061--0.0081 & 0.0034 & 0.0082 & 1.2/1.6 & 0.33/0.36 & 0.76/0.55 & 3.9 \\
%DIFDELCMD < \botrule
%DIFDELCMD < \end{tabular}
%DIFDELCMD < %%%
\DIFdelendFL \DIFaddbeginFL \begin{tabular*}{\textwidth}{@{\extracolsep{\fill}}lrrrrrr@{}}
\toprule
Characteristic & $F$ & Interval & Concave/convex & Single runs & No warm-up & $F/\sigma_\mathrm{LT}$ \\
\midrule
Mid-side (mm) & 0.051 & 0.042--0.061 & 0.050/0.052 & 0.36--1.62 & 0.99 & 1.0/1.4 \\
Corner (mm) & 0.046 & 0.032--0.053 & 0.053/0.030 & 0.27--1.62 & 0.95 & 1.4/1.2 \\
Wavin. (mm) & 0.0028 & 0.0021--0.0050 & 0.0037/0.0021 & 0.21--2.08 & 0.89 & 1.6/0.7 \\
Wall ($^\circ$) & 0.078 & 0.054--0.10 & 0.081/0.078 & 0.30--1.74 & 0.80 & 1.3/0.7 \\
Arm ($^\circ$) & 0.11 & 0.089--0.15 & 0.12/0.100 & 0.40--2.16 & 1.00 & 1.7/0.8 \\
Depth (mm) & 0.011 & 0.0090--0.012 & 0.0087/0.012 & 0.39--1.45 & 0.75 & 1.4/1.6 \\
Dome (mm) & 0.0070 & 0.0052--0.0084 & 0.0034/0.0082 & 0.42--1.61 & 0.93 & 1.2/1.6 \\
\botrule
\end{tabular*}
\DIFaddendFL \footnotetext{Batches of 50 consecutive parts, 95\,\% quantile of the difference of their 10\,\% trimmed means\DIFdelbeginFL \DIFdelFL{; interval from }\DIFdelendFL \DIFaddbeginFL \DIFaddFL{. Interval: 95\,\% over }\DIFaddendFL 2000 \DIFdelbeginFL \DIFdelFL{batch }\DIFdelendFL resamples \DIFdelbeginFL \DIFdelFL{. Ratios per geometry }\DIFdelendFL \DIFaddbeginFL \DIFaddFL{of whole alternatives }\DIFaddendFL (\DIFdelbeginFL \DIFdelFL{concave/convex}\DIFdelendFL \DIFaddbeginFL \DIFaddFL{the narrower interval over resampled batches is in the archive}\DIFaddendFL )\DIFaddbeginFL \DIFaddFL{. Single runs: floor of one run over that of its geometry}\DIFaddendFL , \DIFdelbeginFL \DIFdelFL{classical statistics pooled }\DIFdelendFL \DIFaddbeginFL \DIFaddFL{range }\DIFaddendFL over the \DIFdelbeginFL \DIFdelFL{series as root mean squares}\DIFdelendFL \DIFaddbeginFL \DIFaddFL{18 runs. No warm-up}\DIFaddendFL : \DIFaddbeginFL \DIFaddFL{floor without the first 150 parts over the reference. }\DIFaddendFL $\sigma_\mathrm{LT}$\DIFdelbeginFL \DIFdelFL{, }\DIFdelendFL \DIFaddbeginFL \DIFaddFL{: }\DIFaddendFL standard deviation of the parts of a series \DIFaddbeginFL \DIFaddFL{(concave/convex)}\DIFaddendFL ; $\sigma_\mathrm{b}$, \DIFdelbeginFL \DIFdelFL{between batch means, corrected for $\sigma_\mathrm{w}^2/50$; }\DIFdelendFL $\sigma_\mathrm{w}$ \DIFdelbeginFL \DIFdelFL{, within batches. Drift: floor over that of randomly permuted series}\DIFdelendFL \DIFaddbeginFL \DIFaddFL{and the drift in \esm{Table~S20}}\DIFaddendFL .}
\end{table}
%% END GENERATED floors

\DIFdelbegin %DIFDELCMD < \subsection{What production resolves and how far the simulation misses it}%%%
\DIFdelend \DIFaddbegin \begin{figure}[t]
\centering
\includegraphics[width=\textwidth]{F_batches}
\caption{The 500 parts of one run in production order with the centres of its ten batches of 50 (bars) and the floor: mid-side draw-in (a) and wall angle (b)}\label{fig:batches}
\end{figure}

\subsection{What production resolves and where the simulation dataset lies}\DIFaddend \label{sec:res-sim}

All 63 geometry contrasts exceed the floor and blank-holder force is resolved in 105 of 126 contrasts, but lubrication only in 51 of 126, so a lubrication variant is often a near-replicate of its siblings. The minimum resolvable force change is 21--67\,kN between 100 and 300\,kN for most characteristics and lies beyond either force interval for the wall angle and the dome.
\DIFdelbegin \DIFdel{The nominal simulation misses production by far more than a floor: its median offset reaches 109 floors, its }\DIFdelend \DIFaddbegin 

\DIFadd{The simulation dataset does not contain the produced process (Table~\ref{tab:validation}): at the nominal thickness the measured centre lies within the range spanned by the friction grid in only 11 of the 42 combinations of characteristic, geometry and force. The simulated mid-side draw-in exceeds the measured one at 100\,kN for every friction coefficient (by 1.8--2.0\,mm at friction 0.10), reaches it at 300\,kN only at the edge of the grid on the concave cup, and at 500\,kN at friction 0.105 and 0.116; the corner draw-in exceeds it everywhere, by 2.6--3.5\,mm. The simulated }\DIFaddend force effect on the draw-in is 2.7 to 5.1 times the measured one, and the cut arms of the concave cups bend downwards in 89--100\,\% of the simulations but upwards in every part\DIFdelbegin \DIFdel{. Part of the offset is definitional: the cup depth differs by a median 0.90 and 0.95}\DIFdelend \DIFaddbegin \DIFadd{, whereas the convex arm angles agree within 0.6$^\circ$. Calibrating friction cannot repair this within the grid: all 72 within-family least-squares calibrations of Mc end at its upper edge (0.15), and with a discrepancy term the posterior of Mcd piles at the lower edge (0.05), where the simulated force effect is weakest. The dataset's mapping from oil film to friction implies 0.8--3.1\,mm of mid-side draw-in per g/m$^2$ of oil, against within-run slopes of $-0.01$ to $+0.09$\,mm and median slopes between the lubrication patterns of $-0.17$ to $+0.39$}\DIFaddend \,mm \DIFdelbegin \DIFdel{, consistent with the scanned surface against the shell mid-surface, and removing this offset alone brings the decisive distance of the nominal simulation from }\DIFdelend \DIFaddbegin \DIFadd{per g/m$^2$; noise in the oil readings attenuates the within-run slope, but the mapping is too sensitive by an order of magnitude. An envelope over the measured incoming variation instead of the whole grid halves its width but makes M5 less decisive with a sibling (24 against 7.2 floors), because through this mapping it varies with the lubrication pattern more than production does. A Gaussian-process surrogate of the simulations decides as the simulation does (109 against }\DIFaddend 108 \DIFdelbegin \DIFdel{to }\DIFdelend \DIFaddbegin \DIFadd{floors).
}

\DIFadd{The cup depth is simulated 0.70--1.06\,mm (2.5--3.7\,\%) deeper than measured, and the wall position at 15\,mm depth, set by punch and die, 2.0--2.2\,mm (2.8--3.5\,\%) smaller in radius. The depth offset is not definitional, since bottom and flange are seen from the same side of the sheet in both, nor due to the flange inclination, which would make the measured depth larger (\esm{Table~S19}); a shorter effective drawing depth, a different tool geometry or a scale error of the scanner could each produce these offsets, which the data cannot separate. Without the depth offset the nominal simulation decides at }\DIFaddend 65 \DIFdelbegin \DIFdel{floors}\DIFdelend \DIFaddbegin \DIFadd{instead of 108 floors, a sensitivity only}\DIFaddend .

\subsection{Decisions by evidence relation}\label{sec:res-decisions}

Table~\ref{tab:distances} and Fig.~\ref{fig:decisions} score the rules. The nominal simulation decides at 108 floors\DIFdelbegin \DIFdel{(101--115)}\DIFdelend , erring mostly by rejecting adequate alternatives.

\emph{A new variant with a produced sibling.} Every calibrated rule improves on it, and the \DIFdelbegin \DIFdel{simplest }\DIFdelend \DIFaddbegin \DIFadd{production record }\DIFaddend is the most decisive: the nearest \DIFdelbegin \DIFdel{produced }\DIFdelend setting, the mean \DIFdelbegin \DIFdel{$q_{95}$ }\DIFdelend \DIFaddbegin \DIFadd{$q_p$ }\DIFaddend of the two siblings, is right beyond 3.4 floors (2.0--4.5) \DIFdelbegin \DIFdel{, 0.17\,mm of mid-side draw-in; the Gaussian processes decide at 7.2 }\DIFdelend \DIFaddbegin \DIFadd{and the force trend beyond 3.5, while M5 and M7n decide at 6.8--7.2 }\DIFaddend floors but are safe below one floor. \DIFaddbegin \DIFadd{Stratified by what a team knows when deciding, the nearest setting needs 1.60 floors when the two produced siblings do not differ resolvably from each other (75 cases on 17 alternatives) and 4.40 when they do (51 cases on 16), the force trend 3.05 and 4.20, and M5 6.15 and 8.90 (\esm{Table~S3}, with the strata of the first version, defined by the held-out alternative's own production, as a diagnostic). In physical units the nearest setting is right beyond 0.10\,mm of mid-side draw-in, 0.26$^\circ$ of wall angle and 0.04\,mm of cup depth (Table~\ref{tab:groups}). One produced sibling suffices: with a single alternative at the new design's force the nearest setting decides at 3.5 floors (2.3--4.6), and further alternatives do not improve it, while they make M5 more decisive and safer (\esm{Fig.~S3}).
}\DIFaddend 

\emph{A new process setting.} With every alternative at the new design's force held out, \DIFdelbegin \DIFdel{all rules lose: }\DIFdelend the force trend and the nearest \DIFdelbegin \DIFdel{produced }\DIFdelend setting decide at 8.2 and 8.5 floors \DIFdelbegin \DIFdel{, }\DIFdelend \DIFaddbegin \DIFadd{and }\DIFaddend the interval rules \DIFdelbegin \DIFdel{(M2, M2n, M3, M3n, M5, M5n) only at 13--31}\DIFdelend \DIFaddbegin \DIFadd{only at 15--19}\DIFaddend , though they are safe at \DIFdelbegin \DIFdel{5.9--16. Interpolating in force }\DIFdelend \DIFaddbegin \DIFadd{5--8. Interpolating }\DIFaddend is much easier than extrapolating (the nearest setting 4.7 against 9.5 floors), \DIFdelbegin \DIFdel{and the }\DIFdelend \DIFaddbegin \DIFadd{but at 300\,kN, the midpoint of a balanced calibration set, the nearest setting, the force trend and the centre of M5n coincide by construction, so interpolation cannot separate rules of the mean type. The }\DIFaddend scaled simulation, extrapolating a fitted trend, falsely accepts failing alternatives as far as \DIFdelbegin \DIFdel{164 }\DIFdelend \DIFaddbegin \DIFadd{165 }\DIFaddend floors from their limit\DIFaddbegin \DIFadd{, on the cup depth, while on the mid-side draw-in and the arm angle it is more decisive than the force trend (4.0 against 5.7 and 5.9 against 9.5 floors when extrapolating)}\DIFaddend .

\emph{A new \DIFdelbegin \DIFdel{design }\DIFdelend family.} Scored in the floor of the produced family, as the procedure prescribes, every rule without the simulation fails \DIFdelbegin \DIFdel{in the two transfers }\DIFdelend (130--135 floors): nothing in one family's \DIFdelbegin \DIFdel{production }\DIFdelend record describes the other. With the simulation, \DIFdelbegin \DIFdel{the bias correction decides at 33 }\DIFdelend \DIFaddbegin \DIFadd{Mcd decides at 30 }\DIFaddend floors and M5 at 35, and M3 is safe at 17 floors \DIFdelbegin \DIFdel{, M2 at 26 }\DIFdelend and \DIFdelbegin \DIFdel{M5 at 31; in the floor of the new family, which a team would not have, }\DIFdelend M2 \DIFdelbegin \DIFdel{is safe at 16, M3 }\DIFdelend at \DIFdelbegin \DIFdel{21 and }\DIFdelend \DIFaddbegin \DIFadd{26. Per characteristic, }\DIFaddend M5 \DIFdelbegin \DIFdel{at 28. The best rule was thus safe at about 16--17 floorsin either floor, some 0.8\,mm of mid-side draw-in, but two }\DIFdelend \DIFaddbegin \DIFadd{decides at 5--10 floors for the draw-ins and the wall angle, which the forming model represents directly, but at 20--40 for depth, waviness, arm and dome. Production knowledge still carries over where a drawing nominal transfers: for the wall angle, the design angle plus the produced family's deviation from its own decides at 1.9 floors, against 7.5 for M5. Two }\DIFaddend transfers do not resolve which rule is safest\DIFdelbegin \DIFdel{: the safe distance of M3 minus that of M2 lies between $-16$ and $+7$ floors in the produced family's floor and between $-3$ and $+8$ in the new family's}\DIFdelend .

%% BEGIN GENERATED distances
\begin{table}[t]
\tabsize
\DIFdelbeginFL %DIFDELCMD < \setlength{\tabcolsep}{1.9pt}
%DIFDELCMD < %%%
\DIFdelendFL \DIFaddbeginFL \setlength{\tabcolsep}{3.4pt}
\DIFaddendFL \caption{Decisive and safe distances (floors) by the relation of the new design to the produced evidence}\label{tab:distances}
\DIFdelbeginFL %DIFDELCMD < \begin{tabular}{@{}lrrrrrrrrrr@{}}
%DIFDELCMD < \toprule
%DIFDELCMD <  & \multicolumn{2}{c}{New variant} & \multicolumn{4}{c}{New process setting} & \multicolumn{2}{c}{New family} & \multicolumn{2}{c}{Pooled} \\
%DIFDELCMD < \cmidrule(lr){2-3}\cmidrule(lr){4-7}\cmidrule(lr){8-9}\cmidrule(lr){10-11}
%DIFDELCMD < Rule & $\delta_\mathrm{d}/\delta_\mathrm{s}$ & FA & $\delta_\mathrm{d}/\delta_\mathrm{s}$ & FA & interp. & extrap. & $\delta_\mathrm{d}/\delta_\mathrm{s}$ & FA & variant & setting \\
%DIFDELCMD < \midrule
%DIFDELCMD < M0 & 108 & 17 & 108 & 17 & 109 & 106 & 104 & 28 & 108 & 108 \\
%DIFDELCMD < M1 & 15 & 15 & 24 & 24 & 11 & 25 & 33 & 35 & 27 & 28 \\
%DIFDELCMD < \quad interval & 12--18 &  & 21--25 &  & 10--11 & 24--25 & 31--36 &  &  & \\
%DIFDELCMD < M1s & 4.9 & 4.7 & 19 & 164 & 19 & 19 & 121 & 125 & 15 & 20 \\
%DIFDELCMD < \quad interval & 4.7--6.3 &  & 14--21 &  & 19--20 & 15--19 & 119--121 &  &  & \\
%DIFDELCMD < M2 & 23/0.15 & 0.30 & 18/8.1 & 8.3 & 18/2.7 & 18/9.9 & 41/26 & 27 & 53/0 & 53/1.0 \\
%DIFDELCMD < \quad interval & 22--23 &  & 17--18 &  & 15--18 & 17--19 & 39--42/9.8--32 &  &  & \\
%DIFDELCMD < M5 & 7.2/0.75 & 0.15 & 19/6.8 & 7.6 & 19/2.8 & 19/8.9 & 35/31 & 33 & 11/0.05 & 28/5.4 \\
%DIFDELCMD < \quad interval & 6.6--8.4 &  & 18--20 &  & 18--20 & 18--20 & 31--39/28--32 &  &  & \\
%DIFDELCMD < M5n & 7.2/0.55 & 0.10 & 17/6.1 & 9.0 & 16/0 & 18/8.6 & 133/128 & 127 & 9.8/0.20 & 16/2.7 \\
%DIFDELCMD < \quad interval & 6.4--7.5 &  & 16--18 &  & 15--16 & 16--19 & 131--133/127--128 &  &  & \\
%DIFDELCMD < NN & 3.4 & 2.8 & 8.5 & 10 & 4.7 & 9.5 & 130 & 129 & 3.4 & 8.5 \\
%DIFDELCMD < \quad interval & 2.0--4.5 &  & 5.5--10 &  & 4.3--5.1 & 6.6--11 & 129--131 &  &  & \\
%DIFDELCMD < \midrule M0w & 119/94 & 9.2 & 119/94 & 9.2 & 119/94 & 119/77 & 117/79 & 10 & 119/94 & 119/94 \\
%DIFDELCMD < Mc & 27 & 28 & 39 & 39 & 14 & 40 & 43 & 53 & 32 & 33 \\
%DIFDELCMD < M3 & 23/0 & 0 & 31/16 & 13 & 29/13 & 33/17 & 40/17 & 18 & 25/0 & 32/5.1 \\
%DIFDELCMD < \quad interval & 17--26 &  & 25--36 &  & 24--31 & 25--40 & 36--45/15--19 &  &  & \\
%DIFDELCMD < M3n & 11/0.05 & 0 & 13/7.4 & 9.2 & 15/7.0 & 13/7.8 & 135/126 & 124 & 15/0 & 15/4.9 \\
%DIFDELCMD < M1n & 9.1 & 10 & 10 & 12 & 4.8 & 12 & 131 & 129 & 130 & 130 \\
%DIFDELCMD < M1sn & 3.5 & 3.5 & 8.2 & 9.0 & 4.7 & 8.9 & 131 & 129 & 77 & 111 \\
%DIFDELCMD < \quad interval & 3.1--4.2 &  & 6.7--9.0 &  & 4.3--5.1 & 8.6--9.5 & 129--131 &  &  & \\
%DIFDELCMD < M2n & 19/0 & 0.20 & 15/5.9 & 9.2 & 13/0 & 16/9.1 & 133/128 & 126 & 325/0 & 325/4.0 \\
%DIFDELCMD < M6 & 17/0.05 & 0 & 56/6.2 & 5.8 & 47/0 & 56/8.1 & 35/28 & 28 & -- & -- \\
%DIFDELCMD < \botrule
%DIFDELCMD < \end{tabular}
%DIFDELCMD < %%%
\DIFdelendFL \DIFaddbeginFL \begin{tabular}{@{}lrrrrrrrr@{}}
\toprule
 & \multicolumn{2}{c}{Sibling} & \multicolumn{2}{c}{Interpolated} & \multicolumn{2}{c}{Extrapolated} & \multicolumn{2}{c}{New family} \\
\cmidrule(lr){2-3}\cmidrule(lr){4-5}\cmidrule(lr){6-7}\cmidrule(lr){8-9}
Rule & $\delta_\mathrm{d}/\delta_\mathrm{s}$ & FA/FR & $\delta_\mathrm{d}/\delta_\mathrm{s}$ & FA/FR & $\delta_\mathrm{d}/\delta_\mathrm{s}$ & FA/FR & $\delta_\mathrm{d}/\delta_\mathrm{s}$ & FA/FR \\
\midrule
M0 & 108 & 17/109 & 109 & 17/115 & 106 & 12/109 & 104 & 28/105 \\
M1s$^\dagger$ & 4.9 & 4.7/4.9 & 19 & 30/19 & 19 & 165/7.9 & 121 & 125/118 \\
M2 & 23/$<$1 & $<$1/$<$1 & 18/2.7 & 3.7/2.4 & 18/9.9 & 8.7/10.0 & 41/26 & 27/9.7 \\
M5 & 7.2/$<$1 & $<$1/$<$1 & 19/2.8 & 2.4/3.1 & 19/8.9 & 8.4/9.7 & 35/31 & 33/15 \\
\quad interval of $\delta_\mathrm{d}$ & 6.6--8.4 &  & 18--20 &  & 18--20 &  & 31--39 & \\
M7n$^\ddagger$ & 6.8/$<$1 & $<$1/$<$1 & 7.4/3.4 & 2.2/3.4 & 15/5.9 & 7.5/$<$1 & 132/130 & 129/167 \\
\quad interval of $\delta_\mathrm{d}$ & 6.0--7.3 &  & 7.1--8.1 &  & 15--16 &  & 131--132 & \\
NN$^\dagger$ & 3.4 & 2.8/3.4 & 4.7 & 3.1/4.7 & 9.5 & 11/5.6 & 130 & 129/168 \\
\quad interval of $\delta_\mathrm{d}$ & 2.0--4.5 &  & 4.3--5.1 &  & 6.6--11 &  & 129--131 & \\
M1sn$^\dagger$ & 3.5 & 3.5/3.6 & 4.7 & 3.1/4.7 & 8.9 & 9.3/7.4 & 131 & 129/168 \\
NN+RM$^\ddagger$ & 6.3/$<$1 & $<$1/$<$1 & 9.6/$<$1 & $<$1/$<$1 & 20/$<$1 & 1.5/$<$1 & 293/4.5 & 4.1/44 \\
\midrule Trial, first 20 parts & 1.4 & 1.5/1.2 & 1.4 & 1.5/1.2 & 1.4 & 1.5/1.2 & 1.4 & 1.5/1.2 \\
Trial, first 50 parts & 0.90 & 1.0/0.85 & 0.90 & 1.0/0.85 & 0.90 & 1.0/0.85 & 0.90 & 1.0/0.85 \\
\botrule
\end{tabular}
\DIFaddendFL \footnotetext{Over all seven characteristics; one number for rules that always decide\DIFaddbeginFL \DIFaddFL{; $<$1: below one floor}\DIFaddendFL . FA\DIFaddbeginFL \DIFaddFL{/FR}\DIFaddendFL : safe \DIFdelbeginFL \DIFdelFL{distance }\DIFdelendFL \DIFaddbeginFL \DIFaddFL{distances }\DIFaddendFL on the false-accept \DIFaddbeginFL \DIFaddFL{and the false-reject }\DIFaddendFL side\DIFdelbeginFL \DIFdelFL{(failing alternatives)}\DIFdelendFL . \DIFdelbeginFL \DIFdelFL{Interval: bootstrap 95\,\% interval of the decisive distance }\DIFdelendFL \DIFaddbeginFL \DIFaddFL{126 cases with a sibling }\DIFaddendFL and \DIFdelbeginFL \DIFdelFL{, }\DIFdelendFL for a new family, \DIFdelbeginFL \DIFdelFL{of the safe distance (fits fixed). 126 cases per relation (interp.: }\DIFdelendFL 42 interpolated\DIFdelbeginFL \DIFdelFL{; extrap.: }\DIFdelendFL \DIFaddbeginFL \DIFaddFL{, }\DIFaddendFL 84 extrapolated\DIFdelbeginFL \DIFdelFL{)}\DIFdelendFL ; a new family \DIFdelbeginFL \DIFdelFL{rests on two transfers and is scored }\DIFdelendFL in the \DIFdelbeginFL \DIFdelFL{floor of the }\DIFdelendFL produced family\DIFaddbeginFL \DIFaddFL{'s floor }\DIFaddendFL (\DIFdelbeginFL \DIFdelFL{Section~\ref{sec:floor}}\DIFdelendFL \DIFaddbeginFL \DIFaddFL{in its own floor in \esm{Table~S1}}\DIFaddendFL ). \DIFdelbeginFL \DIFdelFL{Pooled}\DIFdelendFL \DIFaddbeginFL \DIFaddFL{Interval}\DIFaddendFL : \DIFaddbeginFL \DIFaddFL{95\,\% over 1000 resamples of }\DIFaddendFL the \DIFdelbeginFL \DIFdelFL{other geometry's }\DIFdelendFL \DIFaddbeginFL \DIFaddFL{held-out }\DIFaddendFL alternatives\DIFdelbeginFL \DIFdelFL{added to the calibration set. Below the line}\DIFdelendFL , \DIFdelbeginFL \DIFdelFL{the further rules (M1sn}\DIFdelendFL \DIFaddbeginFL \DIFaddFL{fits fixed. Trial}\DIFaddendFL : the \DIFdelbeginFL \DIFdelFL{force trend, the counterpart }\DIFdelendFL \DIFaddbeginFL \DIFaddFL{first $n$ parts }\DIFaddendFL of \DIFdelbeginFL \DIFdelFL{M1s without }\DIFdelendFL the \DIFdelbeginFL \DIFdelFL{simulation) and}\DIFdelendFL \DIFaddbeginFL \DIFaddFL{design's own run}\DIFaddendFL , \DIFdelbeginFL \DIFdelFL{as }\DIFdelendFL a \DIFdelbeginFL \DIFdelFL{robustness variant}\DIFdelendFL \DIFaddbeginFL \DIFaddFL{within-run lower bound. $^\dagger$}\DIFaddendFL , \DIFdelbeginFL \DIFdelFL{M6 (}\DIFdelendFL \DIFaddbeginFL \DIFaddFL{$^\ddagger$: added after }\DIFaddendFL the \DIFdelbeginFL \DIFdelFL{M5 mean with a conformal margin on standardised leave-one-out residuals)}\DIFdelendFL \DIFaddbeginFL \DIFaddFL{first results and for this version}\DIFaddendFL .}
\end{table}
%% END GENERATED distances

\begin{figure}[t]
\centering
\DIFdelbeginFL %DIFDELCMD < \includegraphics[width=\textwidth]{F3_decisions_compare}
%DIFDELCMD < \caption{Verdict shares against the requirement distance $d$ (symmetric logarithmic axis, linear within $\pm$10 floors; 126 cases per panel; a new family in the produced family's floor), with $\delta_\mathrm{d}$ and $\delta_\mathrm{s}$}%%%
\DIFdelendFL \DIFaddbeginFL \includegraphics[width=\textwidth]{F3_decisions_main}
\caption{Verdict shares against the requirement distance $d$ (symmetric logarithmic axis, linear within $\pm10$ floors; 126 cases per panel; a new family in the produced family's floor)}\DIFaddendFL \label{fig:decisions}
\end{figure}

\DIFdelbegin %DIFDELCMD < \subsection{Where the distances come from}\label{sec:res-drivers}
%DIFDELCMD < %%%
\DIFdelend \DIFaddbegin \subsection{Learning given a fair test}\label{sec:res-ml}
\DIFaddend 

\DIFdelbegin \DIFdel{Three features of the design shape these distances. First, the sibling relation mixes replication with variation: in 58 of the 126 casesneither sibling differs resolvably from the held-out alternative, and the nearest settingis then right beyond 0.85 floors, the reproducibility of $q_{95}$ between near-replicate runs; it needs 2.5 floors when one sibling differs resolvably (34 cases) and 4.9 when both do (34 cases), where }\DIFdelend \DIFaddbegin \DIFadd{The Gaussian processes fitted by maximum marginal likelihood reduce to the baselines in this design (Table~\ref{tab:gp}). For a new setting the force length scale of M5n sits at its lower bound in 68--79\,\% of the fits and its noise variance at its bound in every fit, and its centre lies within 0.01 floor of the calibration mean in 93--95\,\% of the cases; with a sibling it lies a median 0.28 floor from the nearest setting. The trend GP, specified for this regime, reduces to }\DIFaddend the force trend \DIFdelbegin \DIFdel{needs 4.5 and M5 9.7. Second, the extrapolation is confounded with stroke speed: the nearest settingdecides at 11 floors at 100\,kN, where the press also ran slower, against 4.9 at 500}\DIFdelend \DIFaddbegin \DIFadd{instead: its centre lies a median 0.33--0.36 floor from it for a new setting, its force length scale stays near the prior's 200}\DIFaddend \,kN\DIFdelbegin \DIFdel{(the force trend at 8.9 at both). Taken together, the nearest setting decided at up to about 5 floors whenever the new design differed resolvably and force and speed did not change together, though each stratum has 34 to 58 cases and each distance rests on a handful of verdicts. Third, the other family's record does not help with a produced sibling: pooling it makes most rules less decisive (M5 from 7.2 to 11 floors), though it raises M5's coverage from 86 to 90\,\%.
}%DIFDELCMD < 

%DIFDELCMD < %%%
\DIFdel{The relation therefore does not govern decision fitness in general. Over the eleven calibrated rules and three relations, the relation explains 67\,\% of the variation of the log decisive distance and the rule 6\,\% (69 }\DIFdelend \DIFaddbegin \DIFadd{, which two force levels cannot inform, }\DIFaddend and \DIFdelbegin \DIFdel{6\,\% in the new family's own floor), but only because rules without the simulation fail }\DIFdelend \DIFaddbegin \DIFadd{its between-run nugget, 0.4--0.5 floors, is three to four times the sampling error of the truths. It is more decisive than M5n }\DIFaddend for a new \DIFdelbegin \DIFdel{family by construction; over the two relations within a family, the rule explains 68\,\% and the relation 15\,\% (71 and 15\,\% with M1sn). Pooled over strata, the best attainable distance grows with novelty: 3.4 }\DIFdelend \DIFaddbegin \DIFadd{setting (14.5 against 17.2 floors; paired difference $-2.7$, $-5.6$ to $-1.1$) and as safe, but less decisive than the nearest setting in every relation (by 3.5 }\DIFaddend floors with a \DIFdelbegin \DIFdel{produced sibling , 8.2 }\DIFdelend \DIFaddbegin \DIFadd{sibling and 6 }\DIFaddend for a new setting\DIFdelbegin \DIFdel{, 33 for a new family. }%DIFDELCMD < 

%DIFDELCMD < \subsection{What the simulation and learning contribute}\label{sec:res-ml}
%DIFDELCMD < 

%DIFDELCMD < %%%
\DIFdel{Pairing each calibrated rule with its counterpart without the simulation isolates what the simulation adds as used. Added as an offset, it adds nothing with a produced sibling and degrades two rules: M5 decides where M5n does (difference 0.0 floors, $-0.4$ to $+1.1$), while M2 and M1 decide 3.9 ($+0.5$ to $+7.4$) and 6.2 floors ($+4.6$ to $+9.7$) further out than M2n and M1n, and M1 13.5 further out }\DIFdelend \DIFaddbegin \DIFadd{). With the simulation in its mean (M7) it gets worse }\DIFaddend for a new setting \DIFaddbegin \DIFadd{(28 floors)}\DIFaddend , because the simulated force \DIFdelbegin \DIFdel{effect is too large. Rescaled, itdoes no better: within a geometry the nominal simulation is a function of the force, and the straight line in force (M1sn) decides 1.4 floors closer than the scaled simulation }\DIFdelend \DIFaddbegin \DIFadd{trend enters it. The archive's study of part-level boosting (regularised, flexible, ridge, quantile loss, inner selection) finds no setting closer than 11.2 floors }\DIFaddend with a sibling \DIFdelbegin \DIFdel{($+0.65$ to $+3.0$) and 11 closer }\DIFdelend \DIFaddbegin \DIFadd{and 13.1 }\DIFaddend for a new setting\DIFdelbegin \DIFdel{($+5.3$ to $+13.8$), in every resample and refit configuration; the 4.2 floors that the scaled simulation gains over the median of the produced alternatives with a sibling come from the force level. Only for a new family is the simulation decisive, by 92 to 98 floors as an offset and 10 rescaled, and even there production knowledge carries over where a drawing nominal transfers: for the wall angle, the design angle plus the produced family's deviation from its own decides at 1.9 floors, against 7.5 }\DIFdelend \DIFaddbegin \DIFadd{.
}

\DIFadd{Prediction and interval separate (Table~\ref{tab:crossing}). A relation-matched interval borrowed from the other family makes even the nearest setting safe for a new setting (safe distance $<$1 floor, against 6.8 }\DIFaddend for M5 \DIFdelbegin \DIFdel{(for the arm angle and the dome, whose nominal is zero in both families, this baseline is M1n, at 44 and 27 floors). A Gaussian-process surrogate of the simulations over force, friction and thickness, with a leave-one-out root-mean-square error below 2 floors for 12 of the 14 characteristics and geometries, decides as the simulation does (109 }\DIFdelend \DIFaddbegin \DIFadd{and 5.2 for M7n) at the decisive distance of the learned intervals (19.5 }\DIFaddend against \DIFdelbegin \DIFdel{108 floors as M0, 15 against 15 as M1 with a sibling).
}%DIFDELCMD < 

%DIFDELCMD < %%%
\DIFdel{Learning adds intervals rather than accuracy. Used as point rules, the centres of the M5}\DIFdelend \DIFaddbegin \DIFadd{18.7 and 14.5 floors)}\DIFaddend , \DIFdelbegin \DIFdel{M5n }\DIFdelend and \DIFdelbegin \DIFdel{M3n intervals decide at 3.3, 3.1 }\DIFdelend \DIFaddbegin \DIFadd{covers 88\,\% of the truths interpolating and 93\,\% extrapolating, against 71 }\DIFaddend and \DIFdelbegin \DIFdel{3.7 floors with a produced sibling, as close as the nearest setting (Table~\ref{tab:intervals}}\DIFdelend \DIFaddbegin \DIFadd{46\,\% for M5 and 52 and 56\,\% for M7n. By the interval score \mbox{%DIFAUXCMD
\citep{gneiting2007strictly} }\hskip0pt%DIFAUXCMD
the relation-matched intervals are the best for a new setting (12--14 interpolating and 19--24 extrapolating, against 14--98 and 41--147 for the rules' own intervals; \esm{Table~S4}}\DIFaddend ). With a sibling the \DIFdelbegin \DIFdel{Gaussian-process intervals cover 85--86\,\% of the true $q_{95}$ against a nominal 90\,\%, but their noise variance sits at its lower bound, the variance between the calibration series of one setting, in 97--99}\DIFdelend \DIFaddbegin \DIFadd{model-based intervals (M5, M5n, M7n, M7, Mcd) cover 83--86}\DIFaddend \,\% \DIFdelbegin \DIFdel{of these fits and in every new-setting and new-family fit, so this coverage comes from the quasi-replicate series rather than from the learner; the largest-residual margin covers 89}\DIFdelend \DIFaddbegin \DIFadd{(cluster intervals 0.71--0.94), compatible with their nominal 90}\DIFaddend \,\%\DIFaddbegin \DIFadd{, }\DIFaddend and the jackknife+ intervals \DIFdelbegin \DIFdel{96\,\% (}\DIFdelend \DIFaddbegin \DIFadd{of }\DIFaddend M3 \DIFdelbegin \DIFdel{) and 89}\DIFdelend \DIFaddbegin \DIFadd{96}\DIFaddend \,\% (\DIFdelbegin \DIFdel{M3n) , above their guarantee of 7/9. For a new setting the calibrated intervals cover 43--55}\DIFdelend \DIFaddbegin \DIFadd{0.93--0.98) against a guarantee of 0.78; the conformal margin of M6, computed on sibling-type residuals, covers 93}\DIFaddend \,\% \DIFdelbegin \DIFdel{(M2n and M5n 93}\DIFdelend \DIFaddbegin \DIFadd{with a sibling at a cost of ten floors of decisiveness against M5 (17 against 7.2) but 61}\DIFaddend \,\% \DIFdelbegin \DIFdel{interpolating), for }\DIFdelend \DIFaddbegin \DIFadd{extrapolating. For }\DIFaddend a new family \DIFdelbegin \DIFdel{the Gaussian processes at most 12\,\%: trees do not extrapolate, two force levels identify at most a linear trend, and in the median new-setting fit the length scale of the force sits at its lower bound.
}\DIFdelend \DIFaddbegin \DIFadd{no interval covers well except at widths that send every design to trial. With one produced alternative of the new family, its offset to the other family's record decides at 8.9 floors in the new family's floor and the family's own record at 10.5; with two or three, its own nearest setting decides at 7.7 and 5.9 floors, while a two-task Gaussian process with a free inter-family correlation is safe at 1.3--1.6 floors but decides only at 13--17.
}\DIFaddend 

\DIFdelbegin %DIFDELCMD < \subsection{Calibration budget}\label{sec:res-budget}
%DIFDELCMD < %%%
\DIFdelend %DIF > % BEGIN GENERATED guards
\DIFaddbegin \begin{table}[t]
\tabsize
\setlength{\tabcolsep}{2.2pt}
\caption{Guard bands of step 6 (floors), their out-of-sample share correct and the trial rate}\label{tab:guards}
\begin{tabular*}{\textwidth}{@{\extracolsep{\fill}}lrrrrrrrrrrrr@{}}
\toprule
 & \multicolumn{4}{c}{Sibling} & \multicolumn{4}{c}{Interpolated} & \multicolumn{4}{c}{Extrapolated} \\
\cmidrule(lr){2-5}\cmidrule(lr){6-9}\cmidrule(lr){10-13}
Rule & $g$ & $g_5$ & OOS & T & $g$ & $g_5$ & OOS & T & $g$ & $g_5$ & OOS & T \\
\midrule
NN$^\dagger$ & 0.25 & 2.0 & 95 & 2 & 2.0 & 4.25 & 95 & 18 & 7.0 & 14.25 & 94 & 55 \\
M1sn$^\dagger$ & 0.75 & 2.75 & 95 & 7 & 2.0 & 4.25 & 95 & 18 & 6.0 & 8.5 & 95 & 48 \\
M2 & 0 & 0 & 99 & 54 & 0.75 & 2.75 & 95 & 53 & none & none & -- & 100 \\
M5 & 0 & 0 & 99 & 23 & 0.5 & 2.75 & 96 & 54 & none & none & -- & 100 \\
M7n$^\ddagger$ & 0 & 6.75 & 98 & 19 & 0.75 & 3.75 & 95 & 27 & 3.25 & 6.75 & 95 & 56 \\
NN+RM$^\ddagger$ & 0 & 0.5 & 99 & 28 & 0 & 0 & 99 & 42 & 0 & 5.25 & 97 & 72 \\
\botrule
\end{tabular*}
\footnotetext{\DIFaddFL{$g$, $g_5$: guard band under the reference prior and under the release-like prior (requirements within $\pm5$ floors); none: no band up to 20 floors reaches 95\,\%. OOS: share correct (\%) among decided verdicts with the band estimated leaving each alternative out within its family. T: share of verdicts sent to trial (\%), reference prior. No rule qualifies for a new family; bands per family, out of sample and under other priors in \esm{Tables~S7--S10}.}}
\end{table}
%DIF > % END GENERATED guards
\DIFaddend 

\DIFdelbegin \DIFdel{A single produced sibling decides (Fig.~\ref{fig:budget}): with one alternative at the new design's force, the bias correction, the envelope rule and the nearest setting coincide and decide at 3.5 floors(2.3--4.6) without a wrong verdict at 10 floors . Further alternatives do not improve the nearest setting; they make the bias correction worse (12--18 floors), the envelope rule safer but less decisive, }\DIFdelend \DIFaddbegin \subsection{Guard bands}\label{sec:res-procedure}

\DIFadd{Step 5 turns the margins into guard bands (Table~\ref{tab:guards}; \esm{Fig.~S4}). Under the reference prior the nearest setting needs 0.25 floors with a sibling, 2.0 interpolating and 7.0 extrapolating, and sends 2, 18 and 55\,\% of the verdicts to trial; M2 }\DIFaddend and \DIFdelbegin \DIFdel{M5, shown from three alternatives because a Gaussian process fits its hyperparameters to them, safer and more decisive (13 to 7.2 floors). Without a sibling , bracketing the new force buys safety but not decisiveness: the envelope rule and }\DIFdelend M5 \DIFdelbegin \DIFdel{are safe at 2.7--4.8 floors but decide only at 15--25, against 4.7--6.0 for the nearest setting, and extrapolation stays unsafe at every budget. With calibration series of 50, 100 or 250 parts instead of 500, no distance within a family moves by more than }\DIFdelend \DIFaddbegin \DIFadd{need no band with a sibling and less than one floor interpolating, but send 23--54\,\% of the verdicts to trial, and they do not qualify for an extrapolated setting. Estimated per family, as step 5 prescribes, the bands differ between the geometries (the nearest setting }\DIFaddend 1.5 \DIFdelbegin \DIFdel{floors, }\DIFdelend and \DIFdelbegin \DIFdel{for a new family by at most 4 in its own floor (Table~\ref{tab:robust}): rules can be calibrated on short runs, but the floors and the held-out truths came from the full series, and whether truths from short runs suffice to qualify the rules was not tested.
}%DIFDELCMD < 

%DIFDELCMD < \begin{figure}[t]
%DIFDELCMD < \centering
%DIFDELCMD < \includegraphics[width=\textwidth]{F4_budget}
%DIFDELCMD < \caption{Decisive and safe distances against the number of produced alternatives of the family, by the relation of the new design to them; bands: 95\,\% intervals over the held-out alternatives}\label{fig:budget}
%DIFDELCMD < \end{figure}
%DIFDELCMD < 

%DIFDELCMD < \subsection{The procedure's decision rule}\label{sec:res-procedure}
%DIFDELCMD < 

%DIFDELCMD < %%%
\DIFdel{What a designer observes for a single verdict is the margin between the predicted interval and the requirement, }\DIFdelend \DIFaddbegin \DIFadd{2.5 floors interpolating, 9.5 and 5.75 extrapolating), }\DIFaddend and \DIFdelbegin \DIFdel{its value as evidence depends on the relation (Fig.~\ref{fig:reliability}): with a produced sibling at least 95}\DIFdelend \DIFaddbegin \DIFadd{M2 and M5 qualify for an extrapolated setting on the convex cup only; moved to the concave cup, their bands give only 84 and 83}\DIFaddend \,\% \DIFdelbegin \DIFdel{of the decided verdictsof }\DIFdelend \DIFaddbegin \DIFadd{correct decided verdicts. Out of sample the bands hold: leaving one alternative out keeps the share of correct decided verdicts at 0.94--0.99 for every rule and relation, and leaving one force cell out at 0.93--0.95 for }\DIFaddend the \DIFdelbegin \DIFdel{interval rules are correct at every margin, for a new setting the share grows with the margin, and for a new family }\DIFdelend \DIFaddbegin \DIFadd{nearest setting and the force trend, but at 0.83 and 0.81 for }\DIFaddend M2 and M5 \DIFdelbegin \DIFdel{are right in at most 79}\DIFdelend \DIFaddbegin \DIFadd{when extrapolating. The 95}\DIFaddend \,\% \DIFdelbegin \DIFdel{of their decided verdicts.
}%DIFDELCMD < 

%DIFDELCMD < %%%
\DIFdel{Step 5 turns this into guard bands (Table~\ref{tab:step6}). Under the reference prior the interval rules need none }\DIFdelend \DIFaddbegin \DIFadd{is an average over characteristics: }\DIFaddend with a sibling \DIFdelbegin \DIFdel{and }\DIFdelend the nearest setting \DIFdelbegin \DIFdel{0.25 floors; for an interpolated setting the nearest setting needs 2.0 floors and then decides at 6.7 floors without trials at 10 floors, whereas M2 and M5 need 0.75 and 0.5 floors but send 31 and 40}\DIFdelend \DIFaddbegin \DIFadd{is right in 89\,\% of its decided verdicts on the dome, 98.5}\DIFaddend \,\% \DIFdelbegin \DIFdel{to trial; for an extrapolated setting it needs 7.0 floors and sends a quarter to trial. M0 and M0w never qualify, nor do M1 and Mc with a sibling, M1s for an interpolated setting, M1, M1s, M2, M5 and Mc for an extrapolated one, or any rule for a new family}\DIFdelend \DIFaddbegin \DIFadd{on the waviness and 92\,\% on the cut part}\DIFaddend .

\DIFdelbegin \DIFdel{These }\DIFdelend \DIFaddbegin \DIFadd{The }\DIFaddend bands belong to the prior as much as to the rules. \DIFdelbegin \DIFdel{A prior concentrated near the truth, where release-level requirements lie (the grid within $\pm5$ floors), raises }\DIFdelend \DIFaddbegin \DIFadd{Under the release-like prior }\DIFaddend the band of the nearest setting \DIFaddbegin \DIFadd{grows }\DIFaddend to 2.0 floors with a sibling, 4.25 interpolating and 14.25 extrapolating\DIFdelbegin \DIFdel{(8.0}\DIFdelend , \DIFdelbegin \DIFdel{5.0 and 14.75 within $\pm3$); a uniform prior within $\pm20$ floors lowers them to 0, 0.25 and 3.5 floors and qualifies M2 and M5 for an extrapolated setting (6.0 and 5.0 floors), and one within $\pm50$ floors qualifies M2, M3 and M5 for a new family (6.25, 6.5 }\DIFdelend and \DIFdelbegin \DIFdel{15 floors). That every new-family design goes to trial thus follows from }\DIFdelend \DIFaddbegin \DIFadd{it decides only 63, 32 and 2\,\% of the verdicts; leaving a force cell out, its share correct drops to 0.72 when extrapolating. Under }\DIFaddend the reference prior \DIFdelbegin \DIFdel{and the 20-floor cap}\DIFdelend \DIFaddbegin \DIFadd{no rule qualifies for a new family}\DIFaddend , as the practice of trying out every new die would have it\DIFdelbegin \DIFdel{. The bands are chosen on the evaluation data, which makes the trial rates optimistic. The }\DIFdelend \DIFaddbegin \DIFadd{: Mcd reaches 95\,\% on one transfer with a band of 18 floors, but that band gives 73\,\% on the other (under a uniform prior within $\pm50$ floors M2, M3 and M5 qualify). The applicability }\DIFaddend guards, reported as diagnostics, refuse indiscriminately: the range guard refuses exactly the extrapolated designs\DIFdelbegin \DIFdel{(84 of the 126 new-setting cases), and }\DIFdelend \DIFaddbegin \DIFadd{, and 96\,\% of its refusals }\DIFaddend at requirements 10 floors from the truth \DIFdelbegin \DIFdel{96\,\% of its refusals }\DIFdelend would have been right for the nearest setting\DIFdelbegin \DIFdel{and 72\,\% for M5 (novelty guard 96 and 73\,\%)}\DIFdelend .

\DIFdelbegin \emph{\DIFdel{A worked example.}} %DIFAUXCMD
\DIFdel{A tryout engineer has to decide whether the concave cup at 300\,kN with medium lubrication, a new setting, can be released without a trial run, with the six alternatives at 100 and 500\,kN produced and }\DIFdelend \DIFaddbegin \subsection{Costs, trials and release-level requirements}\label{sec:res-costs}

\DIFadd{Pricing the verdicts by Eq.~(\ref{eq:cost}) with }\DIFaddend a \DIFdelbegin \DIFdel{requirement, chosen for illustration, that 95\,\% of the parts have a mid-side draw-in, by which tryout sets the blank-holder force , of at most 15.30\,mm. The floor is 0.050\,mm, the relation is interpolation and no guard fires. Alone, the nominal simulation predicts 15.66\,mm and fails the design . The nearest produced setting , the mean of }\DIFdelend \DIFaddbegin \DIFadd{false reject at half the cost of a false accept (Fig.~\ref{fig:costs}), }\DIFaddend the \DIFdelbegin \DIFdel{six alternatives at the neighbouring forces, predicts 14.99}\DIFdelend \DIFaddbegin \DIFadd{production record and the Gaussian processes cost about the same with a sibling while trials are cheap, and the nearest setting is the cheapest rule once a trial costs a tenth of a false accept or more; for a new setting the force trend or the nearest setting is the cheapest at every trial cost; and for a new family trying every design is the cheapest until a trial costs about a tenth of a false accept, beyond which M2 and then M5 are. The cost-minimising band of the nearest setting falls from 5 to 0 floors with a sibling and from 15 to 0 when extrapolating as the trial cost rises from 0.01 to 0.5 of a false accept; the 95}\DIFaddend \,\DIFdelbegin \DIFdel{mm, a margin of 6.2 floors}\DIFdelend \DIFaddbegin \DIFadd{\% bands of step 5 correspond to trials costing a tenth to a quarter of a false accept}\DIFaddend , more than \DIFdelbegin \DIFdel{its guard band of 2.0 floors (5.0 under the grid within $\pm3$ floors), so step 6 returns meets. Widened by their guard bands , M2, M5 }\DIFdelend \DIFaddbegin \DIFadd{a new force on an existing tool costs, so cost-optimal bands would be wider. A trial is not error-free either: scored as a rule, a trial of the first 20 or 50 parts of the design's own run decides at 1.35 }\DIFaddend and \DIFdelbegin \DIFdel{M5n return 14.41--15.41}\DIFdelend \DIFaddbegin \DIFadd{0.90 floors}\DIFaddend , \DIFdelbegin \DIFdel{14.23--15.61 and 14.66--15.33\,mm and send the design to trial . The 500 parts had a 95th percentile of 14.88\,mm, 8.4 floors below the limit. The case is typical, since at 8.5 floorsabove the truth step 6 with the nearest setting is right in all 42 interpolated cases, but its guard band was estimated on data that include it}\DIFdelend \DIFaddbegin \DIFadd{a lower bound, since a trial on another day also carries the variation between runs}\DIFaddend .

\DIFdelbegin %DIFDELCMD < \begin{figure}[t]
%DIFDELCMD < \centering
%DIFDELCMD < \includegraphics[width=\textwidth]{F5_reliability}
%DIFDELCMD < \caption{Share correct among decided verdicts with at least the given margin; reference prior: 47 grid distances within $\pm$20 floors, equal weights; new family in the produced family's floor; filled markers: $g$}\label{fig:reliability}
%DIFDELCMD < \end{figure}
%DIFDELCMD < 

%DIFDELCMD < \subsection{Requirements anchored in capability and in tolerances}\label{sec:res-requirements}
%DIFDELCMD < 

%DIFDELCMD < %%%
\DIFdelend Upper limits at which an alternative has a performance index of 1.0, 1.33, 1.67 or 2.0 lie a median \DIFdelbegin \DIFdel{1.0, 1.7, 2.5 and 3.1 }\DIFdelend \DIFaddbegin \DIFadd{1.03, 1.75, 2.46 and 3.14 }\DIFaddend floors above its 95th percentile\DIFdelbegin \DIFdel{; up to }\DIFdelend \DIFaddbegin \DIFadd{. A limit at }\DIFaddend an index of 1.33\DIFdelbegin \DIFdel{they lie inside every decisive distance except a near-replicate sibling's (0.85 floors). The test is one-sided by construction: every alternative meets these limits, so a verdict can only be right, a false reject or a trial, and a rule biased towards meets scores well. At an index of 1.33 the shares of right verdicts, false rejects and trials are, with a produced sibling and for a new setting, 91/9/0 and 65/35/0\,\% for }\DIFdelend \DIFaddbegin \DIFadd{, a recommended minimum for existing processes and the lower end of the range in which an initial process study may be acceptable after a customer review, therefore lies inside every attainable distance except that of }\DIFaddend the nearest setting \DIFdelbegin \DIFdel{, 44/2/53 and 32/16/52\,\% for M5, and 81/19/0 and 79/21/0\,\% for the scaled simulation, which falsely accepts failing alternatives elsewhere; for a new family no rule is right in more than 54\,\%. Release-level requirements are thus }\DIFdelend \DIFaddbegin \DIFadd{when the produced siblings do not differ resolvably from each other (1.60 floors), which it exceeds by only 0.15 floor: release-level requirements are }\DIFaddend decided at the design stage only with a \DIFaddbegin \DIFadd{near-replicate }\DIFaddend produced sibling, \DIFdelbegin \DIFdel{and then by the production record.
General angular tolerances of ISO~2768-1 \mbox{%DIFAUXCMD
\citep{iso1989tolerances} }\hskip0pt%DIFAUXCMD
give a coarser scenario: only 32 of its 108 decisions lie within 5 floorsof the limit, 6 of them on failing alternatives, and a cost comparison inherits this coarseness. With a false reject at half the cost of a false accept, the cheapest rule is (i) with a produced sibling , the nearest setting, the scaled simulation and }\DIFdelend \DIFaddbegin \DIFadd{if at all.
}

\begin{figure}[t]
\centering
\includegraphics[width=\textwidth]{F_costs}
\caption{Expected cost per verdict with the cost-minimising band (a--d) and that band (e--h) against the cost of a trial; reference prior, $c_\mathrm{FR}=0.5\,c_\mathrm{FA}$}\label{fig:costs}
\end{figure}

\subsection{Sensitivity, selection and robustness}\label{sec:res-robust}

\DIFadd{A between-run component added to every truth lengthens the distances but keeps their order (\esm{Table~S12}): with $\tau=3$ floors, about the largest variation between series, the nearest setting needs 5.8 floors with a sibling and 10 extrapolating, M5 9.3 and 20, a production-record rule (nearest setting, force trend or median) is the most decisive within the family in each of 200 draws, and the band of the nearest setting with a sibling grows from 0.25 to 2.9 floors. Selected on one geometry and evaluated on the other (\esm{Table~S13}), the most decisive rule is the nearest setting with a sibling (2.0 and 3.6 floors), }\DIFaddend the force trend \DIFdelbegin \DIFdel{, tied at zero cost; (ii) for a new setting, }\DIFdelend \DIFaddbegin \DIFadd{or the median interpolating (4.2--5.8) and }\DIFaddend the force trend \DIFdelbegin \DIFdel{at every trial cost, by three false rejects of }\DIFdelend \DIFaddbegin \DIFadd{or }\DIFaddend the nearest setting \DIFdelbegin \DIFdel{among 108 decisions; and (iii)}\DIFdelend \DIFaddbegin \DIFadd{extrapolating (8.9--11); }\DIFaddend for a new family \DIFdelbegin \DIFdel{, a trial of every design while a trial costs at most a tenth of a false accept, the envelope rule at a fifth and M5 beyond.
}%DIFDELCMD < 

%DIFDELCMD < \subsection{Robustness}\label{sec:res-robust}
%DIFDELCMD < 

%DIFDELCMD < %%%
\DIFdel{Table~\ref{tab:robust} repeats the analysis under variants of the data, constants and model specifications. Adjusting for punch temperature changes no calibrated distance by more than 2.6 floors. Dropping the warm-up of the first 150 parts shrinks the floors }\DIFdelend \DIFaddbegin \DIFadd{the rule selected on one transfer decides at 32--33 floors on the other, }\DIFaddend and \DIFdelbegin \DIFdel{so lengthens the distances, and floors taken between series shorten them.
Other floor protocols rescale }\DIFdelend \DIFaddbegin \DIFadd{the safest pre-specified rule, M3, is safe at 17--18. Paired intervals between point rules span about $\pm2$ floors with a sibling and when extrapolating, so smaller differences are not resolved by this design.
}

\DIFadd{Requiring all seven characteristics of a design to meet (\esm{Table~S14}) raises }\DIFaddend the decisive distances \DIFdelbegin \DIFdel{by 0.60 to 1.42, and the Gaussian-process specifications move them }\DIFdelend \DIFaddbegin \DIFadd{of the production-record rules }\DIFaddend by at most \DIFdelbegin \DIFdel{1.2 floors and the safe distances by at most 2.1. In every variant that includes them, the most decisive rule }\DIFdelend \DIFaddbegin \DIFadd{1.8 floors }\DIFaddend within a family \DIFdelbegin \DIFdel{is }\DIFdelend \DIFaddbegin \DIFadd{(}\DIFaddend the nearest setting \DIFdelbegin \DIFdel{or the force trend, and with a produced siblingthe nearest setting (tied once) . The refitting bootstrap widens the intervals of the learned rules most (}\DIFdelend \DIFaddbegin \DIFadd{5.1 floors with a sibling, the force trend 9.5 extrapolating) but those of }\DIFaddend M5 \DIFdelbegin \DIFdel{4.0--11 floors with a sibling, the nearest setting 2.3--4.6), yet in all 200 resamples (48 configurations) }\DIFdelend \DIFaddbegin \DIFadd{by up to 16, and to 39 floors for M5 for a new family. Lower limits on the 5th percentile, calibration and truths from the first 20 or 50 parts, and the decided statistic $\bar m+4\sigma$ (an index of 1.33) leave }\DIFaddend the nearest setting \DIFdelbegin \DIFdel{is more decisive than M5 and M5n, }\DIFdelend \DIFaddbegin \DIFadd{within 0.4 floor of its reference distance }\DIFaddend with a sibling and \DIFdelbegin \DIFdel{for a new setting alike, and the force trend than the scaled simulation (by 1.35--3.6 and 8.8--12 floors); point estimates at the edge of their refit intervals (M1s with a sibling at 4.9 floors, 4.9--7.3, and the safe distances of M2 and M5) reflect what losing a pattern does}\DIFdelend \DIFaddbegin \DIFadd{within 1.2 when extrapolating, and a production-record rule the most decisive within the family (\esm{Table~S2}), as do the robustness variants of data, constants and specifications; other floor protocols rescale the distances by 0.60 to 1.42, and floors recomputed without the held-out alternative's own series move the distances of the nearest setting and the force trend by at most 0.25 floors and those of the other main rules by at most 1.4}\DIFaddend .

%% ----------------------------------------------------------------------
\section{Discussion}\label{sec:discussion}

\subsection{What the evaluation establishes}

\DIFdelbegin \DIFdel{Five findings answer the research questions. Two }\DIFdelend \DIFaddbegin \DIFadd{Two findings }\DIFaddend follow from the definitions and hold wherever the framework is applied: the floor puts every rule and characteristic on a common production-referenced scale, though not on a common scale of nonconforming fractions, and as a within-run quantity it bounds the variation of series production from below (RQ1); and an interval rule trades decisiveness for safety, which the two distances make explicit\DIFdelbegin \DIFdel{, while its distribution-free intervals are guaranteed their level only where the calibration alternatives are exchangeable with the new design. Three }\DIFdelend \DIFaddbegin \DIFadd{. The others }\DIFaddend are observations on \DIFdelbegin \DIFdel{this dataset}\DIFdelend \DIFaddbegin \DIFadd{one laboratory process, with every alternative produced once and a generic simulation dataset, }\DIFaddend and hypotheses for confirmation\DIFdelbegin \DIFdel{: the best }\DIFdelend \DIFaddbegin \DIFadd{. The }\DIFaddend attainable distance grows with the novelty of the design\DIFdelbegin \DIFdel{, from }\DIFdelend \DIFaddbegin \DIFadd{: a production-record rule committed to in advance decided at }\DIFaddend 3.4 floors with a \DIFdelbegin \DIFdel{produced sibling (the nearest setting 0.85 for near-replicates}\DIFdelend \DIFaddbegin \DIFadd{sibling (1.6 when the two produced siblings did not differ resolvably from each other, 4.4 when they did)}\DIFaddend , \DIFdelbegin \DIFdel{4.9 for resolvable siblings ) to 8.2 }\DIFdelend \DIFaddbegin \DIFadd{8.5 }\DIFaddend for a new setting (\DIFdelbegin \DIFdel{the force trend; by stratum the nearest setting needed }\DIFdelend 4.7 \DIFdelbegin \DIFdel{interpolated and 4.9 extrapolated to 500\,kN, the force trend 8.9 to 100\,kN, where the nearest setting needed 11) and 33 }\DIFdelend \DIFaddbegin \DIFadd{interpolating, 9.5 extrapolating) and 30 floors at best }\DIFaddend for a new family, \DIFdelbegin \DIFdel{while }\DIFdelend \DIFaddbegin \DIFadd{and }\DIFaddend within a family the choice of rule \DIFdelbegin \DIFdel{matters }\DIFdelend \DIFaddbegin \DIFadd{mattered }\DIFaddend more than the relation\DIFaddbegin \DIFadd{, also among the pre-specified rules }\DIFaddend (RQ2)\DIFdelbegin \DIFdel{; production evidence buys safety before decisiveness, and interval rules held their level with a produced sibling but not when extrapolating or for a new family; and with this simulator and }\DIFdelend \DIFaddbegin \DIFadd{. For this simulation dataset and for Gaussian processes and boosting fitted to }\DIFaddend six to nine \DIFdelbegin \DIFdel{produced alternativesat two or three force levels}\DIFdelend \DIFaddbegin \DIFadd{alternatives}\DIFaddend , the simulation \DIFdelbegin \DIFdel{, as an offset or rescaled, contributed only }\DIFdelend \DIFaddbegin \DIFadd{contributed little within a family, helping single characteristics when extrapolating in force, and mattered }\DIFaddend across families, \DIFaddbegin \DIFadd{for the characteristics it models directly; the learners added no decisiveness, }\DIFaddend and \DIFdelbegin \DIFdel{M5, M5n and M3n predicted as well as the nearest produced setting with a produced sibling while adding abstention }\DIFdelend \DIFaddbegin \DIFadd{what made rules safe for new settings was an interval calibrated on the relation }\DIFaddend (RQ3). \DIFaddbegin \DIFadd{Qualification by relation also exposed a dangerous extrapolation: the scaled simulation, decisive with a sibling, falsely accepted failing alternatives up to 165 floors from their limit when extrapolating in force.
}\DIFaddend 

%DIF > % guidance table (static)
\begin{table}[t]
\tabsize
\setlength{\tabcolsep}{4pt}
\caption{Guidance for design-stage manufacturability decisions, with its basis}\label{tab:guidelines}
\DIFdelbeginFL %DIFDELCMD < \begin{tabular}{@{}L{2.9cm}L{7.75cm}L{1.85cm}@{}}
%DIFDELCMD < \toprule
%DIFDELCMD < Situation & Guidance & Basis \\
%DIFDELCMD < \midrule
%DIFDELCMD < \multicolumn{3}{@{}l}{\emph{From the framework (prescriptive; usefulness not evaluated)}} \\
%DIFDELCMD < Before relying on a predictor & Measure the floor on consecutive production of the family with a stated protocol and report $\sigma_\mathrm{b}$, $\sigma_\mathrm{w}$ and its lag dependence; treat it as a lower bound of series variation. & framework \\
%DIFDELCMD < Qualifying a rule & Evaluate it for the evidence relation of the coming designs by grouped cross-validation, and split near-replicate siblings from resolvable ones. & framework \\
%DIFDELCMD < Choosing a rule & Compare it with the nearest produced setting and with its own version without the simulation, on decisive and safe distance jointly or by expected cost. & framework \\
%DIFDELCMD < Deciding a single design & Derive the rule's guard band for the distribution of requirement distances the team expects and apply it to the observable margin, in place of range and family guards; a margin inside it means a trial. & framework \\
%DIFDELCMD < Planning a requirement & As a heuristic, plan a trial if the expected margin between requirement and predicted 95th percentile is below the safe distance of the best qualified rule; no decision is reliable closer than a trial reproduces the quantile. & framework \\
%DIFDELCMD < \multicolumn{3}{@{}l}{\emph{Observed on this dataset (18 alternatives, 2 geometries, 7 characteristics)}} \\
%DIFDELCMD < New variant with a produced sibling & nearest produced setting right beyond 3.4 floors (0.85 for near-replicates, 4.9 for resolvable siblings); interval rules safe below one floor, coverage 85--96\,\% & 126 cases \\
%DIFDELCMD < New process setting & nearest setting right beyond 4.7 floors interpolated and 4.9 extrapolated to 500\,kN, force trend 8.9 to 100\,kN (nearest setting 11; speed change); interval rules safe at 5.9--16 floors, coverage 43--55\,\% & 42 + 84 cases \\
%DIFDELCMD < New family & best rule safe at 16--17 floors in either family's floor, which rule not resolved; rules without the simulation fail; a transferable drawing nominal helps & two transfers \\
%DIFDELCMD < Release-level requirements & lie 1.7 floors above the 95th percentile at an index of 1.33, inside every decisive distance except a near-replicate's; decided only with a produced sibling & 126 cases \\
%DIFDELCMD < \botrule
%DIFDELCMD < \end{tabular}
%DIFDELCMD < %%%
\DIFdelendFL \DIFaddbeginFL \begin{tabular}{@{}L{2.9cm}L{7.75cm}L{1.85cm}@{}}
\toprule
Situation & Guidance & Basis \\
\midrule
\multicolumn{3}{@{}l}{\emph{From the framework (prescriptive; usefulness not evaluated)}} \\
Before relying on a predictor & Measure the floor on consecutive production of the family with a stated protocol; report $\sigma_\mathrm{b}$, $\sigma_\mathrm{w}$, its lag dependence and the floors of single runs; treat it as a lower bound of series variation. & framework \\
Qualifying a rule & Evaluate it by grouped cross-validation for the evidence relation of the coming designs, per family, with siblings stratified by whether they differ from each other, and per characteristic. & framework \\
Choosing a rule & Compare it with the nearest produced setting and with its own version without the simulation, on both distances and by expected cost at the team's trial cost. & framework \\
Deciding a single design & Derive the guard band for the requirements the team expects, validate it out of sample, and apply it to the observable margin; a margin inside it means a trial. & framework \\
\multicolumn{3}{@{}l}{\emph{Observed in this study (one process, 18 alternatives produced once, 2 geometries; not general)}} \\
New variant with a produced sibling & the nearest produced setting decided beyond 3.4 floors (1.6 when the produced siblings did not differ resolvably from each other); no learned rule was more decisive & 18 alternatives \\
New process setting & production record 8.2--8.5 floors, interval rules 15--19; intervals matched to the relation made it safe & 6 force cells \\
New family & only rules with the (generic) simulation decided, at 30 floors and beyond; none qualified for single verdicts & two transfers \\
Release-level requirements & within about 2 floors of the 95th percentile: decided at best with a near-replicate sibling, otherwise by trial & 18 alternatives \\
\botrule
\end{tabular}
\DIFaddendFL \end{table}

\DIFdelbegin %DIFDELCMD < \subsection{Using the framework in design}
%DIFDELCMD < %%%
\DIFdelend \DIFaddbegin \subsection{Implications for practice}
\DIFaddend 

\DIFdelbegin \DIFdel{As a planning heuristic (Table~\ref{tab:guidelines}), the expected margin between a requirement and the predicted 95th percentile can be compared with the safe distance of the best rule qualified for the relation; }\DIFdelend \DIFaddbegin \DIFadd{For a process planner or quality engineer the results say three things. Within a family the production record answers most design-stage questions }\DIFaddend if it is \DIFdelbegin \DIFdel{smaller, a trial or more production evidence has to be planned, and the decision itself needs the guard band, because the safe distance is defined on the true margin. Release-level requirements make this the usual case, as the practice of a production trial run for a process change expects: the rules decide reliably where a trial is cheapest, a new force on an existing tool being a short press run, and none does near }\DIFdelend \DIFaddbegin \DIFadd{indexed by how the coming design relates to it: a sibling at the same setting settles a requirement a few floors away, a new force within the produced range almost as well, and an extrapolated force or a new family needs a trial unless the requirement is far. Under the reference prior, of 100 verdicts on new variants with a sibling the nearest setting with its band sends 2 to trial and gets at least 93 of the other 98 right; of 100 on extrapolated settings it sends 55 to trial. And near }\DIFaddend release-level \DIFdelbegin \DIFdel{limits where a trial is dearest, for a new family that may need a new tool. The most decisive rule is arguably }\DIFdelend \DIFaddbegin \DIFadd{limits nothing replaced a trial without a nearly identical produced sibling, so the value of design-stage prediction lies in the many decisions far from the limits rather than in the few close ones. For a developer of an AI-based manufacturability tool the lesson is methodological: a learned model has to be compared with a production-record baseline and with its own version without the simulation, scored by evidence relation in a unit that production defines, its intervals calibrated on cases of the same relation and shown to its user with a guard band. The nearest setting is }\DIFaddend also the most transparent \DIFdelbegin \DIFdel{: the nearest produced setting points }\DIFdelend \DIFaddbegin \DIFadd{rule, pointing }\DIFaddend to the produced alternatives whose record it averages\DIFdelbegin \DIFdel{. The variants that steps 1 to 5 need come from process development and tryout, whose short series of several settings support the sibling and new-setting relations, and from a family's production history across part numbers, which supports the new-family relation; production part approval, whose initial process study covers the released setting, supplies the floor but not the variants}\DIFdelend \DIFaddbegin \DIFadd{, and so the natural default where a learned rule shows no advantage}\DIFaddend .

\subsection{Simulation and machine learning in this regime}

Within a design family the production record carried the decision \DIFdelbegin \DIFdel{, }\DIFdelend even against a \DIFdelbegin \DIFdel{rescaled simulation. The simulated force trend of the two draw-ins was 2.7 to 5.1 }\DIFdelend \DIFaddbegin \DIFadd{calibrated or rescaled simulation, but the simulation here is a generic dataset that does not contain the produced process: its force trend is two to five }\DIFaddend times too large, \DIFdelbegin \DIFdel{and calibrating the friction coefficient did not repair it, which points to the friction law and the blank-holder model rather than to a parameter. A multi-task model cannot help a family with no produced alternative, because the correlation between families cannot be estimated without its data; in the pooled scopes the Gaussian process with a 0/1 geometry input is a restricted intrinsic coregionalisation model (equal variances, one positive correlation between the families, shared length scales), which improved coverage but not decisiveness. These findings concern one simulator of low fidelity, used through offsets and one scale, and learners }\DIFdelend \DIFaddbegin \DIFadd{its friction calibration is censored at an edge of the grid with or without a discrepancy, and its oil-to-friction mapping is too sensitive. A process-calibrated simulation, with friction depending on pressure \mbox{%DIFAUXCMD
\citep{hol2012advanced}}\hskip0pt%DIFAUXCMD
, a hardening model fit for springback \mbox{%DIFAUXCMD
\citep{wagoner2013advanced} }\hskip0pt%DIFAUXCMD
and elastic tools \mbox{%DIFAUXCMD
\citep{pilthammar2018introduction}}\hskip0pt%DIFAUXCMD
, might contribute more; the per-characteristic results show where it already helps, for the draw-ins and the arm angle when extrapolating and for the characteristics it models directly in a new family. The learners failed for a reason the framework exposes: }\DIFaddend with six to nine \DIFdelbegin \DIFdel{produced }\DIFdelend alternatives at two or three force levels \DIFdelbegin \DIFdel{. For evaluating AI-based DFM support the lesson is methodological: a learned model should be compared with a production-record baseline, scored by evidence relation in a unit that production defines, and shown to its user with a guard band. A study with design teams could test two hypotheses on reliance \mbox{%DIFAUXCMD
\citep{lee2004trust}}\hskip0pt%DIFAUXCMD
, with cases of known production outcome: designers shown the guard-banded verdict accept a smaller share of wrong verdicts within the safe distance than designers shown the verdict alone, at an equal rate of right acceptances; and, as the same AI assistance helped some teams and hindered others \mbox{%DIFAUXCMD
\citep{zhang2021cautionary}}\hskip0pt%DIFAUXCMD
, the guard band narrows the spread of that share between teams.
}%DIFDELCMD < 

%DIFDELCMD < \subsection{Implications for engineering design}
%DIFDELCMD < 

%DIFDELCMD < %%%
\DIFdel{For variant and platform design \mbox{%DIFAUXCMD
\citep{pahl2007engineering, simpson2006product} }\hskip0pt%DIFAUXCMD
and the capability databases that serve them \mbox{%DIFAUXCMD
\citep{tata1999process, thornton2004variation}}\hskip0pt%DIFAUXCMD
, the results argue for indexing production records by their relation to the coming design rather than for a more sophisticated model on top. The safe distance is the part of a margin that covers predictive uncertainty, as distinct from the part that absorbs change \mbox{%DIFAUXCMD
\citep{eckert2017safety, eckert2019design}}\hskip0pt%DIFAUXCMD
. The growth of the decisive distance, pooled over strata, from about 3 to 8 and over 30 floors says for which overlaps of development \mbox{%DIFAUXCMD
\citep{krishnan1997model, terwiesch2002exchanging} }\hskip0pt%DIFAUXCMD
preliminary manufacturability information can be released, the new-family safe distance, some 0.8\,mm of draw-in, how far a set-based design \mbox{%DIFAUXCMD
\citep{sobek1999toyota} }\hskip0pt%DIFAUXCMD
can narrow a feasible set before production evidence exists, and the distances and trial rates when a test is cheaper than a decision \mbox{%DIFAUXCMD
\citep{thomke2001sequential, salado2018mathematical}}\hskip0pt%DIFAUXCMD
. The evaluation shows structural validity, that the metrics discriminate between rules and relations on real data \mbox{%DIFAUXCMD
\citep{seepersad2006validation}}\hskip0pt%DIFAUXCMD
; performance validity would require better decisions with the framework}\DIFdelend \DIFaddbegin \DIFadd{a flexible model can identify at most a linear trend, and a well-specified one then becomes that trend. Nonlinear multi-fidelity models \mbox{%DIFAUXCMD
\citep{perdikaris2017nonlinear}}\hskip0pt%DIFAUXCMD
, physics-informed learning \mbox{%DIFAUXCMD
\citep{karniadakis2021physics}}\hskip0pt%DIFAUXCMD
, conformalised quantile regression on the parts \mbox{%DIFAUXCMD
\citep{romano2019conformalized} }\hskip0pt%DIFAUXCMD
and multi-task models over more families \mbox{%DIFAUXCMD
\citep{bonilla2008multi} }\hskip0pt%DIFAUXCMD
are the next candidates, to be scored in the same way}\DIFaddend .

\DIFdelbegin %DIFDELCMD < \begin{table}[t]
%DIFDELCMD < \tabsize
%DIFDELCMD < \setlength{\tabcolsep}{3pt}
%DIFDELCMD < \caption{What the framework requires, mapped onto other processes (an untested proposal)}\label{tab:processes}
%DIFDELCMD < \begin{tabular}{@{}L{1.9cm}L{2.4cm}L{2.5cm}L{2.7cm}L{2.7cm}@{}}
%DIFDELCMD < \toprule
%DIFDELCMD < Process & Batch and run & Main variation & Evidence relations & Floor \\
%DIFDELCMD < \midrule
%DIFDELCMD < Sheet forming (this paper) & consecutive parts of one series & drift in a run; coil, die state & variant, setting, family & batch centres in a run \\
%DIFDELCMD < Injection moulding & shots of one cavity and set-up & cavity, shot, set-up, material lot & cavity, process window, new mould & per cavity; cavity offsets as effects \\
%DIFDELCMD < Machining & consecutive parts between offsets & tool wear, offset-compensated & tool path, cutting data, new part & batch centres within an offset interval \\
%DIFDELCMD < Casting & castings of one pattern and melt & melt, pattern position & alloy, gating, new pattern & within a melt; melts as runs \\
%DIFDELCMD < Additive manufacturing & replicates in one build & build, position, orientation & orientation, settings, new part & within a build, via replicated anchors \\
%DIFDELCMD < \botrule
%DIFDELCMD < \end{tabular}
%DIFDELCMD < \end{table}
%DIFDELCMD < 

%DIFDELCMD < \subsection{Reuse for other processes and AI-based DFM tools}%%%
\DIFdelend \DIFaddbegin \subsection{Reuse for other processes and AI-based tools: a proposal}\DIFaddend \label{sec:reuse}

The framework needs several produced variants of a design family whose relations span those of the coming designs, a series of consecutive parts or replicated batches \DIFdelbegin \DIFdel{of one design }\DIFdelend per variant, a design-stage predictor for every variant \DIFdelbegin \DIFdel{, }\DIFdelend and characteristics measured identically on parts and predictions; here \DIFdelbegin \DIFdel{rules calibrated on about }\DIFdelend \DIFaddbegin \DIFadd{nine produced alternatives per family and 500-part series gave ten batches of }\DIFaddend 50 \DIFdelbegin \DIFdel{parts per variant scored within 1.5 floors of those calibrated on full series , against truths and floors from the full series. Attribute outcomes such as splits need a requirement on a rate, which the framework does not yet cover. Other predictors need three adaptations}\DIFdelend \DIFaddbegin \DIFadd{parts for the floor, one produced sibling sufficed for the nearest setting, and truths from the first 20 or 50 parts moved its distance with a sibling by less than half a floor. The extensions below are untested}\DIFaddend . A classifier with a fixed threshold does not take the requirement as input, so its operating characteristic is estimated by binning test cases by their true margin, \DIFdelbegin \DIFdel{not by the exact estimator of Section~\ref{sec:distances}, unless the requirement is made an input; a probabilistic classifier needs two thresholds to return a trial}\DIFdelend \DIFaddbegin \DIFadd{which needs a continuous, production-referenced outcome for every test design that CAD-based manufacturability datasets usually lack}\DIFaddend . A language-model check \DIFdelbegin \DIFdel{may return verdicts that are not }\DIFdelend \DIFaddbegin \DIFadd{is scored on a grid of requirements with repeated queries, reporting model and version, query date, decoding settings and prompt paraphrases, treating refusals and hedged answers as trials, checking that its verdicts are }\DIFaddend monotone in the requirement and \DIFdelbegin \DIFdel{that vary between queries, so $\kappa$ and $\omega$ are computed on a requirement grid with repeated queries, and }\DIFdelend \DIFaddbegin \DIFadd{testing for contamination, since public datasets and the papers describing them may be in its training data. A requirement on a rate, such as }\DIFaddend the share of \DIFdelbegin \DIFdel{non-monotone verdict sequences is reported. For }\DIFdelend \DIFaddbegin \DIFadd{split or wrinkled parts, is scored in the same way once each alternative's rate is estimated from its series, with a floor over batches of rates; for }\DIFaddend continuous design spaces the relation becomes a distance from the produced designs in descriptor space\DIFdelbegin \DIFdel{or the position relative to their hull, and each stratum needs about as many cases as a relation here: with 126 cases, some 240 verdicts, the 95\,\% threshold already rests on about a dozen. A generative tool's manufacturability check is scored. A minimal report lists
}%DIFDELCMD < \begin{enumerate}
%DIFDELCMD < \item %%%
\DIFdel{the floor with its protocol, variance components and lag dependence;
}%DIFDELCMD < \item %%%
\DIFdel{decisive and safe distancesby relation, with refit intervals and one-sided values;
}%DIFDELCMD < \item %%%
\DIFdel{coverage and half-width of interval rules, and the source of their noise variance;
}%DIFDELCMD < \item %%%
\DIFdelend \DIFaddbegin \DIFadd{. Table~\ref{tab:checklist} is a proposed reporting checklist, modelled on model cards \mbox{%DIFAUXCMD
\citep{mitchell2019model} }\hskip0pt%DIFAUXCMD
and TRIPOD+AI \mbox{%DIFAUXCMD
\citep{collins2024tripod}}\hskip0pt%DIFAUXCMD
; the scoring is released as a small documented function, separate from the pipeline, that takes intervals, truths and floors and returns distances, guard bands and trial rates \mbox{%DIFAUXCMD
\citep{canseven2026archive}}\hskip0pt%DIFAUXCMD
. An attempt on open powder-bed fusion data \mbox{%DIFAUXCMD
\citep{leirmo2021data} }\hskip0pt%DIFAUXCMD
failed for lack of replicated batches within a build: a floor within a build would need at least ten replicated anchors of one design per build and several builds per setting.
}

\begin{table}[t]
\tabsize
\setlength{\tabcolsep}{3pt}
\caption{Proposed reporting checklist for evaluating a design-stage decision rule against production}\label{tab:checklist}
\begin{tabular}{@{}L{3.3cm}L{4.6cm}L{4.4cm}@{}}
\toprule
Item & Why it matters & In this paper \\
\midrule
Floor and its protocol & every distance is in this unit & $B$, $\alpha$, $\sigma_\mathrm{b}$, $\sigma_\mathrm{w}$, lag, single runs (Table~\ref{tab:floors}) \\
Measurement of the characteristics & biases enter absolute comparisons & Table~\ref{tab:qcs}, Section~\ref{sec:qcs} \\
Relations, cases and replication units & alternatives, not cases, replicate & per relation, 18 alternatives \\
Both distances, false-accept side & decisiveness and safety differ & Table~\ref{tab:distances} \\
Baselines & production record, no simulation & NN, M1sn; every pair with and without \\
Intervals: source and level & coverage is a property of the construction & Tables~\ref{tab:rules}, \ref{tab:crossing} \\
Per-characteristic results & pooled distances hide the worst & Table~\ref{tab:groups} \\
Guard band, prior, validation & the deployed part & per family, out of sample (Table~\ref{tab:guards}) \\
Costs and trial rates & deployment is a cost question & Fig.~\ref{fig:costs} \\
Development record & selection optimism & Section~\ref{sec:computation}, cross-fitting \\
Query protocol (language models) & verdicts vary between queries & proposal (Section~\ref{sec:reuse}) \\
\botrule
\end{tabular}
\end{table}

\subsection{Deployment}

\DIFadd{In use, a step-6 verdict would be shown with its margin in floors and in the unit of }\DIFaddend the \DIFdelbegin \DIFdel{nearest-produced-setting }\DIFdelend \DIFaddbegin \DIFadd{characteristic, the relation on which it rests and the produced alternatives behind it, }\DIFaddend and \DIFdelbegin \DIFdel{no-simulation baselines;
}%DIFDELCMD < \item %%%
\DIFdel{the guard bands with their prior, the trial rates and the performance of the guards;
}%DIFDELCMD < \item %%%
\DIFdel{for language-model tools, the prompt and the query protocol. }%DIFDELCMD < \end{enumerate}
%DIFDELCMD < %%%
\DIFdel{Table~\ref{tab:processes}, an untested proposal, maps these requirements onto other processes; distances in floors are comparable within a protocol, not across processes, since a floor between builds, cavities or offset intervals is a between-unit quantity (here the floor between series is 1.5 }\DIFdelend \DIFaddbegin \DIFadd{a trial verdict with the requirement distance at which the rule would decide. A qualification has to be maintained: the realised errors of step-6 verdicts can be tracked against later production on a control chart of decision outcomes, and a rule re-qualified when its share of wrong decided verdicts exceeds 5\,\% beyond what the number of verdicts explains, when the floor of new runs leaves the range of the single-run floors, or when the die is reworked. The evidence has a cost, several 500-part series per family before any rule qualifies, which the results bound from below: one produced sibling settles most verdicts with a sibling, and one produced alternative of a new family brings its decisive distance from 28 }\DIFaddend to \DIFdelbegin \DIFdel{9.1 times that within a run). An attempt on open powder bed fusion data of PA12 \mbox{%DIFAUXCMD
\citep{leirmo2021data} }\hskip0pt%DIFAUXCMD
failed on the second requirement: orientation classes mixed orientations in a batch, eight identical anchors per build were too few for a floor within a build, and three builds give three differences between builds}\DIFdelend \DIFaddbegin \DIFadd{9 floors in its own floor. Whether teams decide better with the framework is an empirical question: a study with cases of known production outcome could test whether designers shown the guard-banded verdict accept fewer wrong verdicts within the safe distance than designers shown the verdict alone, at an equal rate of right acceptances \mbox{%DIFAUXCMD
\citep{lee2004trust}}\hskip0pt%DIFAUXCMD
, and whether the band narrows the spread of that share between teams \mbox{%DIFAUXCMD
\citep{zhang2021cautionary}}\hskip0pt%DIFAUXCMD
}\DIFaddend .

\subsection{Threats to validity}

The evaluation covers one material, one press and two geometries, with relations within a family from a $3\times3$ factorial in which lubrication is mostly not resolved; the procedure is general, the numbers are not. The new family rests on two transfers between geometries that differ in three features at once\DIFdelbegin \DIFdel{, and the intervals of its distances hold the two fits fixed}\DIFdelend . Each alternative was produced once, so the \DIFdelbegin \DIFdel{floor covers drift within a run but not between days, coils or tool states, the reference $q_{95}$ }\DIFdelend \DIFaddbegin \DIFadd{reference $q_p$ }\DIFaddend is that of one run \DIFdelbegin \DIFdel{, and the order and dates of the series are unknown; between-series variation bounds the missing component at 1.2--9.8 floors, }\DIFdelend and \DIFdelbegin \DIFdel{scanner drift cannot be separated }\DIFdelend \DIFaddbegin \DIFadd{the variation between days, coils and tool states is missing; the sensitivity analysis bounds its effect but cannot replace replicated runs. The floor does not separate scanner drift }\DIFaddend from process drift\DIFdelbegin \DIFdel{. The extrapolation to 100\,kN is confounded with stroke speed, which the simulation omits}\DIFdelend \DIFaddbegin \DIFadd{, and absolute values are uncalibrated. The simulation dataset was not calibrated to the process, so the findings on simulation concern this dataset}\DIFaddend . The requirements are one-sided limits on the 95th percentile of \DIFdelbegin \DIFdel{a characteristic averaged over sides}\DIFdelend \DIFaddbegin \DIFadd{dimensional characteristics}\DIFaddend , below the conformance of series release\DIFdelbegin \DIFdel{. Several rules, the grouped schemes, the guard bands and their prior were added or chosen after earlier }\DIFdelend \DIFaddbegin \DIFadd{; attribute outcomes such as splits and wrinkles are not covered; and the verdicts are marginal per characteristic, a design-level verdict being an extension. The guard band was estimated pooled over characteristics, its prior is a reference rather than a team's, and the costs are assumed. Rules, strata, bands and priors were partly chosen after }\DIFaddend results had been seen and are evaluated on the same data\DIFdelbegin \DIFdel{, so the comparison between rules carries selection optimism \mbox{%DIFAUXCMD
\citep{cawley2010over}}\hskip0pt%DIFAUXCMD
; }\DIFdelend \DIFaddbegin \DIFadd{; the cross-fitting bounds the optimism but does not remove it, and }\DIFaddend a single analyst carried out the analysis, with code and data open for replication.

%% ----------------------------------------------------------------------
\section{Conclusions}\label{sec:conclusion}

Design-stage manufacturability decisions can be evaluated against production on a scale production defines: the decisive and safe distances in floors tell how close to a requirement a rule decides and how close it can be relied on not to mislead, and a guard band derived \DIFaddbegin \DIFadd{and validated }\DIFaddend for the requirements a team expects turns them into a rule for single decisions\DIFdelbegin \DIFdel{. On 18 deep-drawn alternatives with matched simulations, the best }\DIFdelend \DIFaddbegin \DIFadd{, priced by the costs of errors and trials. In an exploratory application to one laboratory deep-drawing process, with alternatives produced once and a generic simulation dataset, the }\DIFaddend attainable distance grew with \DIFdelbegin \DIFdel{novelty: }\DIFdelend \DIFaddbegin \DIFadd{the novelty of the design: the nearest produced setting, committed to in advance, decided at }\DIFaddend 3.4 floors with a produced sibling \DIFdelbegin \DIFdel{(the nearest setting 0.85 for near-replicates, 4.9 for resolvable ones), 8.2 }\DIFdelend \DIFaddbegin \DIFadd{and 8.5 }\DIFaddend for a new process setting\DIFdelbegin \DIFdel{(the force trend; the nearest setting 4.7 interpolated and 4.9 extrapolated to 500\,kN, the force trend 8.9 to 100\,kN, with a stroke-speed change)}\DIFdelend , and no rule \DIFdelbegin \DIFdel{safe closer than 16--17 }\DIFdelend \DIFaddbegin \DIFadd{decided closer than 30 }\DIFaddend floors for a new family \DIFdelbegin \DIFdel{, in either family's floor. Within a family the choice of rule mattered more than the relation; requirements at a performance index of 1.33 lay inside every decisive distance except a near-replicate's and need a trial without a produced sibling ; and in this small-data regime the simulation as used, offset or rescaled, contributed only across families, while learned 90\,\% intervals covered 85--96\,\% of truths with a producedsibling, under 50\,\% extrapolating or for a new family}\DIFdelend \DIFaddbegin \DIFadd{(the pre-specified M5 at 35). Release-level requirements lay within these distances, so they needed a trial unless a near-replicate sibling had been produced}\DIFaddend .

For design practice, the production record carried the decision within a family, which argues for indexing such records by \DIFaddbegin \DIFadd{their }\DIFaddend relation to the coming design rather than for a more elaborate model; across families only rules using the simulation carried the change of geometry\DIFdelbegin \DIFdel{, and none decided near release-level limits}\DIFdelend . For AI-based manufacturability support, the \DIFdelbegin \DIFdel{framework offers a standard of evaluation rather than a verdict on machine learning: a tool is scored by evidence relation in floors against a production-record baseline and reported with the checklist of Section~\ref{sec:reuse}}\DIFdelend \DIFaddbegin \DIFadd{evaluation separated what the learners were credited with: Gaussian processes and boosting fitted to six to nine alternatives reproduced the production record, and intervals calibrated on the same evidence relation, not learning, made it safe for new settings}\DIFaddend .

The evidence is one process and two geometries, and the comparison of rules is exploratory. Four extensions follow: series from several days, coils and tool states; \DIFaddbegin \DIFadd{a process-calibrated simulation; }\DIFaddend families with continuous geometric descriptors; \DIFaddbegin \DIFadd{and }\DIFaddend a confirmatory test of the frozen procedure on other data\DIFdelbegin \DIFdel{; and a study of whether design teams decide better with it}\DIFdelend \DIFaddbegin \DIFadd{. Within a design family, the baseline that a learned manufacturability check has to beat is a production record indexed by its relation to the coming design}\DIFaddend .

%% ----------------------------------------------------------------------

\FloatBarrier
\begin{appendices}
\renewcommand{\floatpagefraction}{0.4}
\appendixfloatpages%DIF <  Tables A1--A2 and A3--A5 each fill a float page
%DIF >  the appendix tables fill float pages
\renewcommand{\theHtable}{\thetable}%DIF <  distinct hyperlink targets for Tables A1 and A2
%DIF >  distinct hyperlink targets for the appendix tables

\DIFdelbegin %DIFDELCMD < \section{Notation, further results and computation details}%%%
\DIFdelend \DIFaddbegin \section{Notation, further rules and results, and computation details}\DIFaddend \label{app:appendix}

Table~\ref{tab:notation} summarises the notation\DIFdelbegin \DIFdel{of Section~\ref{sec:method}}\DIFdelend \DIFaddbegin \DIFadd{, Table~\ref{tab:further} defines and scores the further rules}\DIFaddend , Table~\DIFdelbegin \DIFdel{\ref{tab:robust} gives the distances under the robustness variants}\DIFdelend \DIFaddbegin \DIFadd{\ref{tab:gp} gives the fitted Gaussian processes}\DIFaddend , Table~\DIFdelbegin \DIFdel{\ref{tab:intervals} the coverage, width and centre of the interval rules}\DIFdelend \DIFaddbegin \DIFadd{\ref{tab:crossing} the predictors with relation-matched intervals, Table~\ref{tab:groups} the distances per characteristic}\DIFaddend , and Table~\DIFdelbegin \DIFdel{\ref{tab:step6} the procedure's decision rule by relation}\DIFdelend \DIFaddbegin \DIFadd{\ref{tab:validation} the simulation against the measurement}\DIFaddend . Table~\ref{tab:mapping} traces every other number quoted in the text to the script that computes it and the result file that holds it in the archived release \citep{canseven2026archive}, in which \texttt{scripts/run\_all.sh} runs the scripts in order and \texttt{make\_tables.py} and \texttt{figs.py} write every table and figure\DIFdelbegin \DIFdel{from the result files}\DIFdelend \DIFaddbegin \DIFadd{; Online Resource~1 holds the further tables (S1--S20) and figures}\DIFaddend .

%DIF <  the appendix tables come first in the source, so that they fill the two float pages that follow the page
%DIF >  the appendix tables come first in the source, so that they fill the float pages that follow the page
% on which the appendix begins
\begin{table}[p]
%DIF <  float page with Table A2
\tabsizesmall
\setlength{\tabcolsep}{4pt}
\caption{Notation and terms}\label{tab:notation}
\DIFdelbeginFL %DIFDELCMD < \begin{tabular}{@{}L{3.0cm}L{9.75cm}@{}}
%DIFDELCMD < \toprule
%DIFDELCMD < Symbol, term & Meaning \\
%DIFDELCMD < \midrule
%DIFDELCMD < $a$, $c$, $q_p(a,c)$ & design--process alternative, characteristic, $p$-quantile ($p=0.95$) of its parts \\
%DIFDELCMD < $R$, $d$ & requirement (upper limit on $c$ or $|c|$), $d$ floors from the truth: $R=q_p+dF_c$ \\
%DIFDELCMD < $[L_r,U_r]$, $V_r$ & interval of rule $r$ for $q_p$; verdict meets, fails or trial (Eq.~\ref{eq:verdict}) \\
%DIFDELCMD < $B$, $m$, $F_c$ & batch size (50); batch centre (10\,\% trimmed mean); production floor (Eq.~\ref{eq:floor}) \\
%DIFDELCMD < $\sigma_\mathrm{b}$, $\sigma_\mathrm{w}$ & standard deviation between batch means and between parts within a batch \\
%DIFDELCMD < ESR, MRC & effect-to-scatter ratio (Eq.~\ref{eq:esr}); minimum resolvable change, $F_c/|s_c|$ \\
%DIFDELCMD < $u$, $v$ & distances in floors by which an interval lies above or below the truth \\
%DIFDELCMD < $\kappa(\delta)$, $\omega(\delta)$ & shares of correct and wrong verdicts among the feasible verdicts at $|d|=\delta$ \\
%DIFDELCMD < $\delta_\mathrm{d}$, $\delta_\mathrm{s}$ & decisive distance ($\kappa\ge0.95$ beyond it) and safe distance ($\omega\le0.05$ beyond it) \\
%DIFDELCMD < FA, FR side & distances on the false-accept (failing) and false-reject (adequate) side \\
%DIFDELCMD < $g$ & guard band: observable margin beyond which ${\ge}95$\,\% of decided verdicts were correct, under a stated prior of requirement distances, here the 47 grid distances within $\pm20$ floors, equal weights (Section~\ref{sec:procedure}); not a property of rule and relation alone \\
%DIFDELCMD < Calibration alternatives & produced alternatives whose record calibrates a rule \\
%DIFDELCMD < Evidence relation & sibling (same setting), new setting (interpolated, extrapolated), new family \\
%DIFDELCMD < Calibration budget & distances against the number $k$ of produced alternatives, by relation \\
%DIFDELCMD < Guard & family, range and novelty checks that flag a design outside the calibration \\
%DIFDELCMD < Measurement resolution & step between adjacent values that the measuring system reports (gauge sense) \\
%DIFDELCMD < \botrule
%DIFDELCMD < \end{tabular}
%DIFDELCMD < %%%
\DIFdelendFL \DIFaddbeginFL \begin{tabular}{@{}L{3.0cm}L{9.75cm}@{}}
\toprule
Symbol, term & Meaning \\
\midrule
$a$, $c$, $q_p(a,c)$ & design--process alternative, characteristic, $p$-quantile ($p=0.95$) of its parts \\
$R$, $d$ & requirement (upper limit on $c$ or $|c|$), $d$ floors from the truth: $R=q_p+dF_c$ \\
$[L_r,U_r]$, $V_r$ & interval of rule $r$ for $q_p$; verdict meets, fails or trial (Eq.~\ref{eq:verdict}) \\
$c_\mathrm{FA}$, $c_\mathrm{FR}$, $c_\mathrm{T}$ & costs of a false accept, a false reject and a trial (Eq.~\ref{eq:cost}) \\
$B$, $m$, $F_c$ & batch size (50); batch centre (10\,\% trimmed mean); within-run floor (Eq.~\ref{eq:floor}) \\
$\sigma_\mathrm{b}$, $\sigma_\mathrm{w}$, $\sigma_\mathrm{LT}$ & standard deviations between batch means, within batches, and of the parts of a series \\
ESR & effect-to-scatter ratio $|\bar m_a-\bar m_b|/F_c$; resolvable when at least one \\
$\kappa(\delta)$, $\omega(\delta)$ & shares of correct and wrong verdicts among the feasible verdicts at $|d|=\delta$ \\
$\delta_\mathrm{d}$, $\delta_\mathrm{s}$ & decisive distance ($\kappa\ge0.95$ beyond it) and safe distance ($\omega\le0.05$ beyond it) \\
FA side & distance for failing alternatives only, where an error is a false accept \\
$g$ & guard band: observable margin beyond which ${\ge}95$\,\% of decided verdicts were correct, under a stated prior of requirement distances \\
Evidence relation & sibling (same setting), new setting (interpolated, extrapolated), new family \\
Decision-time stratum & siblings resolvable from each other or not, known before deciding \\
P+RM & predictor P with a relation-matched split-conformal interval borrowed from the other family \\
$\tau$ & standard deviation (floors) of a between-run component added to every truth \\
\botrule
\end{tabular}
\DIFaddendFL \end{table}

%DIF < % BEGIN GENERATED robust
%DIF > % BEGIN GENERATED further
\begin{table}[p]
\tabsizesmall
\DIFdelbeginFL %DIFDELCMD < \setlength{\tabcolsep}{1.4pt}
%DIFDELCMD < \caption{Decisive/safe distances (floors) under the robustness variants}\label{tab:robust}
%DIFDELCMD < \begin{tabular}{@{}lrrrrrrrrrr@{}}
%DIFDELCMD < \toprule
%DIFDELCMD <  & \multicolumn{4}{c}{New variant} & \multicolumn{4}{c}{New setting} & \multicolumn{2}{c}{New family} \\
%DIFDELCMD < \cmidrule(lr){2-5}\cmidrule(lr){6-9}\cmidrule(lr){10-11}
%DIFDELCMD < Variant & M1s & M2 & M5 & NN & M2 & M5 & NN & M1sn & M2 & M5 \\
%DIFDELCMD < \midrule
%DIFDELCMD < Reference & 4.9 & 23/0.15 & 7.2/0.75 & 3.4 & 18/8.1 & 19/6.8 & 8.5 & 8.2 & 51/16 & 38/28 \\
%DIFDELCMD < Temperature-adjusted & 5.0 & 23/0.10 & 7.2/0.65 & 3.4 & 17/7.6 & 19/6.8 & 9.4 & 8.9 & 48/15 & 37/27 \\
%DIFDELCMD < Warm-up dropped & 6.5 & 27/0.20 & 9.0/0.45 & 4.1 & 21/8.7 & 24/8.7 & 12 & 8.8 & 63/18 & 39/29 \\
%DIFDELCMD < Floor between series & 2.9 & 11/0 & 1.3/0.20 & 0.75 & 9.9/2.6 & 11/2.3 & 2.7 & 4.5 & 16/3.8 & 12/9.3 \\
%DIFDELCMD < Floor: batches of 25 & 4.6 & 21/0.15 & 6.4/0.70 & 2.9 & 16/6.9 & 17/6.3 & 7.7 & 7.3 & 43/14 & 31/23 \\
%DIFDELCMD < Floor: batches of 100 & 5.3 & 24/0.15 & 8.2/0.80 & 3.5 & 19/8.5 & 21/7.9 & 8.9 & 8.9 & 57/17 & 46/29 \\
%DIFDELCMD < Floor: 90\,\% quantile & 6.2 & 31/0.20 & 8.8/0.90 & 4.0 & 21/10 & 23/8.6 & 10 & 10 & 61/20 & 45/34 \\
%DIFDELCMD < Floor: 99\,\% quantile & 3.4 & 14/0.10 & 5.2/0.50 & 2.4 & 11/4.4 & 13/4.4 & 5.7 & 6.1 & 37/11 & 29/15 \\
%DIFDELCMD < Calibration: 50 parts & 5.0 & 22/0.20 & 6.8/0.70 & 3.4 & 18/7.5 & 19/6.2 & 8.4 & 7.7 & 50/16 & 37/28 \\
%DIFDELCMD < Calibration: 100 parts & 4.9 & 22/0.40 & 6.8/0.40 & 3.3 & 18/7.5 & 19/6.2 & 8.7 & 7.8 & 50/16 & 38/28 \\
%DIFDELCMD < Calibration: 250 parts & 4.9 & 22/0.45 & 7.0/0.45 & 3.3 & 18/7.5 & 19/6.2 & 8.7 & 8.3 & 50/16 & 38/28 \\
%DIFDELCMD < Conformance $p=0.90$ & 4.9 & 23/0.20 & 7.2/0.70 & 3.3 & 18/7.5 & 18/6.4 & 8.4 & 8.1 & 51/15 & 38/28 \\
%DIFDELCMD < Conformance $p=0.99$ & 5.1 & 23/0 & 7.5/0.70 & 3.3 & 18/8.3 & 20/6.9 & 8.7 & 8.4 & 51/16 & 38/28 \\
%DIFDELCMD < Interval level 0.80 &  & 23/0.15 & 6.4/1.0 &  & 18/8.1 & 17/8.1 &  &  & 50/17 & 38/28 \\
%DIFDELCMD < Interval level 0.95 &  & 23/0.15 & 8.0/0.25 &  & 18/8.1 & 20/6.2 &  &  & 51/16 & 39/27 \\
%DIFDELCMD < Response surface & 4.8 & 25/0.30 & 7.2/0.70 & 3.4 & 18/7.6 & 20/6.9 & 8.5 & 8.2 & 49/25 & 37/27 \\
%DIFDELCMD < GP: sampling bound &  &  & 6.7/1.4 &  &  & 18/7.4 &  &  &  & 38/30 \\
%DIFDELCMD < Unit: $q_{95}$ floor & 5.3 & 17/0.15 & 6.9/0.40 & 2.8 & 15/6.0 & 18/6.2 & 6.9 & 7.1 & 41/13 & 34/30 \\
%DIFDELCMD < \midrule Refit, $\delta_\mathrm{d}$ & 4.9--7.3 & 22--23 & 4.0--11 & 2.3--4.6 & 16--18 & 18--20 & 6.9--9.2 & 7.3--10 & 49--53 & 34--39 \\
%DIFDELCMD < Refit, $\delta_\mathrm{s}$ & 4.9--7.3 & 0.15--1.4 & 0.75--3.0 & 2.3--4.6 & 6.8--8.4 & 6.2--7.8 & 6.9--9.2 & 7.3--10 & 15--22 & 27--33 \\
%DIFDELCMD < \botrule
%DIFDELCMD < \end{tabular}
%DIFDELCMD < %%%
\DIFdelendFL \DIFaddbeginFL \setlength{\tabcolsep}{3pt}
\caption{The further rules and their decisive/safe distances (floors) by relation}\label{tab:further}
\begin{tabular}{@{}lL{5.0cm}rrrr@{}}
\toprule
Rule & Prediction of $q_p$ & Sibling & Interp. & Extrap. & New family \\
\midrule
M1 & M0 plus the median offset of the calibration alternatives & 15 & 11 & 25 & 33 \\
M1n$^\dagger$ & median $q_p$ of the calibration alternatives & 9.1 & 4.8 & 12 & 131 \\
M2n$^\dagger$ & M1n with the largest-residual margin & 19/$<$1 & 13/$<$1 & 16/9.1 & 133/128 \\
M5n$^\dagger$ & GP of $q_p$ over the descriptors, ML-II, noise-variance bound & 7.2/$<$1 & 16/$<$1 & 18/8.6 & 133/128 \\
M7$^\ddagger$ & M7n with the nominal simulation in the mean (AR(1) multi-fidelity) & 7.1/$<$1 & 22/16 & 28/17 & 140/104 \\
Mc$^\dagger$ & friction per pattern fitted to the draw-ins, least squares, plus offset & 27 & 14 & 40 & 43 \\
Mcd$^\ddagger$ & friction calibrated with a GP discrepancy; posterior mixture & 7.2/$<$1 & 22/1.6 & 24/5.8 & 30/19 \\
M0w$^\dagger$ & range of the simulation over thickness and friction & 119/94 & 119/94 & 119/77 & 117/79 \\
M3 & matched simulation plus boosted discrepancy, parts; jackknife+ & 23/$<$1 & 29/13 & 33/17 & 40/17 \\
M3n & boosted model of the characteristic, parts; jackknife+ & 11/$<$1 & 15/7.0 & 13/7.8 & 135/126 \\
M6$^\dagger$ & M5 mean with a conformal margin on standardised residuals & 17/$<$1 & 47/$<$1 & 56/8.1 & 35/28 \\
\botrule
\end{tabular}
\DIFaddendFL \footnotetext{\DIFdelbeginFL \DIFdelFL{Rules refitted for every variant; empty cells: }\DIFdelendFL \DIFaddbeginFL \DIFaddFL{Pairs with and without }\DIFaddendFL the \DIFdelbeginFL \DIFdelFL{variant does not affect }\DIFdelendFL \DIFaddbeginFL \DIFaddFL{simulation: M1/M1n, M2/M2n, M3/M3n, M5/M5n and M7/M7n, as M1s/M1sn (Table~\ref{tab:rules}); M3n is M4 in }\DIFaddendFL the \DIFdelbeginFL \DIFdelFL{rule}\DIFdelendFL \DIFaddbeginFL \DIFaddFL{code}\DIFaddendFL ; \DIFaddbeginFL \DIFaddFL{M6 is }\DIFaddendFL a \DIFaddbeginFL \DIFaddFL{robustness variant. A }\DIFaddendFL new family in \DIFdelbeginFL \DIFdelFL{its own floor (in }\DIFdelendFL the produced family's floor\DIFdelbeginFL \DIFdelFL{in Table~\ref{tab:distances})}\DIFdelendFL . \DIFdelbeginFL \DIFdelFL{The floor between series contains the lubrication effect }\DIFdelendFL \DIFaddbeginFL \DIFaddFL{Pooled scopes }\DIFaddendFL and \DIFdelbeginFL \DIFdelFL{, for a new variant, the series that NN averages, which makes it close to circular there}\DIFdelendFL \DIFaddbeginFL \DIFaddFL{intervals in \esm{Table~S1}}\DIFaddendFL .\DIFdelbeginFL \DIFdelFL{Warm-up}\DIFdelendFL \DIFaddbeginFL }
\end{table}
%DIF > % END GENERATED further

%DIF > % BEGIN GENERATED gp
\begin{table}[p]
\tabsizesmall
\setlength{\tabcolsep}{2.6pt}
\caption{Fitted Gaussian processes by relation: hyperparameters and where their centres lie}\label{tab:gp}
\begin{tabular}{@{}llrrrrrrrrr@{}}
\toprule
Relation & Rule & Cases & $\ell_F$ & At bound & Noise & $\Delta$mean & ${\le}0.01$ & $\Delta$NN & $\Delta$M1sn & HW \\
\midrule
Sibling & M5 & 126 & 0.13 & 45 & 1.00 & 2.6 & 0 & 0.27 & 0.92 & 2.2 \\
 & M5n & 126 & 1.76 & 22 & 1.00 & 2.1 & 20 & 0.28 & 0.88 & 2.3 \\
 & M7n & 126 & 1.43 & -- & 3.6 & 2.6 & 0 & 0.40 & 1.0 & 1.6 \\
Interpolated & M5 & 14 & 0.10 & 64 & 1.00 & 4.1 & 0 & 4.1 & 4.1 & 6.4 \\
 & M5n & 14 & 0.10 & 79 & 1.00 & 0 & 95 & 0 & 0 & 5.8 \\
 & M7n & 14 & 1.92 & -- & 4.0 & 0.36 & 2 & 0.36 & 0.36 & 2.0 \\
Extrapolated & M5 & 28 & 0.10 & 54 & 1.00 & 6.8 & 0 & 6.1 & 4.9 & 4.6 \\
 & M5n & 28 & 0.10 & 68 & 1.00 & 0 & 93 & 1.3 & 3.9 & 3.7 \\
 & M7n & 28 & 1.95 & -- & 3.1 & 4.2 & 0 & 2.9 & 0.33 & 2.8 \\
New family & M5 & 14 & 0.21 & 43 & 1.00 & 24 & 0 & 21 & 22 & 1.9 \\
 & M5n & 14 & 1.77 & 21 & 1.00 & 1.9 & 13 & 0.22 & 0.66 & 2.0 \\
 & M7n & 14 & 1.51 & -- & 3.7 & 2.2 & 1 & 0.37 & 0.83 & 1.2 \\
\botrule
\end{tabular}
\footnotetext{\DIFaddFL{$\ell_F$}\DIFaddendFL : \DIFdelbeginFL \DIFdelFL{first 150 parts}\DIFdelendFL \DIFaddbeginFL \DIFaddFL{median force length scale (100\,kN)}\DIFaddendFL . \DIFdelbeginFL \DIFdelFL{Calibration}\DIFdelendFL \DIFaddbeginFL \DIFaddFL{At bound}\DIFaddendFL : \DIFdelbeginFL \DIFdelFL{calibration $q_{95}$ from the first $n$ parts, truths and floors from the full series}\DIFdelendFL \DIFaddbeginFL \DIFaddFL{fits with $\ell_F$ at its lower bound (\%)}\DIFaddendFL . \DIFdelbeginFL \DIFdelFL{Response surface}\DIFdelendFL \DIFaddbeginFL \DIFaddFL{Noise}\DIFaddendFL : \DIFdelbeginFL \DIFdelFL{simulations from a quadratic fit }\DIFdelendFL \DIFaddbeginFL \DIFaddFL{fitted noise }\DIFaddendFL over \DIFdelbeginFL \DIFdelFL{thickness and friction; other GP kernels and descriptors move }\DIFdelendFL \DIFaddbeginFL \DIFaddFL{its bound (}\DIFaddendFL M5\DIFdelbeginFL \DIFdelFL{by at most 0.4 floors. Unit: the $q_{95}$ floor}\DIFdelendFL , \DIFdelbeginFL \DIFdelFL{95\,\% quantile of the difference between }\DIFdelendFL \DIFaddbeginFL \DIFaddFL{M5n) or nugget over }\DIFaddendFL the \DIFdelbeginFL \DIFdelFL{$q_{95}$ of two batches of 100 parts. Refit: 95\,\% intervals }\DIFdelendFL \DIFaddbeginFL \DIFaddFL{sampling error }\DIFaddendFL of the \DIFdelbeginFL \DIFdelFL{decisive and safe distances over 200 draws of each geometry's lubrication patterns with replacement }\DIFdelendFL \DIFaddbeginFL \DIFaddFL{truths }\DIFaddendFL (\DIFdelbeginFL \DIFdelFL{Section~\ref{sec:evaldesign}}\DIFdelendFL \DIFaddbeginFL \DIFaddFL{M7n}\DIFaddendFL )\DIFaddbeginFL \DIFaddFL{, median}\DIFaddendFL . \DIFdelbeginFL %DIFDELCMD < \MBLOCKRIGHTBRACE
%DIFDELCMD < \end{table}
%DIFDELCMD < %%%
%DIF < % END GENERATED robust
%DIFDELCMD < 

%DIFDELCMD < %%%
%DIF < % BEGIN GENERATED intervals
%DIFDELCMD < \begin{table}[p]
%DIFDELCMD < \tabsizesmall
%DIFDELCMD < \setlength{\tabcolsep}{2.6pt}
%DIFDELCMD < \caption{Interval rules by relation: coverage, half-width and the interval centre as a point rule}\label{tab:intervals}
%DIFDELCMD < \begin{tabular*}{\textwidth}{@{\extracolsep{\fill}}lrrrrrrrrrrrr@{}}
%DIFDELCMD < \toprule
%DIFDELCMD <  & \multicolumn{3}{c}{Sibling produced} & \multicolumn{3}{c}{Interpolated setting} & \multicolumn{3}{c}{Extrapolated setting} & \multicolumn{3}{c}{New family} \\
%DIFDELCMD < \cmidrule(lr){2-4}\cmidrule(lr){5-7}\cmidrule(lr){8-10}\cmidrule(lr){11-13}
%DIFDELCMD < Rule & Cov & HW & Ctr & Cov & HW & Ctr & Cov & HW & Ctr & Cov & HW & Ctr \\
%DIFDELCMD < \midrule
%DIFDELCMD < M2 & 89 & 7.6 & 8.6 & 71 & 5.7 & 10 & 44 & 4.8 & 13 & 47 & 7.6 & 34 \\
%DIFDELCMD < M2n & 90 & 5.9 & 9.1 & 93 & 5.0 & 4.8 & 33 & 3.1 & 12 & 12 & 5.7 & 131 \\
%DIFDELCMD < M3 & 96 & 9.3 & 9.9 & 43 & 7.0 & 19 & 43 & 8.9 & 21 & 40 & 9.7 & 29 \\
%DIFDELCMD < M3n & 89 & 3.0 & 3.7 & 36 & 3.4 & 11 & 48 & 2.3 & 9.9 & 7 & 3.3 & 130 \\
%DIFDELCMD < M5 & 86 & 2.2 & 3.3 & 71 & 6.4 & 9.8 & 46 & 4.6 & 13 & 12 & 2.3 & 34 \\
%DIFDELCMD < M5n & 85 & 2.3 & 3.1 & 93 & 5.8 & 4.7 & 36 & 3.7 & 12 & 0 & 2.2 & 130 \\
%DIFDELCMD < \botrule
%DIFDELCMD < \end{tabular*}
%DIFDELCMD < \footnotetext{%%%
\DIFdelFL{Cov}\DIFdelendFL \DIFaddbeginFL \DIFaddFL{$\Delta$mean, $\Delta$NN, $\Delta$M1sn}\DIFaddendFL : \DIFdelbeginFL \DIFdelFL{share }\DIFdelendFL \DIFaddbeginFL \DIFaddFL{median distance (floors) }\DIFaddendFL of \DIFaddbeginFL \DIFaddFL{the centre from the calibration mean, NN and M1sn; ${\le}0.01$: }\DIFaddendFL cases \DIFdelbeginFL \DIFdelFL{whose true $q_{95}$ lies in }\DIFdelendFL \DIFaddbeginFL \DIFaddFL{within 0.01 floor of }\DIFaddendFL the \DIFdelbeginFL \DIFdelFL{interval }\DIFdelendFL \DIFaddbeginFL \DIFaddFL{calibration mean }\DIFaddendFL (\%\DIFdelbeginFL \DIFdelFL{; nominal 90}\DIFdelendFL ). HW: median half-width (floors).\DIFdelbeginFL \DIFdelFL{Ctr}\DIFdelendFL \DIFaddbeginFL }
\end{table}
%DIF > % END GENERATED gp

%DIF > % BEGIN GENERATED crossing
\begin{table}[p]
\tabsizesmall
\setlength{\tabcolsep}{2.6pt}
\caption{Point predictors with a relation-matched interval, against the rules' own intervals}\label{tab:crossing}
\begin{tabular}{@{}lrrrrrrrr@{}}
\toprule
 & \multicolumn{2}{c}{Sibling} & \multicolumn{2}{c}{Interpolated} & \multicolumn{2}{c}{Extrapolated} & \multicolumn{2}{c}{New family} \\
\cmidrule(lr){2-3}\cmidrule(lr){4-5}\cmidrule(lr){6-7}\cmidrule(lr){8-9}
Predictor + interval & $\delta_\mathrm{d}/\delta_\mathrm{s}$ & Cov & $\delta_\mathrm{d}/\delta_\mathrm{s}$ & Cov & $\delta_\mathrm{d}/\delta_\mathrm{s}$ & Cov & $\delta_\mathrm{d}/\delta_\mathrm{s}$ & Cov \\
\midrule
NN+RM & 6.3/$<$1 & 84 & 9.6/$<$1 & 88 & 20/$<$1 & 93 & 293/4.5 & 86 \\
M1sn+RM & 7.1/$<$1 & 88 & 9.6/$<$1 & 88 & 18/$<$1 & 88 & 293/4.3 & 86 \\
M1s+RM & 11/$<$1 & 90 & 43/6.9 & 86 & 31/4.8 & 86 & 243/$<$1 & 93 \\
M5n+RM & 6.1/$<$1 & 88 & 9.6/$<$1 & 88 & 24/$<$1 & 89 & 293/4.5 & 86 \\
M5+RM & 6.0/$<$1 & 86 & 21/$<$1 & 93 & 26/2.2 & 85 & 52/$<$1 & 91 \\
M3n+RM & 8.1/$<$1 & 88 & 21/$<$1 & 86 & 20/$<$1 & 93 & 292/4.8 & 86 \\
M7n+RM & 6.5/$<$1 & 91 & 11/$<$1 & 90 & 18/$<$1 & 89 & 293/5.1 & 86 \\
M7+RM & 6.7/$<$1 & 89 & 45/6.4 & 86 & 31/$<$1 & 89 & 248/$<$1 & 93 \\
\midrule \multicolumn{9}{@{}l}{\emph{The rules with their own intervals}} \\
M2 & 23/$<$1 & 89 & 18/2.7 & 71 & 18/9.9 & 44 & 41/26 & 47 \\
M3 & 23/$<$1 & 96 & 29/13 & 43 & 33/17 & 43 & 40/17 & 40 \\
M5 & 7.2/$<$1 & 86 & 19/2.8 & 71 & 19/8.9 & 46 & 35/31 & 12 \\
M7n & 6.8/$<$1 & 85 & 7.4/3.4 & 52 & 15/5.9 & 56 & 132/130 & 0 \\
\botrule
\end{tabular}
\footnotetext{\DIFaddFL{P+RM}\DIFaddendFL : \DIFdelbeginFL \DIFdelFL{decisive distance }\DIFdelendFL \DIFaddbeginFL \DIFaddFL{the prediction }\DIFaddendFL of \DIFaddbeginFL \DIFaddFL{P (for interval rules }\DIFaddendFL the interval centre\DIFdelbeginFL \DIFdelFL{used as a point rule (floors}\DIFdelendFL ) \DIFdelbeginFL \DIFdelFL{. A new family in the floor of }\DIFdelendFL \DIFaddbeginFL \DIFaddFL{with }\DIFaddendFL the \DIFdelbeginFL \DIFdelFL{produced family. The GP noise variance sits at its lower bound in 97 and 99}\DIFdelendFL \DIFaddbeginFL \DIFaddFL{90}\DIFaddendFL \,\% \DIFaddbeginFL \DIFaddFL{split-conformal quantile }\DIFaddendFL of \DIFdelbeginFL \DIFdelFL{the M5 and M5n fits with a sibling and in every new-setting and new-family fit.}%DIFDELCMD < \MBLOCKRIGHTBRACE
%DIFDELCMD < \end{table}
%DIFDELCMD < %%%
%DIF < % END GENERATED intervals
%DIFDELCMD < 

%DIFDELCMD < %%%
%DIF < % BEGIN GENERATED step6
%DIFDELCMD < \begin{table}[p]
%DIFDELCMD < \tabsizesmall
%DIFDELCMD < \setlength{\tabcolsep}{2.2pt}
%DIFDELCMD < \caption{The procedure's decision rule (step 6) by relation: guard band, distances and trial shares}\label{tab:step6}
%DIFDELCMD < \begin{tabular*}{\textwidth}{@{\extracolsep{\fill}}lrrrrrrrrrrrr@{}}
%DIFDELCMD < \toprule
%DIFDELCMD <  & \multicolumn{4}{c}{Sibling produced} & \multicolumn{4}{c}{Interpolated setting} & \multicolumn{4}{c}{Extrapolated setting} \\
%DIFDELCMD < \cmidrule(lr){2-5}\cmidrule(lr){6-9}\cmidrule(lr){10-13}
%DIFDELCMD < Rule & $g$ & $\delta_\mathrm{d}/\delta_\mathrm{s}$ & $T_{10}$ & $T_{20}$ & $g$ & $\delta_\mathrm{d}/\delta_\mathrm{s}$ & $T_{10}$ & $T_{20}$ & $g$ & $\delta_\mathrm{d}/\delta_\mathrm{s}$ & $T_{10}$ & $T_{20}$ \\
%DIFDELCMD < \midrule
%DIFDELCMD < M1 & none & -- & 100 & 100 & 8.75 & 20/2.1 & 44 & 4 & none & -- & 100 & 100 \\
%DIFDELCMD < M1s & 2.0 & 6.9/2.9 & 1 & 0 & none & -- & 100 & 100 & none & -- & 100 & 100 \\
%DIFDELCMD < M2 & 0 & 23/0.15 & 33 & 8 & 0.75 & 19/2.0 & 31 & 4 & none & -- & 100 & 100 \\
%DIFDELCMD < M5 & 0 & 7.2/0.75 & 1 & 0 & 0.5 & 20/2.3 & 40 & 4 & none & -- & 100 & 100 \\
%DIFDELCMD < M5n & 0 & 7.2/0.55 & 0 & 0 & 0 & 16/0 & 25 & 0 & 6.75 & 25/1.9 & 48 & 15 \\
%DIFDELCMD < NN & 0.25 & 3.6/3.1 & 0 & 0 & 2.0 & 6.7/2.7 & 0 & 0 & 7.0 & 17/2.5 & 26 & 1 \\
%DIFDELCMD < M3 & 0 & 23/0 & 45 & 6 & 17.75 & 47/0 & 94 & 69 & 16.75 & 51/0.05 & 85 & 69 \\
%DIFDELCMD < M3n & 0 & 11/0.05 & 9 & 0 & 6.0 & 21/1.0 & 55 & 6 & 6.5 & 19/1.3 & 47 & 3 \\
%DIFDELCMD < M1sn & 0.75 & 4.3/2.8 & 0 & 0 & 2.0 & 6.7/2.7 & 0 & 0 & 6.0 & 15/2.9 & 21 & 0 \\
%DIFDELCMD < \botrule
%DIFDELCMD < \end{tabular*}
%DIFDELCMD < \footnotetext{%%%
\DIFdelFL{$g$: guard band }\DIFdelendFL \DIFaddbeginFL \DIFaddFL{P's absolute errors }\DIFaddendFL (floors) \DIFdelbeginFL \DIFdelFL{under }\DIFdelendFL \DIFaddbeginFL \DIFaddFL{over }\DIFaddendFL the \DIFdelbeginFL \DIFdelFL{reference prior (Section~\ref{sec:procedure}); no applicability guard is applied; $\delta_\mathrm{d}/\delta_\mathrm{s}$: distances }\DIFdelendFL \DIFaddbeginFL \DIFaddFL{cases }\DIFaddendFL of the \DIFdelbeginFL \DIFdelFL{rule with its interval widened by $g$; $T_{10}$}\DIFdelendFL \DIFaddbeginFL \DIFaddFL{same relation in the other family (sibling 63}\DIFaddendFL , \DIFdelbeginFL \DIFdelFL{$T_{20}$}\DIFdelendFL \DIFaddbeginFL \DIFaddFL{interpolated 21, extrapolated 42, new family 63). Cov}\DIFaddendFL : \DIFdelbeginFL \DIFdelFL{share }\DIFdelendFL \DIFaddbeginFL \DIFaddFL{coverage }\DIFaddendFL of \DIFdelbeginFL \DIFdelFL{verdicts sent to trial }\DIFdelendFL \DIFaddbeginFL \DIFaddFL{the true $q_p$ }\DIFaddendFL (\%)\DIFdelbeginFL \DIFdelFL{at requirements 10 and 20 floors from the truth. None: no band up to 20 floors reaches 95\,\%, so the procedure sends every design to trial}\DIFdelendFL ; \DIFdelbeginFL \DIFdelFL{so also M0 }\DIFdelendFL \DIFaddbeginFL \DIFaddFL{cluster intervals }\DIFaddendFL and \DIFdelbeginFL \DIFdelFL{M0w }\DIFdelendFL \DIFaddbeginFL \DIFaddFL{interval scores }\DIFaddendFL in \DIFdelbeginFL \DIFdelFL{every relation, Mc with a sibling and for an extrapolated setting (for an interpolated one $g=18$), and every rule for a }\DIFdelendFL \DIFaddbeginFL \DIFaddFL{\esm{Table~S4}. A }\DIFaddendFL new family \DIFdelbeginFL \DIFdelFL{, scored }\DIFdelendFL in the produced family's floor.}
\end{table}
%DIF < % END GENERATED step6
%DIF > % END GENERATED crossing
\DIFaddbegin 

%DIF > % BEGIN GENERATED groups
\begin{table}[p]
\tabsizesmall
\setlength{\tabcolsep}{1pt}
\caption{Decisive distance per characteristic in floors and, in brackets, in its unit}\label{tab:groups}
\begin{tabular*}{\textwidth}{@{\extracolsep{\fill}}lrrrrrrr@{}}
\toprule
 & \multicolumn{3}{c}{NN} & M1sn & M1s & M5 & \\
\cmidrule(lr){2-4}
Characteristic & Sibling & Interp. & Extrap. & Extrap. & Extrap. & New family & $P_\mathrm{pk}$ \\
\midrule
Mid-side (mm) & 2.0 (0.099) & 3.8 (0.19) & 6.5 (0.32) & 5.7 (0.28) & 4.0 (0.20) & 5.5 (0.27) & 1.3 \\
Corner (mm) & 4.5 (0.13) & 1.9 (0.058) & 6.7 (0.23) & 2.8 (0.083) & 2.9 (0.100) & 9.9 (0.31) & 2.5 \\
Wavin. (mm) & $<$1 ($<$0.0028) & 4.8 (0.014) & 11 (0.040) & 9.3 (0.027) & 21 (0.077) & 22 (0.047) & 0.9 \\
Wall ($^\circ$) & 3.4 (0.26) & $<$1 ($<$0.078) & 1.9 (0.15) & 1.8 (0.14) & 2.0 (0.16) & 7.5 (0.58) & 1.7 \\
Arm ($^\circ$) & 2.8 (0.28) & 5.9 (0.59) & 5.4 (0.53) & 9.5 (0.94) & 5.9 (0.69) & 39 (3.9) & 1.6 \\
Depth (mm) & 3.3 (0.039) & 4.4 (0.038) & 15 (0.13) & 9.9 (0.086) & 165 (1.4) & 20 (0.20) & 2.4 \\
Dome (mm) & 7.0 (0.038) & 5.3 (0.044) & 4.3 (0.020) & 3.0 (0.020) & 15 (0.056) & 40 (0.21) & 4.0 \\
\botrule
\end{tabular*}
\footnotetext{\DIFaddFL{18 cases per characteristic with a sibling and for a new family, 6 interpolated, 12 extrapolated; units through the pooled floor of the characteristic; a new family in the produced family's floor. $P_\mathrm{pk}$: median performance index of a requirement at the decisive distance of NN with a sibling. All rules and the safe distances in \esm{Table~S15}.}}
\end{table}
%DIF > % END GENERATED groups

%DIF > % BEGIN GENERATED validation
\begin{table}[p]
\tabsizesmall
\setlength{\tabcolsep}{1pt}
\caption{The simulation dataset against the measurement, over the six geometry--force cells}\label{tab:validation}
\begin{tabular*}{\textwidth}{@{\extracolsep{\fill}}lrrrrrrr@{}}
\toprule
Characteristic & Measured & Sim.\,$-$\,meas. & In floors & Inside & Closest $\mu$ & $x$--$y$ & Roughness \\
\midrule
Mid-side (mm) & 14.73--16.08 & 0.12--2.01 & 2--38 & 2 & 0.11--0.15 & 0.1--0.2 & 0.1 \\
Corner (mm) & 4.62--10.00 & 2.59--3.48 & 49--116 & 0 & 0.15 & -- & 0.1--0.2 \\
Wavin. (mm) & 0.052--0.103 & $-$0.052 to +0.011 & $-$14 to +5 & 3 & 0.05--0.14 & -- & 0.4--4.0 \\
Wall ($^\circ$) & 11.43--21.60 & $-$2.14 to $-$0.37 & $-$27 to $-$5 & 0 & 0.05--0.15 & 1.8--2.3 & 1.7--3.1 \\
Arm ($^\circ$) & $-$4.95 to +0.76 & $-$4.64 to +0.42 & $-$40 to +4 & 3 & 0.05--0.11 & 0.5--1.1 & 0.9 \\
Depth (mm) & 28.29--28.82 & 0.71--1.06 & 58--116 & 0 & 0.05--0.15 & -- & 0.7--0.8 \\
Dome (mm) & 0.030--0.116 & $-$0.087 to +0.104 & $-$11 to +30 & 3 & 0.08--0.15 & -- & 2.4--9.8 \\
\botrule
\end{tabular*}
\footnotetext{\DIFaddFL{Measured: range of the cells' centres (signed). Sim.\,$-$\,meas.: simulation at 0.99\,mm and friction 0.10 minus the measured centre, in the unit and in floors. Inside: cells whose measured centre lies within the simulated range over the friction grid (of 6). Closest $\mu$: grid friction nearest the measurement. $x$--$y$, Roughness: standard deviation of the difference between the two directions of the quarter model and roughness along the friction grid (floors), the extraction noise of the simulation. Per cell in \esm{Table~S18}.}}
\end{table}
%DIF > % END GENERATED validation
\DIFaddend 

\begin{table}[p]
%DIF <  float page with Tables A3 and A4
\tabsizesmall
\setlength{\tabcolsep}{3pt}
\caption{Where the numbers quoted in the text come from: scripts and result files (\texttt{results/*.csv})}\label{tab:mapping}
\DIFdelbeginFL %DIFDELCMD < \begin{tabular}{@{}L{0.95cm}L{4.2cm}L{7.5cm}@{}}
%DIFDELCMD < \toprule
%DIFDELCMD < Section & Results & Script: result files \\
%DIFDELCMD < \midrule
%DIFDELCMD < 3.4, 4.1, 4.3 & failing-side weight, run metadata and oil film, scan-line scatter, measurement resolution & panel.py: panel\_failing\_weight, panel\_runs, panel\_resolution\_steps; measurement.py: meas\_scan \\
%DIFDELCMD < 5.1 & time order, lag profile, normal model, subgroups, part-to-part share, between series, $q_{95}$ reproducibility, in-line drift & alternatives.py: alt\_floor; sensitivity.py: sens\_drift; panel.py: panel\_variogram, panel\_floor\_protocol; measurement.py: meas\_floor; robustness.py: rob\_floor\_between; inline.py: inline\_drift \\
%DIFDELCMD < 5.2 & resolvable contrasts and changes, simulation offsets, force effect, arm sign, cup depth & alternatives.py: alt\_effects, alt\_mrc; robustness.py: rob\_bhf\_pairs; decisions.py: dec\_alternatives; panel.py: panel\_force\_effect, panel\_arm\_sign, panel\_m0\_split \\
%DIFDELCMD < 5.3--5.4 & distances in both floors, paired differences, strata, force levels, decomposition, pooled coverage & decisions.py: dec\_resolution, dec\_paired, dec\_coverage; panel.py: panel\_source\_floor(\_paired), panel\_sibling\_strata, panel\_force\_levels, panel\_decomposition \\
%DIFDELCMD < 5.5 & simulation pairs, drawing nominal, surrogate, coverage, GP fits & decisions.py: dec\_paired, dec\_coverage, dec\_gp; panel.py: panel\_nominal\_offset, panel\_source\_floor(\_paired); design\_space.py: ds\_surrogate(\_rule) \\
%DIFDELCMD < 5.6, 5.9 & budget, short series, robustness variants, refit, $q_{95}$ unit & budget.py: budget\_design, budget\_resolution; robustness.py: rob\_decisions, rob\_refit(\_paired); panel.py: panel\_q95\_unit \\
%DIFDELCMD < 3.6, 5.7 & margins, guard bands and priors, uncapped bands, guards, worked example & panel.py: panel\_step6(\_curves), panel\_guard\_priors; guard.py: dec\_guard; decisions.py: dec\_intervals \\
%DIFDELCMD < 5.8 & capability-anchored limits, ISO 2768-1 scenario, costs & panel.py: panel\_capability; scenarios.py: scen\_strata, scen\_scores, scen\_cost\_map \\
%DIFDELCMD < 6.5 & PA12 attempt & case2\_pa12.py: case2\_floor, case2\_effects \\
%DIFDELCMD < \botrule
%DIFDELCMD < \end{tabular}
%DIFDELCMD < %%%
\DIFdelendFL \DIFaddbeginFL \begin{tabular}{@{}L{0.95cm}L{4.2cm}L{7.5cm}@{}}
\toprule
Section & Results & Script: result files \\
\midrule
3.4, 4.1, 4.3 & failing-side weight, run metadata, scan redundancy and biases, resolution & panel.py: panel\_failing\_weight, panel\_runs, panel\_resolution\_steps; measurement.py: meas\_scan; simcheck.py: meas\_characteristics, meas\_cut\_walls \\
4.2 & simulation noise of the characteristics & measurement.py: meas\_sim \\
5.1 & time order, lag profile, single-run floors, warm-up, components, tooling reference, in-line drift & alternatives.py: alt\_floor; sensitivity.py: sens\_drift; panel.py: panel\_variogram, panel\_floor\_protocol; measurement.py: meas\_floor; simcheck.py: meas\_run\_floors, meas\_tooling\_reference; inline.py: inline\_drift \\
5.2 & resolvable contrasts, validation, friction fits, oil mapping, depth offset, envelope, force effect, arm sign, surrogate & alternatives.py: alt\_effects, alt\_mrc; simcheck.py: sim\_validation, sim\_mc\_friction, sim\_oil\_mapping, sim\_depth\_offset, sim\_envelope\_variant; learners.py: learn\_mcd\_posterior; panel.py: panel\_force\_effect, panel\_arm\_sign, panel\_m0\_split; design\_space.py: ds\_surrogate\_rule \\
5.3 & distances, paired differences, strata, per characteristic, units & decisions.py: dec\_resolution, dec\_paired; learners.py: learn\_resolution, learn\_source\_floor, learn\_paired; panel.py: panel\_source\_floor; qualify.py: qual\_strata, qual\_per\_characteristic \\
5.4 & GP diagnostics, crossing, coverage, interval score, M6, tuning study, few-shot transfer & learners.py: learn\_gp\_diagnostics, learn\_fewshot; qualify.py: qual\_coverage, qual\_m6\_coverage; decisions.py: dec\_gp; tuning.py: tune\_decisions \\
5.5 & guard bands per family, out of sample, priors, per characteristic & qualify.py: qual\_guard\_family, qual\_guard\_oos, qual\_guard\_prior\_width, qual\_guard\_characteristic; panel.py: panel\_step6, panel\_guard\_priors \\
5.6 & costs, trial, capability-anchored limits & qualify.py: qual\_costs, qual\_trial; panel.py: panel\_capability \\
5.7 & between-run sensitivity, cross-fit, resolvable difference, design level, variants, floors without the held-out series, robustness, budget & qualify.py: qual\_between\_run, qual\_crossfit, qual\_resolvable\_difference, qual\_design\_level, qual\_variants, qual\_decomposition, qual\_floor\_loao; robustness.py: rob\_decisions, rob\_refit; budget.py: budget\_design \\
6.4 & scoring function; PA12 attempt & opchar.py (self-check against dec\_resolution); case2\_pa12.py: case2\_floor \\
\botrule
\end{tabular}
\footnotetext{\DIFaddFL{M3n is named M4 in the released code.}}
\DIFaddendFL \end{table}

\emph{Computation details.} Batches need at least 80\,\% valid values. M5 and M5n are Gaussian-process regressions (constant times anisotropic squared-exponential kernel plus white noise, normalised targets, hyperparameters by maximum marginal likelihood without hyperpriors, three restarts, length scales in $[0.1, 100]$) on a geometry indicator, the force in 100\,kN and the lubrication rank; the lower bound of the noise variance \DIFdelbegin \DIFdel{, }\DIFdelend \DIFaddbegin \DIFadd{is }\DIFaddend the larger of the mean sampling variance of the calibration alternatives' \DIFdelbegin \DIFdel{$q_{95}$ }\DIFdelend \DIFaddbegin \DIFadd{$q_p$ }\DIFaddend and the pooled variance between their series of one geometry and force\DIFdelbegin \DIFdel{, is expressed on the normalised scale}\DIFdelend \DIFaddbegin \DIFadd{. M7n and M7 work in pooled floors: the mean basis is an intercept, the force (or, for M7, the nominal simulation) and the lubrication rank, plus the geometry indicator when both families calibrate; the coefficients have a flat prior and are integrated out (restricted likelihood); the hyperparameters are maximum a posteriori under log-normal priors with medians of 200\,kN and one pattern step for the length scales, two floors for the GP and one floor for the nugget standard deviation, and log-scale standard deviations of 1, 1.5 and 1; the predictive variance adds the nugget and the mean sampling variance of the calibration $q_p$. Mcd has a uniform prior over the 11 friction values of the grid, a GP discrepancy with a constant mean and the same priors, a posterior from the two draw-ins, and a 90\,\% interval from the posterior mixture}\DIFaddend . NN takes the mean \DIFdelbegin \DIFdel{$q_{95}$ }\DIFdelend \DIFaddbegin \DIFadd{$q_p$ }\DIFaddend at the nearest force of the family, both neighbours when equally near; M1sn fits \DIFdelbegin \DIFdel{$q_{95}=a+b\,$}\DIFdelend \DIFaddbegin \DIFadd{$q_p=a+b\,$}\DIFaddend force by least squares. Mc calibrates, per lubrication pattern, the friction coefficient whose simulated draw-ins are closest to the calibration centres\DIFaddbegin \DIFadd{, without a discrepancy}\DIFaddend . M3 and M3n use histogram gradient boosting (depth 3, 200 iterations, rate 0.08) on geometry, force, lubrication, sheet thickness and oil film, with 500 incoming conditions drawn from calibration parts of the same pattern\DIFdelbegin \DIFdel{; their jackknife}\DIFdelend \DIFaddbegin \DIFadd{, and jackknife}\DIFaddend + \DIFdelbegin \DIFdel{interval takes the leave-one-out predictions for the new alternative, shifted by the median leave-one-out residual, minus and plus the absolute residuals about that median, with its bounds at the extreme order statistics when $n$ is too small for the 0.9 quantiles}\DIFdelend \DIFaddbegin \DIFadd{intervals}\DIFaddend . Guard bands lie on a grid of 0.25 floors. \DIFdelbegin \DIFdel{The refitting bootstrap draws the three lubrication patterns of each geometry with replacement 200 times, at least two of them distinct, and does not restrict the noise floor to distinct series. }\DIFdelend Software: Python 3.11 with NumPy, pandas, SciPy, scikit-learn \citep{pedregosa2011scikit} and Matplotlib.

\DIFaddbegin \FloatBarrier
\DIFaddend \end{appendices}

\backmatter

\DIFaddbegin \bmhead{Supplementary information}
\DIFadd{Online Resource~1 holds the further tables (S1--S20) and figures referred to in the text: all rules with their intervals, robustness and sensitivity variants, coverage per characteristic and stratum, the step-6 decision rule, guard bands per family, out of sample and under other priors, costs, cross-fitting, few-shot transfer, design-level verdicts, the full validation of the simulation and the redundancy of the scans.
}

\DIFaddend \bmhead{Acknowledgements}
The author thanks the authors of the RDDAC and DDACS datasets at the University of Stuttgart for publishing the data on which this work rests.

\section*{Declarations}

\bmhead{Funding}
No specific funding was received for this work.

\bmhead{Competing interests}
The author declares no competing interests.

\bmhead{Ethics approval and consent}
Not applicable.

\bmhead{Data availability}
All data are public. RDDAC: doi:10.18419/DARUS-5589 (version 2.0); DDACS: doi:10.18419/DARUS-4801 (version 3.0); both licensed CC BY 4.0. The extracted per-part and per-simulation feature tables and all result files are archived with the code \citep{canseven2026archive}.

\bmhead{Code availability}
The complete analysis pipeline, which reproduces every number in this paper from the public data, is archived with the feature tables and result files \DIFdelbegin \DIFdel{of this version }\DIFdelend on Zenodo \citep{canseven2026archive}; Table~\ref{tab:mapping} maps the numbers to its scripts and files\DIFaddbegin \DIFadd{, and }\texttt{\DIFadd{scripts/opchar.py}} \DIFadd{is the scoring function for other datasets}\DIFaddend . A working copy is at \url{https://github.com/htcanseven/Claude} (directory \texttt{scripts/}).

\bmhead{Author contributions}
H.T.C. conceived the method, designed and carried out the analysis, and wrote the paper.

\bmhead{Use of generative AI}
A large language model (Claude, Anthropic) \DIFdelbegin \DIFdel{assisted }\DIFdelend \DIFaddbegin \DIFadd{was used }\DIFaddend under the author's direction \DIFaddbegin \DIFadd{to write and debug analysis code, retrieve and check literature, run analyses and draft text. It was also used to simulate reviews of earlier versions of the manuscript in five rounds, acting as referees and editor; most rules and analyses added after the first results were added in response to these simulated reviews, among them the nearest-setting baseline NN, M1s}\DIFaddend , \DIFdelbegin \DIFdel{as }\DIFdelend \DIFaddbegin \DIFadd{M1sn, M7n, M7, Mcd, the jackknife+ and relation-matched intervals, the guard bands with their priors and their validation, the decision-time strata, the costs, the between-run sensitivity and the cross-fitting (}\DIFaddend Section~\ref{sec:computation}\DIFdelbegin \DIFdel{describes}\DIFdelend \DIFaddbegin \DIFadd{). None of these reviews is a review by a journal}\DIFaddend . The author designed the study, verified \DIFdelbegin \DIFdel{all analyses}\DIFdelend \DIFaddbegin \DIFadd{every analysis against the archived results}\DIFaddend , checked every reference against its published record, and takes full responsibility for the content\DIFaddbegin \DIFadd{; the language model is not an author}\DIFaddend . All figures are produced by the released code; no AI-generated images are used.


%% ----------------------------------------------------------------------
\FloatBarrier
{\small
\bibliography{refs}}

\end{document}
